% 本文件由 LaTeX 排版 Agent 根据有效 translation.json 自动生成。
% author=Clement of Alexandria; work_id=0209; title=《导师》（The Instructor）

\chapter{《导师》（The Instructor）}

% structure: p0001 source=h1 target=omitted reason=page-title-already-rendered-by-parent

第一卷\newline{}第二卷\newline{}第三卷\par

\SourceRecord{Translated by William Wilson.；Ante-Nicene Fathers；Vol. 2.；Edited by Alexander Roberts, James Donaldson, and A. Cleveland Coxe.；Buffalo, NY: Christian Literature Publishing Co.,；1885.；<http://www.newadvent.org/fathers/0209.htm>.}

% structure: p0001 source=h1 target=section reason=page-title

\section{\WorkTitle{师者（第一卷）}}

\Emph{（师者）}\par

% structure: p0003 source=h2 target=subsection reason=relative-heading-rank; base=h2

\subsection{第一章 师者的职责}

在人身上，有这三件事：习惯、行为与情欲。习惯乃是\Emph{劝诫}性论述的领域，而劝诫论述是敬虔的向导，如同船的龙骨，安放在信德建造的根基之下；我们因这信德而欣喜异常，摒弃旧见，借着救恩焕发青春，同那预言般的诗歌一同歌唱：\Quote{天主对待心地正直的以色列人，何其美善！}行为则属于\Emph{诫命}性论述的领域；而\Emph{劝导}性论述则致力于医治情欲。然而，同一位圣言既将人从他自幼养成的世俗习俗中解救出来，又训练他走上唯一的天主救恩中的信德之路。\par

因此，当天上的向导——圣言——邀请人归向救恩时，祂被恰当地称为\Emph{劝诫者}；同一圣言被称为激励者（以部分代整体）。因为整个敬虔都是劝诫性的，在相似的理性官能中孕生对现世及来世真实生命的渴望。但现在，祂既是医治者，又是诫命者，循着自己的脚步，将所定规之事变为劝导的对象，应许医治我们内在的情欲。因此，我们适宜地用同一个名字称呼这圣言：导师（或\OriginalTerm{师者}{Pædagogue}，或师者）。\par

师者是实践的，而非理论的；祂的目标因此是改善灵魂，而非教导；是训练灵魂达到有德的生活，而非有知识的生活。尽管这同一圣言也是教导性的，但在此处并非如此。因为那在教义问题上解释并启示的圣言，其职责才是教导。但我们的教育者乃是实践的，祂首先劝勉人获得正当的心态与品格，然后说服我们积极地履行责任，赋予我们纯洁的诫命，并向后来者展示那些先前在错误中游荡之人的榜样。两者都具有最高的益处——那种以规劝形式引导人顺从的，以及那种以榜样形式呈现的；后者又分为两种，对应于前一种二元性：其一是让我们选择并模仿善者，其二是让我们拒绝并远离对立者。\par

由此，随着那些榜样的缓解，我们的情欲的医治随之而来；师者增强我们的灵魂，并以祂仁慈的命令，如同温和的药物，引导病者达至对真理的圆满认识。\par

健康与知识之间有很大的区别；因为知识由学习产生，健康由医治产生。一个患病的人，在他完全康复之前，不会学习任何学问。因为对于学习者与病人，每条命令的表达方式并非总是一样；对前者是引导至知识，对后者是引导至健康。因此，正如我们中那些身体患病的人需要医生，同样，灵魂患病的人需要一个师者来医治我们的疾病；然后，当灵魂能够接纳圣言的启示时，便需要一个教师来训练并引导它获得一切必要的知识。因此，那全善的圣言渴望通过有益于救恩的渐进方式，按照适于有效训练的顺序，完善我们；祂首先劝勉，然后训练，最后教导。\par

% structure: p0009 source=h2 target=subsection reason=relative-heading-rank; base=h2

\subsection{第二章 师者对待我们的罪}

如今，我的孩子们，我们的师者如同祂的天主父一样，是祂的儿子，无罪、无可指摘，且灵魂毫无情欲；天主成为人的形象，无玷污，是祂父旨意的仆役，是圣言——祂是天主——在父内，在父的右边，以天主的形象而为天主。祂是我们无瑕的肖像；我们当竭尽全力使我们的灵魂与祂相似。祂完全超脱人的情欲；因此唯有祂是审判者，因为唯有祂无罪。然而，就我们所能的，让我们尽量少犯罪。因为首要之事，莫过于摆脱情欲和失调，其次是遏制我们陷入已成习惯的罪。因此，最佳的是根本不在任何方面犯罪，我们断言这唯独是天主的特权；其次是避免故意的过犯，这是智者的特征；第三是不犯许多非自愿的过失，这是那些受过卓越训练之人的特点。不要长期停留在罪中，让这一点排在最后。但这对那些被召回补赎、重新开始争战的人也是有益的。\par

而师者，如我所想，非常优美地通过梅瑟说：\Quote{若有人突然死在他身旁，他祝圣的头就立刻被玷污，应剃去头发，}\BibleRef{户 6:9}这里将非自愿之罪称为突然的死亡。祂说这罪污秽灵魂，使之不洁；因此祂规定要迅速医治，劝告立即剃去头发；也就是说，劝告将遮蔽理性的无知之发剃干净，使理性（其处于大脑）脱离罪恶的厚重物质而畅通无阻，迅疾走向补赎。然后，稍加几句，祂补充说：\Quote{先前的日子不算为无效，}\BibleRef{户 6:12}这显然是指那与理性相悖的罪。非自愿的行为祂称为\Quote{\Emph{突然}}，罪祂称为\Quote{非理性的}。因此，圣言——师者——接管了我们，目的是预防那与理性相悖的罪。\par

因此，请思量圣经的表述：\Quote{因此主说这些事。}先前所犯之罪被随后的\Quote{因此}一词所谴责，公义的审判也随之而来。先知们清楚地表明了这一点，他们说：\Quote{你们若没有犯罪，祂就不会发出这些警告。}\Quote{因此，主这样说：}\Quote{因为你们没有听从这些话，因此主这样行；}以及\Quote{因此，看哪，主说。}因为预言是因着顺服和悖逆而赐下的：为顺服，使我们得救；为悖逆，使我们受纠正。\par

我们的导师——圣言，因此藉着劝诫治愈灵魂不自然的激情。因为身体疾病的帮助被最恰当地称为医术——一种由人类技艺获得的技艺。但父的圣言是人类软弱的唯一医治者，是患病灵魂的圣洁慰藉者。经上说：\Quote{上主，我的天主，求祢拯救祢那信赖祢的仆人。上主，求祢怜悯我，因为我终日向祢呼号。}德谟克利特说：\Quote{医术治疗身体的疾病；智慧使灵魂脱离激情。}但善的导师、智慧、圣父的圣言，祂创造了人，并关怀祂所造之物的全部本性；全能的医师、人类的救主，治愈身体和灵魂。祂对瘫子说：\Quote{起来，拿起你所躺卧的床，回家去吧。}\BibleRef{谷 2:11} 瘫子立刻恢复了力量。祂对死者说：\Quote{拉匝禄，出来！}\BibleRef{若 11:43} 死者就从棺材里出来，如同死前一样，经历了复活。此外，祂藉着诫命和恩赐来治愈灵魂本身——诫命是在时间过程中，但祂慷慨地赐予恩赐，对我们罪人说：\Quote{你的罪赦了。}\BibleRef{玛 9:2}\par

然而，我们，一旦祂产生了这个想法，就成了祂的儿女，由祂有序的安排获得了最好、最稳固的地位。这安排首先环绕世界、诸天和太阳的运行，并为了人的益处而关注其余星辰的运动，然后忙于人本身——祂的全部关怀都集中在他身上；视他为最伟大的作品，用智慧和节制管理他的灵魂，用美丽和匀称调和身体。而人类行为中一切正直和有序的，都是祂正直和秩序之启发的成果。\par

% structure: p0015 source=h2 target=subsection reason=relative-heading-rank; base=h2

\subsection{第三章 导师的仁爱}

主施予一切善和一切帮助，既作为人又作为天主：作为天主，祂赦免我们的罪；作为人，祂训练我们不犯罪。因此，人理当为天主所爱，因为祂是祂的工程。其他受造之物，祂仅以命令之言创造，但人却是祂亲手亲自塑造，并向人吹入祂独有的气息。那么，由祂所塑造、且肖似祂的受造物，要么因其本身为可欲而被天主亲自创造，要么为了其他缘故而形成。如果人本身是可欲的，那么好者爱善，爱之魅力就在人内，且正是那被称为天主之神（或气息）的东西；但如果人是因其他缘故而成为可欲的，那么天主创造人的理由无非是：除非人存在，否则天主就无法成为善的创造者，人也无法达到认识天主。因为天主不会通过创造人以外的方式完成祂创造人的目的；天主在愿意时拥有的隐秘能力，祂在创造行为中通过向外发出力量而完全实现，从人那里得到祂所造的人；祂看见了祂所造的人，祂所愿的都成就了；没有什么天主不能做的。因此，天主所造的人，其本身是可欲的；因他而成为可欲的，就与他所关涉的对象结合；而这后者也是可接受和可爱的。\par

但什么是可爱的，却不为祂所爱呢？人已被证明是可爱的；因此人蒙天主所爱。因为人若不被爱，缘何独生子从父怀中被派遣而来，那信德的圣言，那超乎丰盛的信德呢？主亲自明确承认并说：\Quote{因为父自己爱你们，因为你们爱了我。}\BibleRef{若 16:27} 又说：\Quote{祢爱了他们，如同爱了我一样。}\BibleRef{若 17:23} 那么，导师所愿、所宣告的，祂在言语和行为上是如何的，祂如何命令当作的事，禁止相反的事，已经显明。\par

显然，那另一种教导性的论说强而有力，且属神，注重精确，默观奥迹。但如今暂且搁下。现在我们理当回报那慈爱地引导我们走向最佳生命者的爱，并按照祂旨意的训令生活，不仅履行所命者，或防范所禁者，还要远离某些榜样，尽己所能效法另一些榜样，如此按祂的肖像完成主的工程，并实现圣经关于我们按祂的形像和肖像受造的教导。因为我们在生命中徘徊，如在深暗中，需要一位不会失足或迷失的向导；我们的向导是至善者，并非盲目的，如圣经所说：\BibleQuote{盲人领盲人，一同掉入坑里。} \BibleRef{玛 15:14} 但圣言目光敏锐，洞察内心的隐秘。因此，正如那不能照亮的不算光，不能移动的不算动，不能爱的不算爱，同样，那无益于救恩、不引导得救的也不算善。那么，让我们以主的工程为目标，完成诫命；因为圣言亲自公开地成了血肉，\BibleRef{若 1:14} 展示了实践与默观的同一德性。为此，让我们视圣言为法律，视祂的命令与劝诫为通往永生的短而直的道路；因为祂的训诫充满说服，而非恐惧。\par

% structure: p0019 source=h2 target=subsection reason=relative-heading-rank; base=h2

\subsection{第四章  男女同受导师的管教}

那么，让我们日益拥抱这美好的服从，将自己交付于主；紧握最稳妥之物，即对祂信德的缆绳，并理解男人与女人的美德是相同的。因为若两人的天主是同一的，两人的主人也是同一的；同一教会，同一节制，同一谦逊；他们的食物相同，婚姻是平等的轭；呼吸、视觉、听觉、知识、望德、服从、爱德，全都相同。那些生活相同的人，共享相同的恩宠和相同的救恩；爱德与训练对他们都是共同的。\Quote{因为在这个世界，} 祂说，\BibleQuote{人们嫁娶，} \BibleRef{路 20:34} 在此唯一分男女；\Quote{但在那个世界，却不再如此。} 在那里，这基于婚姻结合的社会性、圣洁的生活的赏报，并非为男或女，而是为人——那分割人类的性欲被除去。因此，人的名字对男女也是共通的。为此，我认为阿提卡人不仅称男孩，也称女孩为 \OriginalTerm{小孩}{\GreekText{παιδάριον}}，用作通性词；如果喜剧诗人米南德在 \WorkTitle{被打的女子} 中的话对任何人而言是足够的权威，他这样说——\par

我的小女儿；因为本性上\newline{}\OriginalTerm{小孩}{\GreekText{παιδάριον}}最是可爱。\par

\OriginalTerm{羔羊}{\GreekText{Ἄρνες}} 这个词，指羔羊，也是雄性和雌性动物单纯的通称。\par

主自己将永远牧养我们，如同祂的羊群。阿们。但若没有牧人，无论是羊还是任何别的动物都无法生存，孩子没有导师也不能活，家仆没有主人也不能活。\par

% structure: p0024 source=h2 target=subsection reason=relative-heading-rank; base=h2

\subsection{第五章  凡按真理而行者都是天主的子女}

那么，教子学（\OriginalTerm{儿童教导}{\GreekText{παίδων} \GreekText{ἀγωγή}}）这个词本身已清楚说明。剩下的，我们要思考圣经所指的儿童；然后让导师管理他们。我们就是儿童。圣经以多种方式赞扬我们，并以多种比喻描述我们，用不同的名称赋予单纯信德以多样性。因此，在福音中，主站在岸上，对门徒说——他们恰好在打渔——且大声喊道：\BibleQuote{孩子们，你们有吃的吗？} \BibleRef{若 21:4} ——祂称呼那些已在门徒位置的人为孩子们。\Quote{有人带儿童到祂跟前，要祂按手祝福他们；门徒却阻止他们，耶稣说：\InnerQuote{让儿童来，不要阻止他们到我跟前，因为天国正是属于这样的人。}} \BibleRef{玛 19:14} 这说法是什么意思，主自己会阐明，祂说：\BibleQuote{你们若不悔改，变得像儿童一样，绝不能进入天国；} \BibleRef{玛 18:3} 在那里并非比喻地谈论重生，而是以儿童身上的纯朴作为我们效法的榜样。\par

先知之神也分别我们为儿女。经上说：\Quote{摘取}橄榄枝或棕榈枝，孩子们出去迎接主，喊着说：\BibleQuote{贺撒纳于达味之子！奉主名而来的，当受赞美！}\BibleRef{玛 21:9}光明、荣耀和赞美，伴随着向主的恳求；因为\Quote{贺撒纳}这个词译成希腊文就是这个意思。在我看来，圣经为了责备那些无心的人，暗指刚才提到的预言，说：\Quote{你们没有念过：\InnerQuote{从婴孩和吃奶的口中，你完善了赞美}吗？}主在福音中就这样激励祂的门徒，敦促他们留心听祂，因祂正赶着往父那里去；祂暗示自己不久将要离去，使听者更加热切，并表明他们应当比以往更不吝惜地把握真理，因为圣言即将升天。因此，祂又称他们为儿女；因为祂说：\BibleQuote{孩子们，我同你们在一起的时候不多了。}\BibleRef{若 13:33}祂又将天国比作坐在街市上的儿童，说：\Quote{我们向你们吹笛，你们却不跳舞；我们哀哭，你们却不捶胸。}以及其他相应的话。不仅福音持这些看法，预言也与之相符。达味因此说：\Quote{ \Emph{孩子们，赞美}主；\Emph{赞美主的}名。}依撒意亚也说：\BibleQuote{ \Emph{看，我与天主所赐给我的孩子们。} }\BibleRef{依 8:18}那么，你听到属于万邦的人在主眼中是儿子，就惊讶吗？那么，你似乎没有留意阿提卡方言，从中你可以学到，美丽、匀称、生来自由的年轻女子仍被称为\OriginalTerm{少女}{\GreekText{παιδίσκαι}}，而婢女被称为\OriginalTerm{婢女}{\GreekText{παιδισκάρια}}；后者也因青春焕发，被谄媚地称为少女。\par

当祂说：\BibleQuote{让我的羊羔站在我的右边}\BibleRef{玛 25:33}，祂暗示单纯的儿女，如同绵羊和羔羊的本性，而非成人；祂更偏爱羔羊，是由于祂特别看重人性中构成纯洁的那种温柔和单纯的气质。再者，当祂说：\Quote{像吃奶的牛犊}时，祂再次以比喻指我们；而\BibleQuote{像纯洁温和的鸽子}\BibleRef{玛 10:16}，所指的也是我们。此外，通过\Person{梅瑟}，祂命令\Quote{献两只雏鸽或两只斑鸠为赎罪祭}；从而说明这些幼鸟的无害、纯洁和温顺的本性取悦于天主，并解释同类为同类赎罪。再者，斑鸠的胆怯象征对罪的畏惧。\par

圣经证明祂称我们为小鸡：\BibleQuote{如同母鸡把小鸡聚集在翅膀底下}\BibleRef{玛 23:37}。因此，我们是主的小鸡；圣言如此奇妙而奥秘地描述了童年的单纯。因为有时祂称我们为儿女，有时为小鸡，有时为婴孩，有时为儿子，以及\Quote{新的人民}和\Quote{近来的人民}。\BibleQuote{我的仆人将被称为一个新名字}\BibleRef{依 65:15}（一个新名字，祂说，新鲜而永恒，纯洁而单纯，童稚而真实），这名字在地上要蒙福。再者，祂比喻性地称我们为未被罪恶驾驭、未被邪恶驯服的小马驹；而是单纯、欢跃地只奔向父；不是那样的马，\BibleQuote{如那些向邻舍的妻子嘶叫，套着轭、疯狂求雌的马}\BibleRef{耶 5:8}；而是自由而新生，藉信德欢跃，预备奔向真理，迅速奔赴救恩，践踏世物。\par

\BibleQuote{熙雍女子，喜乐吧！耶路撒冷女子，欢呼吧！看，你的君王来到你这里，正义、温和，带着救恩；祂的确温和，骑着一匹牲畜，一匹年幼的驴驹。}\BibleRef{匝 9:9} 只说驴驹还不够，祂还加上\Emph{年幼}，以显示基督内人性的青春，以及单纯的永恒，它不知衰老。我们这些小孩子既是这样的驴驹，就由我们神圣的驯驴师养育。但如果圣经中的新人由驴来代表，那么这驴也是驴驹。经上说：\Quote{祂把驴驹拴在葡萄树上}，将这单纯而童稚的子民拴在圣言上，祂比喻性地以葡萄树代表他们。因为葡萄树产酒，正如圣言产血，两者都为人的健康而饮——酒为身体，血为精神。\par

祂也称我们为羔羊，圣神曾藉依撒意亚的口作无可指摘的见证说：\BibleQuote{祂必像牧人牧养自己的羊群，用臂膀聚集羔羊。}\BibleRef{依 40:11}——用比绵羊更为柔弱的羔羊的比喻说法，来表达纯朴。我们也确实用从\Quote{孩童}一词得来的称呼来尊崇生活中最美好、最完美的事物，而将训练称为\OriginalTerm{培训}{\GreekText{παιδεία}}，将管教称为\OriginalTerm{管教}{\GreekText{παιδαγωγία}}。我们宣称\OriginalTerm{管教}{\GreekText{παιδαγωγία}}是从童年引导人走向德行的正确指引。因此，我们的主更明确地向我们启示了\Quote{孩童}这个称呼所象征的含义。当宗徒们中间起了争论，\Quote{他们中谁该是最大的}时，耶稣让一个小孩站在他们中间，说：\BibleQuote{凡谦卑自己像这个小孩的，他在天国里就是最大的。}\BibleRef{玛 18:4}祂并不是因为孩童因年龄而理解力极其有限才使用这一称呼，像有些人所想的那样。祂说：\BibleQuote{你们若不变成像这些小孩，绝不能进入天国。}这话也不应理解为\Quote{不学习}。我们这些婴孩，不再在地上打滚，也不像从前那样如蛇在地上爬行，用整个身体贴附着无意义的欲望；而是灵魂向上伸展，摆脱了世界和我们的罪，用脚尖触地，仿佛仍在这世界，却追求圣洁的智慧——尽管这在那些心思锐利只为作恶的人看来是愚拙。那么，那些认识独一天主为父的人，那些纯朴、婴孩般、无伪的人，那些喜爱独角兽之角的人，就正当地被称为小孩了。\par

因此，对于那些在圣言上已有长进的人，祂宣布了这句话，命令他们抛开对这世界之事的焦虑，劝勉他们效法孩童，唯独依附圣父。因此，在接下来的经文中祂说：\BibleQuote{不要为明天忧虑，因为明天自有明天的忧虑；一天的难处一天当就够了。} \BibleRef{玛 6:34} 这样，祂命令他们放下今生的一切挂虑，唯独依赖圣父。凡遵守这诫命的人，实际上对天主和世界而言是孩童和儿子——对前者而言是被欺骗的，对后者而言是被爱的。如果我们在天上只有一位师傅，正如圣经所说，那么地上的人理所当然应被称为门徒。因为真理就是这样：完美在于主，祂常教导；而幼稚和孩童气在于我们，我们常学习。因此，预言尊重\Emph{完美}，用\Emph{人}这个称谓来指称它。例如，通过达味，祂论及魔鬼说：\BibleQuote{上主憎恨流血的人；}祂称他为人，因为他在邪恶上是完美的。主被称为人，因为祂在义德上是完美的。宗徒的实例正切合此意，他在致格林多人书中写道：\BibleQuote{因为我曾把你们许配给一个丈夫，把你们如同贞洁的童女献给了基督。} \BibleRef{格后 11:2} 无论是作为孩童还是圣人，都唯独归于主。而在致厄弗所人书中，他极清楚地阐明了这一问题，说道：\BibleQuote{直到我们众人对天主的信仰和认识达成一致，达到成年人的身量，达到基督圆满年龄的程度；使我们不再做小孩子，被各种教义之风摇动，随从人的诡计和欺骗的法术；而要在爱德中持守真理，在各方面长进到祂里面，} \BibleRef{弗 4:13} ——他说这些话是为了建树基督的身体，基督是头，是唯一在义德上成全的人；而我们这些孩童，要防范异端之风，这些风吹动我们使我们自高自大；不信任那些教导我们不同道理的父辈，当我们成为教会，接受了基督为头时，我们才得以成全。然后，有必要留意关于\Quote{婴孩}（\OriginalTerm{婴孩}{\GreekText{νήπιος}}）这个称谓，\OriginalTerm{那婴孩}{\GreekText{τὸ} \GreekText{νήπιον}}并不是指愚蠢的人；因为愚蠢的人被称为\OriginalTerm{愚蠢的人}{\GreekText{νηπύτιος}}；而\OriginalTerm{婴孩}{\GreekText{νήπιος}}是\OriginalTerm{幼小者}{\GreekText{νεήπιος}}（因为温柔的人被称为\OriginalTerm{温柔}{\GreekText{ἤπιος}}），意指在行为上刚刚变得柔和温顺的人。蒙福的保禄最清楚地指出了这一点，他说：\BibleQuote{我们作为基督的宗徒，本可受人敬重，却在你们中间成了\OriginalTerm{温柔}{\GreekText{ἤπιοι}}的人，如同哺乳的母亲抚爱自己的孩子。} \BibleRef{得前 2:6} 因此，\OriginalTerm{孩童}{\GreekText{νήπιος}}就是\OriginalTerm{温柔}{\GreekText{ἤπιος}}，因而更为柔嫩、精致、纯朴、无邪、无伪，心地正直坦诚，这是纯朴与真理的基础。因为祂说：\BibleQuote{我必垂顾那温柔安静的人。} \BibleRef{依 66:2} 因为童贞的话语就是如此，柔嫩而无欺诈；因此，童女常被称为\Quote{温柔的新娘}，孩童常被称为\Quote{柔心的人}。而我们这些柔嫩的人，易受劝说的力量影响，容易被引向良善，性情温和，没有恶意和悖逆的玷污，因为古时的各族是悖逆和硬心的；但我们这婴孩的队伍，即我们这新民族，如同孩童一样细腻。为着纯朴之心，宗徒在致罗马人书中承认自己喜乐，并提供了某种孩童的定义，他说：\BibleQuote{我愿你们在善事上智慧，在恶事上天真。} \BibleRef{罗 16:19} 因为\Quote{孩童}这个名称，即\OriginalTerm{孩童}{\GreekText{νήπιος}}，我们并不理解为否定意义，尽管文法家们把\OriginalTerm{否定前缀}{\GreekText{νη}}当作否定小词。他们若称我们这些追随孩童状态的人为愚蠢，看哪，他们竟对主说出亵渎的话，将那些投靠天主的人视为愚笨。但如果他们自己理解为\Quote{纯朴者的孩童}这个名称——这倒是真正的含义，我们就以这名称为荣。因为那些新近变得智慧、照新约而生的人心，在旧日的愚妄中仍是幼稚的。于是，在末世，天主因基督的来临而被认识：\BibleQuote{除了子，没有人认识天主；除了子和子所愿意启示的人，也没有人认识天主。} \BibleRef{玛 11:27}\par

因此，与旧人相对，新人被称为年轻的，因为他们学会了新的恩赐；我们拥有生命清晨的丰盛，在这不知衰老的青春中，我们总是在理智上走向成熟，永远年轻，永远温和，永远新颖：因为那些分享新圣言的人必然是新生的。而那参与永恒的事物，惯于相似于不朽之物；因此，我们配得上童年的称呼，一个终生的春天，因为我们内在的真理，以及浸透了真理的习惯，不会被衰老触及；但智慧永远绽放，始终如一，永不改变。\Quote{他们的孩子，}经上说，\Quote{必被抱在怀中，放在膝上抚弄；就如母亲安慰人，我也要这样安慰你们。}\BibleRef{依 66:12}母亲将孩子拉近自己；我们寻求我们的母亲教会。凡是软弱柔嫩的，因软弱而需要帮助的，都被友善看待，甜美可亲，对于这样的，愤怒变为帮助：因为马驹、牛犊、狮崽、鹿羔以及人的孩童，都被他们的父母喜悦地看待。同样，宇宙之父对投奔祂的人也怀有慈爱；祂藉着祂的圣神重生他们，使他们获得子女的义子名分，认识他们是温良的，唯独爱他们，帮助他们，为他们争战；因此祂赐给他们孩子的名字。我也将\Person{依撒格}一词与孩子相联。依撒格意为\OriginalTerm{欢笑}{laughter}。他被窥探的国王看见与妻子及助手\Person{黎贝加}嬉戏\BibleRef{创 26:8}。那位名叫\Person{阿彼默肋客}的国王，在我看来，代表一种超越世俗的智慧，默观嬉戏的奥秘。他们将黎贝加解释为\OriginalTerm{坚忍}{endurance}。啊，智慧的嬉戏，欢笑也由坚忍辅助，国王是观众！那些在基督内的孩子们的心灵，他们的生活以坚忍为秩序，便欢欣鼓舞。这就是神圣的嬉戏。\Person{赫拉克利特}说：\Quote{这样的嬉戏，\Person{宙斯}独自嬉戏。}因为除了在坚忍中为善而嬉戏欢庆，并在善的管理中与天主一同庆节，还有什么是智慧而完美的人所宜做的呢？先知所预示的可以有不同的解释——即我们因救恩而喜乐，正如依撒格。他也从死亡中被救出，欢笑、嬉戏，与他的配偶一同喜乐，这配偶是我们救恩的帮助者——教会的预像，教会得到了坚忍这个稳固的名字；确实，因为她独自存留于万代，永远喜乐，因着我们信徒——基督的肢体——的坚忍而存在。那些坚忍到底者的见证，以及因他们而有的喜乐，就是神秘的嬉戏，以及伴随着体面慰藉的救恩，这救恩给我们带来帮助。\par

\Person{王}，就是\Person{基督}，从高处俯瞰我们的欢笑，透过窗户观看（如经上所说），注视着感恩、祝福、欢欣、喜乐，还有与之合作的坚忍以及他们的拥抱：注视祂的教会，只显示祂的面容，这面容原是教会所欠缺的，因她的元首君王而得以完美。那么，主显示自己的门在哪里呢？就是祂所显现的肉身。祂就是\Person{依撒格}（因为叙述可有不同解释），他是主的预像，是一个作为儿子的孩子；因为他是\Person{亚巴郎}的儿子，如同基督是天主之子，也是如同主一样的祭品，但他并未像主那样被祭杀。依撒格只背负了祭献的木柴，如同主背负了十字架的木柴。他神秘地欢笑，预言主要用喜乐充满我们这些因主宝血而从败坏中被救赎的人。依撒格做了除受苦之外的一切事，这是恰当的，将受苦的优先地位让给了\OriginalTerm{圣言}{Word}。此外，主未被杀害也暗示了祂的天主性。因为\Person{耶稣}在埋葬后复活了，没有受到任何伤害，如同依撒格从祭献中被释放。为了证明这一点，我要提出另一个非常重要的考虑。\Person{圣神}亲自称主为一个孩子，通过\Person{依撒意亚}预言说：\Quote{看，一个孩子为我们诞生了，一个儿子赐给了我们，权柄必担在他的肩上；他的名字被称为伟大的谋士的使者。}同一位先知宣告了祂的伟大：\Quote{奇妙的、谋士、全能的天主、永远的大父、和平之王；为使祂的训诲得以成就：祂的和平没有终结。}\BibleRef{依 9:6}啊，伟大的天主！啊，完美的孩子！子在父内，父在子内。这个孩子的训诲怎能不完美呢？它普及万有，如同导师带领我们这些小孩子——祂的小儿女。祂向我们伸出了那双特别值得信赖的手。为了这个孩子，还有\Person{若翰}作证，他说：\Quote{在妇女所生者中最大的先知，}\BibleRef{路 7:28}请看天主的羔羊！因为圣经称婴孩为羔羊，所以也称那为了我们而成为人、并愿意在各方面相似我们的天主圣言为\Quote{天主的羔羊}——祂就是天主子，圣父的孩子。\par

% structure: p0034 source=h2 target=subsection reason=relative-heading-rank; base=h2

\subsection{第六章 孩子之名并非指初等原理的教导}

我们有充足的方法应对那些好挑剔的人。因为我们被称为孩童和婴儿，并非如那些因知识而自高自大者诽谤的那样，指我们的教育幼稚可鄙。恰恰相反，从我们重生那一刻起，我们就获得了我们所渴望的完美。因为我们受了光照，即认识天主。认识完美者的人，他本身并不残缺。当我承认认识天主时，请不要责备我；因为这样对圣言说话是合宜的，而祂是自由的。因为在主受洗的那一刻，从天上有声音为那蒙爱者作证说：\BibleQuote{祢是我的爱子，我今天生了祢。}那么，让我们问那些智者：今天所生的基督，已是完美的呢，还是——这最为荒谬——不完美的？若是后者，祂还需要有所增添。但认为祂在知识上有所增添是荒谬的，因为祂是天主。因为没有人能胜过圣言，也没有人能教导那唯一的导师。难道他们不会勉强承认，完美的圣言出自完美的圣父，乃是按照经济的预旨完美地受生吗？如果祂是完美的，那完美的祂为何要受洗呢？他们说，这是为了完成属于人性的职分。极好。那么我断言，在祂受若翰洗的同时，祂就变得完美吗？显然如此。那么祂没有从他学到更多东西？当然没有。但祂单凭洗\OriginalTerm{礼}{baptism}之水得以完美，并因圣神的降临而圣化吗？正是这样。同样的事也发生在我们身上，基督成了我们的模范。我们受洗，就受了光照；受了光照，就成为儿子；成为儿子，就成为完美；成为完美，就成为不朽的。\BibleQuote{我}祂说\BibleQuote{说过：你们是神，全都是至高者的儿子。}这工程有多种名称：恩宠、光照、完美和洗涤：洗涤，借此我们洗净罪过；恩宠，借此犯罪应受的惩罚得以赦免；光照，借此看见那神圣的救恩之光，亦即借此我们清楚地看见天主。我们称完美为毫无欠缺。因为认识天主的人还缺什么呢？因为把不完美的东西称为天主的恩宠（或恩赐）实在是荒谬的。祂既是完美的，自然施予完美的恩赐。正如万物因祂的命令而造成，同样，祂只要愿意施恩，祂的恩宠之完美随之成就。因为将来的时间被祂意志的力量所预取。\par

脱离邪恶是救恩的开始。那么只有我们，首先触及生命的边界，已经是完美的了；已经活着，因为我们脱离了死亡。因此，救恩就是跟随基督：因为生命就在祂内。\BibleRef{若 1:4}\BibleQuote{我实实在在告诉你们：那听我的话，并相信派遣我来者的，便有永生，不受审判，而已出死入生。}\BibleRef{若 5:24}因此，单单相信和重生就是生命中的完美；因为天主从不软弱。因为正如祂的意愿就是工作，而工作被称为世界；同样，祂的谋划就是人的救恩，这被称为教会。所以祂认识祂所召叫的，也认识祂所拯救的；而且祂在同一时刻召叫并拯救了他们。\BibleQuote{因为你们}宗徒说\BibleQuote{是受天主教导的。}\BibleRef{得前 4:9}所以不允许认为由祂所教导的是不完美的；从祂所学到的，就是永恒救主的永恒救恩，愿感谢归于祂，直到永远。阿们。而那仅仅重生——正如名称所必然表明的——并受了光照的人，立刻脱离黑暗，并在瞬间接受光明。\par

就如那些摆脱了睡眠的人立刻在内里完全清醒；或者更好说，如同那些试图除去眼睛上薄膜的人，并不从外面供给它们所没有的光，而是除去眼睛的障碍，让瞳孔自由；同样，我们这些领受圣洗的人，擦去了遮蔽天主圣神之光的罪，拥有了圣神之眼，自由、无碍且充满光明，惟独藉此我们默观天主，圣神从上倾注于我们。这是视力的永恒调整，能够看见永恒的光，因为相似者爱相似者；那圣洁的，爱慕那从它而出之圣洁，其适切地被称为光。\BibleQuote{从前你们是黑暗，如今在主内你们是光。} \BibleRef{弗 5:8} 因此我认为古人称人为\OriginalTerm{光}{\GreekText{φώς}}。但他们说，人尚未领受圆满的恩赐。我也同意这一点；但他在光中，黑暗不曾胜过他。光与暗之间没有中间物。但结局保留给那些相信者的复活；它并非领受别的什么东西，而是获得先前所许的恩许。因为我们不说两者同时发生——既到达终点，又预期那到达。因为永恒与时间不同，努力与最终结果也不同；但两者都指向同一件事，并且同一个人涉及两者。信德，可以说，是在时间中产生的努力；最终结果是获得恩许，永恒确保。主自己已最清楚地揭示了救恩的平等，当祂说：\BibleQuote{这是我父的旨意，凡看见子并信祂的人，得享永生；在末日我要叫他复活。} \BibleRef{若 6:40} 在此世（祂意指末日，且保存至该终止的时刻），我们相信得以成全。因此祂说：\BibleQuote{信子的人有永生。} \BibleRef{若 3:36} 如果因此那些相信的人已拥有生命，那么拥有永生之外还有什么呢？信德一无所缺，因它本身完美圆满。若它缺少什么，就非完全完美。但信德没有任何方面是瘸腿的；在我们离世之后，它也不使我们这些已相信并毫无区别地领受了未来美物凭据的人等待；而是以信德预先抓住那未来的，在复活后我们便视之为现在，为使那所说的得以应验：\BibleQuote{照你的信德成就吧。} \BibleRef{玛 9:29} 那里有信德，那里就有恩许；恩许的完成就是安息。因此，在光照中我们领受的是知识，知识的终点是安息——那被构思为渴望对象的最终之物。因此，正如无经验因经验而终止，困惑因找到清晰出路而终止，同样，黑暗必因光照而消散。黑暗是无知，我们藉此陷入罪中，对真理视而不见。因此，知识是我们领受的光照，它使无知消散，并赐予我们清晰的视力。此外，舍弃恶就是采纳更善。因为无知所捆绑的恶，被知识妥善释放；那些束缚藉人的信德和天主的恩宠迅速松开，我们的过犯被一种医神之药（派翁的）除去，就是圣言的圣洗。我们从一切罪中洗净，不再被恶缠累。这是光照的唯一恩宠：我们的品性不同于洗前。既然知识随光照而生，将其光芒照入心智，我们一听闻，便从无知者成为门徒。请问，这是否在指导来临时发生？你不能说出时间。因为指导引向信德，信德与圣洗一起由圣神训练。因为那信德是人类唯一的普遍救恩，在公义而慈爱的天主面前有同一平等，祂与众人有同一共融，宗徒最清楚地表明了，说：\Quote{在信德来临之前，我们被看守在法律之下，被禁锢于那将要启示的信德，因此法律成了我们的启蒙师，领我们归于基督，好使我们因信德成义；但信德来临后，我们便不再在启蒙师之下了。} 你没听见我们不再在那伴有恐惧的法律之下，而是在圣言——自由选择的导师——之下吗？然后他补充了毫无偏私的话：\BibleQuote{因为你们众人藉着在基督耶稣内的信德，都是天主的子女。凡领洗归于基督的，就是穿上了基督。不再有犹太人或希腊人，不再有奴隶或自由人，不再有男人或女人，因为你们众人在基督耶稣内合而为一。} \BibleRef{迦 3:26-28} 因此，在同一个圣言中，并没有一些是\Quote{光照者（诺斯替派）}，另一些是\Quote{属血气（或属自然）的人}；但所有放弃肉身欲望的人，在主面前都是平等和属神的。他又在另一处写道：\BibleQuote{因为我们都因同一圣神受洗，成为一个身体，无论是犹太人或希腊人，奴隶或自由人，我们都饮于同一杯。} \BibleRef{格前 12:13} 若采用那些称更美之事的回忆为圣神过滤的人的表达，也并非荒谬，藉过滤理解那更美之事的回忆所带来的卑微之物的分离。在回忆更美之事的人身上，必然随之而来对更坏之事的悔改。因此，他们承认在悔改中，圣神折返其步。同样地，我们也悔改我们的罪，弃绝我们的过犯，藉圣洗洁净，飞快回到永恒的光中，作为子女归于父。因此，耶稣在圣神内欢欣说：\BibleQuote{父啊，天地的主宰，我感谢祢，因为祢将这些事瞒住了智慧及明达的人，而启示给了小孩子；} \BibleRef{路 10:21} 师傅与导师称我们为小孩子，我们比世上的智者更乐于接纳救恩，他们自以为智慧，却充满骄傲。祂在欢喜和极大的喜乐中呼喊，似乎与孩子们咿呀学语：\BibleQuote{是的，父啊！因为祢的美意本是如此。} \BibleRef{路 10:21} 因此，那些向这现世的智慧及明达的人所隐藏的事，却启示给了小孩子。我们真是天主的子女，脱去了旧人，剥去了邪恶的衣服，穿上了基督的不朽；使我们成为新圣洁的子民，藉重生保持人无玷。一个小孩子，作为天主的小家伙，从奸淫和邪恶中被洁净。最有福的保禄在他写给格林多人的第一封书信中极清楚地为我们解答了这个问题，这样写道：\BibleQuote{弟兄们！在见识上不要做小孩子，但要在恶事上做小孩子，在见识上要做成年人。} \BibleRef{格前 14:20} 而\Quote{我做孩子的时候，说话像孩子，思想像孩子}这句话，指出他按法律的生活，在其中他想孩子的事，说话像孩子，亵渎了圣言，因他尚未达到孩子的纯朴，而是处在愚昧中；因为\OriginalTerm{婴孩}{\GreekText{νήπιον}}一词有两个意思。\Quote{当我成了人}保禄又说，\Quote{便丢弃了孩子的事。} \BibleRef{格前 13:11} 宗徒本人自称是幼稚事的传道者，当他仿佛将其放逐时，所暗示的并非身量的不完满，也非固定的时间尺度，亦非在更成人、更完善之事上的附加秘密教导；而是称那些在法律之下的人为\Quote{孩子}，他们因恐惧而害怕，如同孩子被怪物惊吓；称我们为\Quote{成人}，我们服从圣言，作自己的主人，相信并因自由选择得救，理性而非非理性地被恐惧惊吓。宗徒本人将为此作证，称犹太人为按首份盟约的继承人，我们为按恩许的继承人：\BibleQuote{\Emph{我说：继承人几时还是孩子，虽是全业的主人，却与奴隶没有分别，仍属于监护人和管家，直到父亲预定的时期。同样，我们以前做孩子的时候，也是被世界的要素所奴役；但时期一满，天主就派遣了自己的儿子，生于女人，生于法律之下，为要把在法律之下的人赎出来，使我们获得义子的地位。}} \BibleRef{迦 4:1-5} 藉着祂。看祂如何接纳那些在恐惧和罪恶之下的人为孩子；但将成人赐给那些在信德之下的人，称他们为儿子，与在法律之下的孩子相对：\Quote{所以你不再是奴隶，而是儿子；如果是儿子，也藉着天主成为继承人。} \BibleRef{迦 4:7} 那么，儿子在继承之后还缺什么呢？因此，\Quote{我做孩子的时候}这句话可优雅地解释为：即当我还是犹太人时（他生来是希伯来人），我思想像孩子，当我遵循法律时；但成了人以后，我不再怀有孩子的情感，即法律的情感，而是成人的，即基督的情感，圣经单独称祂为人，如我们先前所说。\Quote{我丢弃了孩子的事。} 但在基督里的孩子是成熟，相对于法律而言。达到这一点，我们必须捍卫我们的孩子身份。我们还需解释宗徒的话：\BibleQuote{\Emph{我给你们喝的是奶，是在基督里；}} \BibleRef{格前 3:2} 因为在我看来，这话不应按犹太意义理解；因为我要用另一圣经反驳它：\BibleQuote{我要领你们进入那流奶与蜜的美地。} \BibleRef{出 3:8} 关于这些圣经的比较，当我们思考时，产生了一个极大的难题。因为如果那以奶为特征的孩子身份是在基督内信德的开始，那么它被贬低为幼稚和不完善。那在饭食之后的安息，即完美而赋有知识之人的安息，又如何再次以婴儿之奶来区分呢？难道这不正如同解释比喻，意味着类似的事吗？这话是否应大致如下阅读：\BibleQuote{\Emph{我给你们喝的是奶，是在基督里；}}然后略作停顿，我们加上\Quote{如同孩子，}这样在读时通过分开词语，我们可得出这样的意思：我用简单、真实而自然的滋养——即属神的滋养——在基督内教导你们；因为奶正是从爱之乳房涌出的滋养物质。因此，整体可如此理解：如同乳母用奶喂养新生婴儿，同样我也用圣言，即基督的奶，将属神的养分注入你们心中。\par

因此，完美的奶便是完美的营养，并引向那不能止息的圆满。为此，在安息中也应许了同样的奶和蜜。所以，主又向义人应许奶，为清楚显示圣言既是\BibleQuote{阿耳法和敖默加，起初和终结}\BibleRef{默 1:8}，圣言被象征性地描绘为奶。荷马也违其本意，似神谕般宣称，称义人为\OriginalTerm{喝奶的}{\GreekText{γαλακτοφάγοι}}。同样，我们也可理解圣经：\BibleQuote{弟兄们，我从前不能向你们说话，如同对属神的人，只能如同对属血肉的人，如同对在基督内的婴孩}\BibleRef{格前 3:1}，因此，属血肉的人可理解为刚受教导、仍是在基督内的婴孩。因为他称那些已信了圣神的人为属神的，而称那些新受教导、尚未洁净的人为属血肉的；他公正地称他们仍是属血肉的，因为他们与外邦人一样，思念肉身的事：\BibleQuote{因为在你们中间有嫉妒和纷争，你们岂不是属血肉的，并且按照人的样子行事吗？}\BibleRef{格前 3:3} 他又说：\BibleQuote{因此，我给你们奶喝}，意思是：我已将那种从教导中滋养到永生的知识灌输给你们。但\Quote{我给你们喝}这一表达（\OriginalTerm{我给你们喝}{\GreekText{ἐπότισα}}）是完全领受的象征。因为长成的人被称为喝，婴孩则是吮吸。主说：\BibleQuote{因为我的血是真正的饮料}\BibleRef{若 6:55}。因此，当他说\BibleQuote{我给你们奶喝}时，岂不是指明了真理的知识，即在圣言——那奶——中的完美喜乐吗？紧接着的话：\Quote{不是固体食物，因为你们还不能吃}，可能指示将来世界里清晰的启示，如同食物，面对面。\BibleQuote{因为现在我们对着镜子观看}，同一宗徒说，\BibleQuote{但那时就要面对面了}\BibleRef{格前 13:12}。因此他又加上：\Quote{就是现在你们还不能，因为你们仍是属血肉的}，思念肉身的事——贪爱、嫉妒、忿怒、怨恨。\BibleQuote{因为我们已不在肉身内}\BibleRef{罗 8:9}，如有些人所认为。因为他们说，同这肉身一起，具有如同天使的面容，我们将面对面地看见应许。然则，如果那真是我们离世后的应许，他们怎么说他们知道\Quote{眼所未见，人心也未曾想到的}呢？他们并非借着圣神领悟，而是从教导领受了\Quote{耳所未闻的}，或那只\Quote{被提到第三层天}的耳朵？但即使那时，他也奉命保持沉默。\par

如果人的智慧——我们仍需明白——是以知识为荣，那么请听圣经的法律：\BibleQuote{智者不要因他的智慧而夸耀，勇士不要因他的勇力而夸耀；夸耀的，应因主而夸耀。}然而我们是受天主教诲的，因基督的名而夸耀。那么，我们怎能不认为宗徒将这一意义赋予了婴儿的奶呢？如果我们这些治理教会的人，按照善牧的形象是牧人，而你们是羊群，那么当主也谈到羊群的奶时，我们岂不认为祂在比喻用法上保持了一致吗？其次，我们可以将这一意义应用于这句话：\BibleQuote{我给你们喝的是奶，不是食物，因为你们还不能吃。}认为肉与奶并非不同，而是本质相同。因为同一个圣言，或像奶一样流动柔和，或像肉一样坚实紧密。持此看法，我们可以将普遍传播的福音宣讲视为奶；而将信德视为肉，它由教导凝聚成根基，比听闻更坚固，故被比作肉，并将这种滋养同化于灵魂本身。主在 \BibleBook{若望}{福音} 中曾以象征表达这一点，说：\BibleQuote{你们吃我的肉，喝我的血；}\BibleRef{若 6:34}清楚地以隐喻描述了信德与应许的可饮用的特性，教会借此，如同由许多肢体组成的人，获得滋养而成长，被焊接并凝聚于两者——作为身体的信德和作为灵魂的望德；正如主也是由肉和血构成。实际上，信德的血是望德，信德如同被生命原则束缚于其中。当望德消失时，如同血流出；信德的生命力也遭到破坏。如果有人反对，说奶是指初步教训——就像最初的食物——而肉是指那些通过提升知识而获得的属神认知，让他们明白：当他们说肉是固体食物、是耶稣的肉和血时，他们因自己的虚荣智慧而被引向真正的单纯。因为血在人体中被发现是原始产物，有人竟敢称其为灵魂的实质。这血在母亲怀孕时，通过自然的同化过程，因父母之爱的交感而绽放并老化，以免孩子受到恐惧。血也是肉中较湿润的部分，是一种液态的肉；而奶是血中较甜、较精细的部分。无论是供给胎儿、通过母亲肚脐送去的血，还是被阻于正常通道之外、经自然扩散、由滋养万物的创造天主命令、流向已经膨胀的乳房、并被气息的热力转化的经血，转变的都是血。因为所有肢体中，乳房与子宫最有交感。分娩时，输送血液给胎儿的血管被切断：血流受阻，血液涌向乳房；大量血液涌入后，乳房膨胀，将血转化为奶，类似于溃疡中血转化为脓。另一方面，如果怀孕期间乳房附近静脉中的血液被注入乳房的天然空隙，而相邻动脉排出的气息与之混合，血液本质仍然纯净，如同波浪般被搅动而变白；通过这种中断，它像发泡一样变化，类似大海，在风的冲击下，诗人们说：\Quote{吐出咸味的泡沫。}然而本质仍由血液提供。\par

同样，河流奔腾而下，与周围空气接触而激荡，喃喃地涌出泡沫。我们口中的水分也被气息变白。那么，不承认血被气息转化为那种非常洁白明亮的物质，是何等荒谬！它所遭受的变化是在性质上，而非本质上。你必定找不到任何比奶更有营养、更甜或更白的东西。因此，它在各方面都像属神的滋养，因恩宠而甘甜，如同生命般滋养，如同基督之日般明亮。\par

圣言的血也曾被显示为乳。在生育中，乳这样被预备好，供给婴孩；而原本直向丈夫的乳房，现在弯向孩子，被教导以易于接受的方式供应自然所制成的、有益于滋养的物质。因为乳房不像满是乳汁的泉源，预先准备好而流出；而是通过改变养料，在自身中形成乳并排出。对于新生和重生的婴孩，那适宜而健康的养料由天主——所有被生和重生者的养育者和圣父——所精心制成，如同玛纳，天使的天上食粮，从天降给古希伯来人。事实上，就连现今，乳母们也将初次挤出的乳汁称为与那食粮相同的名字——玛纳。此外，怀孕的妇女成为母亲时，排出乳汁。但主基督，童贞女的果实，并没有称女人的乳房为有福，也没有选择它们来给予滋养；但当仁慈而慈爱的圣父降下圣言时，祂自己便成为善人的神性滋养。神秘的奇事！万有的圣父是唯一的，万有的圣言也是唯一的；圣神处处是同一的，且唯一的是那童贞母亲。我爱称她为教会。这位母亲，当她独自时，没有乳汁，因为独自时她不是女人。但她同时是童贞和母亲——纯洁如童贞，慈爱如母亲。她将儿女唤到自己身边，用圣洁的乳汁——即为孩童的圣言——哺育他们。因此她没有乳汁；因为乳汁就是这美丽俊美的孩童，基督的身体，祂用圣言滋养那祂自己在肉身疼痛中产下的幼小群羊，那祂自己用祂的宝血包裹的群羊。啊，惊人的诞生！啊，圣洁的襁褓！圣言对孩童就是一切，既是父亲又是母亲，既是导师又是乳母。祂说：\Quote{你们吃我的肉，并喝我的血。}\BibleRef{若 6:53}这就是主所供应的适宜食物，祂献出自己的肉，倾流自己的血，为孩童的成长毫无欠缺。啊，惊人的奥迹！我们被命令要脱去旧有的、属血肉的败坏，以及旧有的养料，换取另一种新的生活方式，就是基督的生活方式，接受祂，如果我们能的话，将祂藏在内心；这样，我们将救主供奉在灵魂中，就能矫正我们肉体的情感。\par

但你们并不倾向于这样理解，或许更倾向于一般性的理解。也请以下面的方式来听。肉身比喻性地代表圣神；因为肉身是由祂所造。血指向我们指出圣言，因为如同丰富的血，圣言被注入生命；而二者的结合就是主，婴孩的食物——那既是圣神又是圣言的主。食物——即主耶稣——即天主的圣言，成为肉身的圣神，被圣化的天上肉身。营养是父的奶，只有靠它我们这些婴孩才得喂养。那么圣言本身，即那蒙爱者，我们的喂养者，为我们倾流了自己的血，以拯救人类；藉着祂，我们信从天主，逃向主的\Quote{抚慰忧患的怀抱}。唯有祂，恰如其分地，用爱的奶供给我们孩童，只有吮吸这怀抱的才真是有福的。因此伯多禄也说：\BibleQuote{所以你应摒除一切恶毒、一切诡诈、虚伪、嫉妒和诽谤，如同初生的婴孩，渴慕那圣言的奶，好使你们由此得以成长，直到救恩；既然你们已尝过主是基督。}若有人承认肉与奶有所不同，那么他们如何能避免因缺乏对自然的思考而被自己的唾沫刺穿呢？因为在冬天，当空气凝缩，阻止了内部封闭的热量散发，食物经过转化、消化和变成血，进入血管，这些血管在没有蒸发的情况下大大膨胀，并显示出强烈的搏动；因此，乳母那时奶水最充足。我们刚才也已表明，怀孕时血液通过不改变其本质的变化而变成奶，就像老年人黄发变灰一样。但在夏天，身体毛孔更张开，为食物的发汗作用提供了更大的便利，而奶水最少，因为血不充盈，全部营养也不被保留。因此，如果食物的消化结果产生血，而血又变成奶，那么血是奶的准备，就像血对于人，葡萄对于葡萄藤一样。因此，我们一出生就受到主营养的奶的哺育；一旦我们重生，我们就有幸听到安息希望的好消息，即上面的耶路撒冷，据说那里有奶和蜜如雨降下，通过有形的食物领受神圣食物的预许。\BibleQuote{食物是为肚腹，}\BibleRef{格前 6:13}正如宗徒自己所说；但这奶的滋养引领我们升天，培养天国的公民和天使歌咏团的成员。既然圣言是涌流的生命之泉，又被称为橄榄油的河流，保禄使用恰当的比喻语言，称祂为奶，补充说：\BibleQuote{我给你们喝的是奶；}\BibleRef{格前 3:2}因为我们饮用圣言，即真理的营养。的确，流质食物也称为饮料；同一物可在不同角度下同时是肉和饮料，正如奶酪是奶的固化或固化的奶；因为我不在乎挑选一个精妙的表达，只说一种物质提供两种食物。此外，对于哺乳的婴儿，仅奶就够了；它既作肉又作饮料。\BibleQuote{我，}主说，\BibleQuote{有食物吃，是你们不知道的。我的食物就是承行派遣我者的旨意。}\BibleRef{若 4:32}你们看到另一种食物，类似于奶，比喻性地代表天主的旨意。此外，祂把自己苦难的完成借喻为\Quote{杯}，那时祂必须独自饮尽。因此，对基督来说，承行父的旨意就是食物；对我们这些饮天上圣言之奶的婴孩，基督自己就是食物。因此寻求被称为吮吸；因为对于那些寻求圣言的婴孩，父的爱之乳房供应奶。\par

圣言进一步宣告自己是天上的食粮。\BibleQuote{梅瑟，}祂说，\BibleQuote{没有给你们从天降下的食粮，而是我父给你们从天上降下的真正食粮。因为天主的食粮就是那从天上降下来，并赐生命给世界的。我所要赐给的食粮，就是我的肉，是为世界的生命而赐给的。}这里要注意食粮的奥秘，因为祂称其为肉，并且作为肉，是通过火而升起的，正如麦子从朽坏和萌芽中生长出来；事实上，它为了教会的喜乐而通过火升起，如同烘烤的饼。但这一点将在论复活的章节中更清楚地解释。但既然祂说：\BibleQuote{我所要赐给的食粮，就是我的肉，}而肉被血湿润，血比喻性地称为酒，我们被告知要知道，正如饼被揉碎放入酒和水的混合物中，它抓住酒而留下水性的部分，同样，基督的肉，天上的食粮，吸收了血；也就是说，吸收那些属天的人们，滋养他们直到不朽，而只将肉体的情欲留给毁灭。\par

因此，圣言以多种方式被比喻性地描述为肉、肉身、食物、食粮、血和奶。主就是这一切，为了使信祂的我们得享喜悦。那么当我们说主的血被比喻性地表示为奶时，不要有人觉得奇怪。因为它不是也被比喻性地表示为酒么？经上说：\BibleQuote{他在酒中洗了衣服，在葡萄血中洗了袍子。}\BibleRef{创 49:11}祂说，以祂自己的圣神，祂将装饰圣言的身体；同样地，祂也将以祂自己的圣神滋养那些饥渴圣言的人。\par

血是圣言，这有义人\Person{亚伯尔}的血在天主前代为求情为证。因为血若不被视为圣言，就绝不会发出声音：因为古时的义人是新义人的预像；古时代求的血，代替新血代求。那作为圣言的血向天主呼喊，因为它暗示了圣言将要受苦。\par

再者，这肉和其中的血，通过相互的亲和，由乳汁滋润而增长。受精时种子的形成过程，是在它与经血的纯净残渣混合之后发生的。因为种子中的力量将血的物质凝结，如同凝乳酶使奶凝固，从而完成形成过程的主要部分。因为适当的混合有利于繁殖；而极端则有害，导致不育。因为当大地本身因过多的雨水淹没时，种子被冲走，而因雨水稀少则干枯；但当树液粘稠时，它保持种子，使其发芽。有些人还持有这样的假说：动物的种子实质上是血的泡沫，它被男性的自然热力搅动并摇出，变成泡沫，沉积在精索静脉中。因为\Person{阿波洛尼亚的第欧根尼}认为，由此衍生出\OriginalTerm{爱欲}{aphrodisia}一词。\par

因此，从这一切可以明显看出，人体的根本原则是血。胃中的内容物起初是乳状的，是液体的凝结；然后同样的凝结物质变成血；但在子宫中由于自然温暖的精气而形成紧凑的实体时，就成为有生命的受造物。此外，孩子出生后也由同样的血滋养。因为乳汁的分泌是血的产物；而营养的来源是乳汁；由此表明女人生了孩子，真正成为母亲，也由此她获得了有力的爱的魅力。因此，圣神借着宗徒，用主的声音神秘地说：\BibleQuote{我给你们喝的是奶。}\BibleRef{格前 3:2} 因为如果我们已重生归于基督，那重生我们的主就用祂自己的奶，即圣言，滋养我们；因为产生者应当立即为被产生者提供营养。正如重生是符合属灵的，人的营养也是属灵的。因此，在各方面、在一切事上，我们都与基督结合，通过祂的血而与祂有亲属关系，我们靠这血得救赎；也因从圣言流出的滋养而有同情；并通过祂的引导而获得永生：——\par

\Quote{在人中，抚养孩子\newline{}常比生育他们产生更强烈的爱之冲动。}\par

因此，主的同一血和奶是主苦难和教训的象征。因此，我们每一个婴孩都被允许在主里夸耀，同时我们宣告：——\par

\Quote{然而我夸耀自己出身于高贵的父亲和高贵的血统。}\par

而奶是由血经过变化而生成的，这已经很清楚；但我们还可以从羊群和牛群得知。因为这些动物，在一年中我们称为春天的时节，当空气更为潮湿，青草和草地多汁湿润时，首先充满血液，这从肿胀血管的静脉扩张可以看出；然后血中流出更丰富的奶。但在夏天，血被热力烧灼干枯，阻止了变化，因此它们的奶就少了。\par

此外，奶与水有最自然的亲和力，正如属灵的洗涤与属灵的营养一样。因此，那些吞咽一点冷水加上上述奶的人，立刻就感到益处；因为奶与水结合而防止变酸，这不是由于它们之间的任何反感，而是由于水在奶消化时友善地对待奶。\par

正如圣言与圣洗的结合，是奶与水的和谐；因为奶唯独接受水，并允许与水混合，为洁净的目的，正如圣洗为赦免罪恶。奶也自然地与蜂蜜混合，这也是为了洁净并带来甜美的营养。因为圣言与爱混合，既治愈我们的情欲，又洁净我们的罪恶；并且有话说：\par

\Quote{比蜜更甜的话语之流涌流。}\par

在我看来，这话似乎是指着圣言说的，祂就是蜜。而先知书常赞美祂\BibleQuote{比蜜和蜂房更甘甜}。\par

此外，奶与甜酒混合；这样的混合是有益的，正如苦难与杯混合以达永生。因为奶被酒凝结而分离，其中任何掺杂都被滤掉。同样，信德与受苦之人的属灵共融，将肉体的情欲如同浆液般排除，将人连同那些神圣者一起交付于永恒，使他成为不朽。\par

此外，许多人也将奶脂（称为黄油）用于灯，以此谜比喻清晰地指示圣言丰富的傅油，因为唯独祂滋养婴孩，使他们成长并光照他们。因此圣经论及主说：\BibleQuote{祂以田野的出产喂养他们；他们从磐石中吸蜜，从坚石中吸油，牛的黄油，羊的奶，还有羔羊的脂肪；}\BibleRef{申 32:13}以及随后祂赐予他们的。但预言那孩子诞生的人说：\BibleQuote{祂要吃黄油与蜂蜜。}\BibleRef{依 7:15}我心中惊异，何以有些人敢自称完美者和\OriginalTerm{诺斯替者}{gnostics}，自视高于宗徒，膨胀而自夸，因为连\Person{保禄}也自称说：\BibleQuote{这不是说我已经达到，或已经成全；而是我追求，或许可以夺得基督耶稣夺得我的目标。弟兄们，我并不以为自己已经夺得；我只顾一件事：忘记背后，努力面前的，向着标竿直跑，要得天主在基督耶稣里从高天召我来得的奖赏。}\BibleRef{斐 3:12}然而他自认为成全，因为他已从先前的生活中获得释放，并追求更好的生活，并非在知识上成全，而是渴求成全。因此他又加说：\BibleQuote{我们凡成全的人，就当存这样的心。}\BibleRef{斐 3:15}明确地将成全描述为弃绝罪恶，重生进入唯一成全者的信德，并忘记我们先前的罪。\par

% structure: p0058 source=h2 target=subsection reason=relative-heading-rank; base=h2

\subsection{第七章 导师是谁，以及关于祂的教导}

既然我们已经证明，我们众人都被圣经称为儿女；不仅如此，我们这些跟随基督的人，也被比喻称为婴孩；又因为万有的圣父独自是成全的，因为圣子在祂内，圣父也在圣子内；现在是时候按着次序来说我们的导师是谁了。\par

祂被称为耶稣。有时祂称自己为牧人，说：\BibleQuote{我是善牧。}\BibleRef{若 10:11}根据从牧人来比喻，他们引导羊群，由此理解导师，祂引导儿女——那牧养婴孩的牧人。因为婴孩是单纯的，被比喻为羊。经上说：\BibleQuote{他们都要成为一群，属于一个牧人。}\BibleRef{若 10:16}那么，引导儿女走向救恩的圣言，适当地被称为\OriginalTerm{导师}{Paedagogue}。\par

因此，圣言借欧瑟亚极其清楚地论及自己说：\Quote{我是你们的导师。}如今孝爱就是教导，是学习事奉天主，是接受真理知识的训练，是引向天堂的正确指引。而\Quote{教导}一词有多种用法。有被引导和学习的教导，有引导和教导的教导；第三是引导本身；第四是所教导的，如所吩咐的诫命。\par

那出于天主的教导，是将真理正确导向对天主的默观，并在永恒的坚忍中展现圣善的行为。\par

因此，正如将军指挥方阵，顾及士兵的安全，舵手驾驶船只，渴望拯救乘客；同样，导师出于对我们的关怀，引导儿女走上救恩的行为；并且，一般而言，凡我们按理性向天主祈求为我们成就的事，必临到信从导师的人。正如舵手并不总是顺从风向，而是有时调转船头，对抗飓风的全部力量；同样，导师从不屈服于这世界吹来的狂风，也不将孩子像船只一样交付给它们，使之在狂野放荡的生活中遭遇海难；而是乘着真理之神的顺风，牢牢握住孩子的船舵——我是指他的耳朵——直到他安全抵达天堂的锚地。\par

人们所谓的祖传习俗转瞬即逝，但天主的引导是永存的产业。\par

据说\Person{福尼克斯}是\Person{阿基里斯}的导师，\Person{阿德剌斯托斯}是\Person{克罗伊斯}子女的导师；\Person{列奥尼达}是\Person{亚历山大}的导师，\Person{瑙西托斯}是\Person{腓力}的导师。但福尼克斯好色，阿德剌斯托斯是逃亡者。列奥尼达没有遏制亚历山大的骄傲，瑙西托斯也没有改造那个酗酒的佩拉人。色雷斯人\Person{佐皮鲁斯}同样未能阻止\Person{阿尔基比亚德}的淫乱；但佐皮鲁斯是买来的奴隶，而\Person{西基诺斯}，\Person{地米斯托克利}子女的家庭教师，是个懒惰的家仆。据说他还发明了西基诺斯舞。那些被称为波斯王的导师的人，也没有逃过我们的注意；波斯王从所有波斯人中极其谨慎地选出四人，安置在他们儿子身边。但这些孩子只学射箭，成年后与姐妹、母亲、女人、妻子和无数娼妓发生性关系，像野猪一样交合。\par

但我们的导师是圣洁的天主耶稣，圣言，全人类的向导。慈爱的天主自己就是我们的导师。圣神曾在某处诗歌中论到祂说：\BibleQuote{祂在旷野中充分供给百姓。在干旱之地的炎夏干渴中，祂引领他，教导他，保护他如同眼中的瞳人，又如鹰搅动巢窝，在雏鹰以上展翅，接取雏鹰，背在两翼之上。主独自引领他们，并无外邦神与他们同在。}\BibleRef{申 32:10}显然，我认为，圣经在描述祂的引导时，展示了这位导师。\par

再者，当祂以自己位格发言时，祂承认自己是训导者：\BibleQuote{我是上主你的天主，曾领你出埃及地。}\BibleRef{出 20:2} 那么，谁有权能领入领出呢？岂不是训导者吗？祂就是向\Person{亚巴郎}显现并对他说\BibleQuote{我是你的天主，在我面前蒙悦纳；}\BibleRef{创 17:1}并且以最适宜训导者的方式，将他塑造成一个忠信的孩子，说：\BibleQuote{你要做成全的人，我要在你和我之间，以及你的后裔之间，订立我的盟约。}这便是训导者友谊的相通。祂最明显地显现为\Person{雅各伯}的训导者。祂因此对他说：\BibleQuote{看，我与你同在，在你所行的路上保护你，并将领你回到这地方；因为在我实行了我对你所说的话以前，我决不离开你。}\BibleRef{创 28:15}并且据说祂曾与祂角力。\BibleQuote{只剩下\Person{雅各伯}一人，有一个人（训导者）与他角力，直到破晓。}\BibleRef{创 32:24}这人就是那引领、带领、角力、并膏立了这斗士\Person{雅各伯}以对抗邪恶的那一位。既然圣言同时是\Person{雅各伯}的教练和人类的训导者，由此可见——据说，\Quote{祂问了}，\Quote{祂的名，并对他说：请告诉我你的名字。}祂说：\Quote{你为什么问我的名字？}因为祂为新的子民——即婴孩——保留了那新名字；当时还未命名，因为上主天主尚未成为人。但\Person{雅各伯}给那地方起名叫\Quote{天主的面}。他接着说：\BibleQuote{因为我见了天主的面，我的性命还保全了。}\BibleRef{创 32:30}此后他也被称为\Person{以色列}，因为他见了上主天主。后来，又是天主、圣言、训导者对他说道：\BibleQuote{不要害怕下到埃及去。}\BibleRef{创 46:3}请看，训导者如何跟随义人，如何膏立这斗士，教导他绊倒对手。\par

也是祂教导\Person{梅瑟}作训导者。因为主说：\BibleQuote{若有人得罪我，我要从我的书上抹去他的名字；但现在你去领这人民到我告诉你的地方去。}\BibleRef{出 32:33}这里祂是训导艺术的教师。因为实际上是主藉着\Person{梅瑟}作了旧子民的训导者；但祂亲自面对面地作新子民的训导者。\Quote{看，}祂对\Person{梅瑟}说，\Quote{我的天使要在你前面行走，}这代表了圣言的福音性和命令性权能，同时又守护着主的特权。\Quote{在我眷顾他们的日子，}\BibleRef{出 32:33}祂说，\Quote{我要将他们的罪归到他们身上；}即在我作为审判者坐席的日子，我要按他们的罪施行报应。因为同一位是训导者也是审判者，审判那些不顺从祂的人；而仁爱的圣言不会默默放过他们的过犯。祂谴责他们，好使他们悔改。因为\Quote{主愿意罪人悔改，胜过他的死亡。}让我们如同婴孩，听到别人的罪过，因畏惧警告而避免类似的过犯，免得我们也遭受同样的苦楚。那么，他们所犯的罪是什么呢？\BibleQuote{因为他们一怒杀人，任性砍断牛腿。愿他们的愤怒受诅咒。}\BibleRef{创 49:6}那么，有谁比祂更慈爱地训练我们呢？从前，旧子民拥有旧盟约，法律以恐惧训导人民，圣言是一位天使；但对于这新鲜的新子民，也赐给了新盟约，圣言显现了，恐惧转化为爱，那位神秘的天使诞生了——\Person{耶稣}。因为同一位训导者那时说：\BibleQuote{你要敬畏上主你的天主；}\BibleRef{申 6:2}但对我们，祂劝诫说：\BibleQuote{你要爱主你的天主。}\BibleRef{玛 22:37}因此，这也是对我们的命令：\Quote{停止你们自己的工作，离开你们的旧罪；}\Quote{学习行善；}\Quote{离恶行善；}\Quote{你爱了正义，恨了不义。}这就是用旧文字写成的我的新盟约。因此，圣言的新颖不可成为指责的理由。但主也在耶肋米亚中说：\BibleQuote{不要说我是少年；我还没有在母腹中形成你以前，我就认识了你；在你还没有出母胎以前，我就圣化了你。}\BibleRef{耶 1:7}这样的暗示，预言能向我们发出；我们在天主眼中，于世界奠基之前就被预定要得到信德；但现在我们成为婴孩，因着天主旨意最近得以实现，我们如今诞生于召叫和救恩。因此祂又说：\BibleQuote{我立了你作万民的先知，}\BibleRef{耶 1:5}说他必须说预言，好使\Quote{少年}这个称呼不致成为那些被称为婴孩之人的羞辱。\par

如今，法律是藉着圣言经由梅瑟所赐下的古代恩宠。因此圣经也说：\BibleQuote{法律是藉着梅瑟赐下的}，\BibleRef{若 1:17} 不是由梅瑟，而是由圣言，并通过祂的仆人梅瑟。因此它只是暂时的；但永恒的恩宠和真理却是藉着耶稣基督而来的。请注意圣经的表达：关于法律，只说是\BibleQuote{赐下的}；但真理既是圣父的恩宠，便是圣言的永恒工程；它不是被说成\Emph{赐下的}，而是藉着耶稣\Emph{成为的}，\BibleQuote{万物是藉着祂而有的}。\BibleRef{若 1:3} 因此，梅瑟先知性地让位给完美的导师圣言，预言了导师的名号和职分，将服从的命令托付给百姓，并把导师摆在他们面前。\BibleQuote{一位先知}，他说，\BibleQuote{如同我一样，上主你的天主要从你们兄弟中给你兴起一位}，指着农的儿子若苏厄暗示性地指出天主子耶稣；因为法律中所预言的耶稣这名是基督的预像。因此，他为了人民的益处接着说道：\BibleQuote{你们要听从祂}，\BibleRef{申 18:15} 并且，\BibleQuote{那人不听从这先知的}，\BibleRef{申 18:19} 祂便要惩罚他。他预言这样一位名号属于导师，祂是救恩的创造者。因此，预言将一根棍杖赋予祂，一根纪律、统治、权柄的棍杖；为让那不被劝导之言治愈的，被威胁治愈；不被威胁治愈的，被棍杖治愈；不被棍杖治愈的，被火吞灭。\Quote{要生出一根棍杖}，经上说，\Quote{从叶瑟的根茎生出}。\par

请看导师的关怀、智慧和权能：\BibleQuote{祂不按意见审判，也不按报告审判；祂要为谦卑者伸张正义，谴责地上的罪人。} 藉着达味：\BibleQuote{上主教导了我，祂教导了我，没有将我交给死亡。} 因为受主的惩戒和教导，就是脱离死亡。同一位先知又说：\BibleQuote{你必用铁杖治理他们。} 同样，宗徒在致格林多人书，受感动说：\BibleQuote{你们愿意怎样呢？要我带着棍杖到你们那里去呢？还是怀着爱，以温柔的精神呢？} \BibleRef{格前 4:21} 还有，\BibleQuote{上主必从熙雍伸出权杖}，祂藉另一位先知说。而同一根教鞭，\BibleQuote{你的棍杖和木杖都安慰了我}，另一个人说。这就是导师的权能——神圣的、抚慰的、救赎的。\par

% structure: p0071 source=h2 target=subsection reason=relative-heading-rank; base=h2

\subsection{第八章　驳斥那些认为正义并非善的人}

这时有些人起来说，主因着棍杖、威胁和恐惧，不是善的；他们显然误解了圣经的话：\BibleQuote{敬畏上主的人必归向自己的心}；\BibleRef{德 21:6} 最严重的是，他们忘记了祂的爱，即祂为我们成为人。为此，先知更恰当地祈祷说：\BibleQuote{求祢记念我们，我们不过是尘土}；意思是：同情我们；因为祢亲自从受苦经验中知道肉身的软弱。因此，主导师是最善的、无可指摘的，祂因着极其伟大的爱而同情每个人的本性。\BibleQuote{因为上主所恨的，什么也没有}。\BibleRef{智 11:24} 因为祂确实不恨任何事物，却又希望祂所恨的存在。祂也不希望任何事物不存在，却又成为祂不希望存在的事物的原因。祂也不希望任何已存在的事物不存在。那么，如果圣言恨什么，祂就不希望它存在。但没有任何事物的存在原因不是由天主提供的。所以，没有什么被天主恨，也没有被圣言恨。因为两者是一体——就是天主。因为祂说过：\BibleQuote{在起初，圣言就在天主内，圣言就是天主}。\BibleRef{若 1:1} 因此，如果祂不恨祂所造的任何一个，那么祂就爱它们。祂尤其，且有理地，更爱人类——祂所造的最尊贵的受造物，一个爱天主的存有。所以天主是爱；因此圣言是爱。\par

但爱任何事物的人都愿意对它行善。而行善者必须在各方面优于不行善者。但没有什么比善更好。因此，善行善。而天主被认为是善的。所以天主行善。善，因着它的本质是善，不行别的事，只行善。所以天主行一切善。祂对人行善而不关心他，这是不可能的；祂不会\Emph{关心}他而不\Emph{照顾}他，也是不可能的。因为有意行善的，优于无意行善的。但没有什么比天主更好。有意行善，无非就是照顾人。所以天主关心人，并照顾他。祂在实践中表明这一点，藉着圣言教导人，圣言是天主爱人的真正合作者。但善之所以被称为善，不是因为拥有德性；正如正义之所以被称为善，也不是因为它拥有德性——因为它本身就是德性——而是因为它在自身内并藉着自身是善的。\par

另一方面，有用之物被称为善，不是因其悦乐，而是因其行善。因此，所有正义之物，既是善的，作为德性本身和因其自身而被欲求，而非因其提供快乐；因为正义不是为讨好而审判，而是按各人所应得的给予。而有益之物紧随有用之物。因此，正义具有与善被考察的所有方面相对应的特性，同样拥有等同的属性。而具有等同属性的事物彼此相等且相似。所以正义是善的。\par

\Quote{那么，}他们说，\Quote{如果主爱人，且是善的，祂为何发怒并惩罚？}因此我们必须尽可能简洁地论述这一点；因为这种处理方式有益于儿童的正当训练，充当一种必要的帮助。因为许多情欲都能通过惩罚、严厉诫命的灌输以及某些原则的教导来治愈。因为责斥仿佛是灵魂情欲的手术；而情欲仿佛是真理的脓肿，必须用责斥的手术刀切开。\par

责斥如同用药，化解情欲的硬结，清除生活中淫秽的污秽；此外，还削减骄傲的赘疣，使病人恢复健康而真实的人性状态。\par

告诫仿佛是患病灵魂的养生法，规定它必须摄取什么，禁止它不应接触什么。而这一切都趋向于救恩和永恒的健康。\par

再者，军队的将军通过处以罚金、用锁链和极端耻辱实施体罚，有时甚至以死刑惩罚个人，这都为了好处，是为了告诫他手下的军官。\par

同样，我们伟大的将军、圣言，宇宙的统帅，通过告诫那些抛弃祂法约束的人，使他们从奴役、错误和敌人的囚禁中获得释放，从而和平地引导他们进入神圣的公民和谐。\par

因此，正如除了说服性话语外，还有规劝和安慰的形式，同样，除了赞美外，还有指责和责难的形式。后者构成了责备的艺术。而责备是善意的标志，而非恶意。因为朋友和非朋友都会责斥，但敌人是出于轻蔑，朋友则是出于善意。所以，主并非出于仇恨而责备人；因为祂为我们受苦，而祂本可以因我们的过错毁灭我们。因为导师，由于祂是善的，以精湛的技巧通过训斥巧妙地转入责备；用锋利的话语如同鞭子一般唤醒心灵的迟钝。然后祂又转而尽力劝诫同一些人。因为那些不受赞美诱导的人，被责备所激励；而那些不被责备唤向救恩犹如死亡的人，被指责唤醒到真理面前。\Quote{因为智慧的鞭责与管教常在一切时候。}\BibleQuote{教导愚人如同粘连陶片；使绝望的愚钝者敏锐如同使泥土有感觉。}\BibleRef{德 22:6}因此祂明确地加上说，\Quote{将沉睡者从深睡中唤醒，}这在一切事物中最像死亡。\par

此外，主非常清楚地显示祂自己，当祂比喻性地描述祂那多样且多方面有益的教育时，祂说：\BibleQuote{我是真葡萄树，我父是园丁。}然后祂加上：\BibleQuote{凡在我身上不结果实的枝子，祂就剪去；凡结果实的，祂就修剪，使它结更多的果实。}\BibleRef{若 15:1}因为不被修剪的葡萄树会长成木头。人也是如此。圣言——刀——清除放荡的枝条；迫使灵魂的冲动结果实，而非沉溺于私欲。现在，对罪人的责斥以他们的救恩为目的，话语和谐地适应每个人的行为；时而收紧，时而放松绳索。因此梅瑟非常明确地说：\BibleQuote{你们要鼓起勇气：天主已临近来试验你们，使祂的敬畏常存在你们中间，使你们不犯罪。}\BibleRef{出 20:20}而柏拉图，从这源泉学到，优美地说：\Quote{因为所有受惩罚的人实际上都受到善待，因他们受益；因为公正受罚者的精神得到改善。}如果那些受纠正的人从正义手中得到好处，且根据柏拉图，正义被承认为善，那么恐惧本身也产生益处，并且被发现对人有好处。\BibleQuote{因为敬畏主的灵魂必生活，因为他们的希望在于拯救他们的那一位。}\BibleRef{德 34:14}而这同一位施加惩罚的圣言就是审判者；关于祂，依撒意亚也说：\BibleQuote{主把祂归于我们的罪过，}\BibleRef{依 53:6}显然作为罪的纠正者和改造者。因此，唯独祂能赦免我们的不义，祂被圣父指派为我们众人的导师；唯独祂能分辨不顺服与顺服。当祂警告时，祂显然不愿意施行恶来实现祂的警告；而是通过激发人的敬畏，切断通往罪的道路，并显示祂对人的爱，仍然延迟，并宣告如果他们继续犯罪将受何苦，并且不像蛇，一旦咬住猎物就吞噬它。\par

\BibleQuote{因为我的箭，}祂说，\BibleQuote{将使他们灭绝；他们将被饥饿耗尽，被飞鸟吞食；并有不治的角弓反张。我要打发野兽的牙齿到他们身上，连同地上爬行的蛇的忿怒。外有刀剑使他们丧子；内室之中有恐怖。} \BibleRef{申 32:23}因为天主并非像有些人所想的那样发怒；祂却时常抑制，总是劝勉人类，并指示当行之事。这是良策：以恐吓阻止我们犯罪。\BibleQuote{敬畏主驱除罪恶；不敬畏者不能成义，} \BibleRef{德 1:21}圣经这样说。天主并非出于忿怒而施行惩罚，而是为了正义的目的；因为正义不应因我们而被忽略。我们中每一个犯罪的人，都凭自己的自由意志选择了惩罚，过错在于选择者。天主是无辜的。\BibleQuote{但如果我们的不义彰显了天主的正义，我们该说什么呢？难道降怒的天主不义吗？断乎不是！} \BibleRef{罗 3:5}因此，祂警告说：\BibleQuote{我要磨利我的剑，我的手要执掌审判；我要向我的敌人复仇，报复恨我的人。我要使我的箭饮血，我的剑要吞食血肉，就是被伤者的血。} \BibleRef{申 32:41}由此可见，那些不与真理为敌、不恨圣言的人，不会恨自己的救恩，而会逃避为敌的惩罚。\BibleQuote{智慧之冠，}正如\BibleBook{}{智慧篇}说，\BibleQuote{是敬畏主。} \BibleRef{德 1:18}因此，主藉着先知亚毛斯非常清楚地揭示了祂的行事方式，说：\BibleQuote{我倾覆了你们，如同天主倾覆了索多玛和哈摩辣；你们好像从火中抽出的柴；但你们仍不归向我——主说。} \BibleRef{亚 4:11}\par

请看天主如何因祂的爱善之心而寻求悔改；并藉着祂默默警告的计划，显示祂对人的爱。\BibleQuote{我要转面，}祂说，\BibleQuote{离开他们，并指示他们将要遭遇的事。} \BibleRef{申 32:20}因为天主的面容所视之处，就有平安与喜乐；而祂转面之处，便引入邪恶。因此，主不愿看邪恶之事，因为祂是善的。但当祂转面时，邪恶便因人的不信而自发产生。\BibleQuote{因此，请看}，保禄说，\BibleQuote{天主的慈爱与严厉：对跌倒的人，是严厉；对你，却是慈爱，只要你存留在祂的慈爱中，} \BibleRef{罗 11:22}即在于对基督的信德。\par

如今，恨恶邪恶属于善人，因其本性为善。因此，我承认祂惩罚悖逆者（因为惩罚对受罚者有益处，是对顽固臣民的纠正），但我不承认祂想要报复。复仇是为报复者的利益而对邪恶施行的报应。祂教训我们\BibleQuote{要为那迫害我们的人祈祷}，\BibleRef{玛 5:44}故祂不会要我们复仇。天主是善的，众人皆愿承认；而同一的天主是正义的，我不需多言，只需引述主的福音言辞：祂论及祂为一，说：\BibleQuote{愿众人都合而为一！父啊！就如祢在我内，我在祢内，使他们也在我们内合而为一，为叫世界相信是祢派遣了我。祢赐给我的光荣，我已赐给了他们，为叫他们合而为一，如同我们原为一体一样。我在他们内，祢在我内，使他们完全合而为一。} \BibleRef{若 17:21}天主是唯一的，超越唯一本身，超越元一。因此，\Quote{祢}这一词，带有指示性的强调，指出那唯一真实存在者，\Quote{今在、昔在、将来永在的}，在这三个时间分段中，那唯一的名字（\OriginalTerm{那自有者}{\GreekText{ὀ} \GreekText{ὤν}}）；\BibleQuote{自有者}，\BibleRef{出 3:14}有其位置。那唯一的天主也是唯一真正正义的，我们的主在福音中亲自见证，说：\BibleQuote{父啊！祢赐给我的人，我愿他们在那里与我同在，为叫他们享见祢赐给我的光荣，因为祢在创世之前就爱了我。公义的父啊！世界没有认识祢，我却认识祢，这些人也知道是祢派遣了我。我已将祢的名宣示给他们，我还要宣示。} \BibleRef{若 17:24}这就是那\BibleQuote{追讨父亲的罪于子孙身上，直到恨祂的人；对爱祂的人却施仁慈，} \BibleRef{出 20:5}。因为祂将一些人安置\Quote{在右边，另一些在左边}，作为父而被理解为善，被称为祂所是的那唯一——\BibleQuote{善}；\BibleRef{玛 19:17}但作为在父内的圣子，即祂的圣言，从彼此的关係中，能力的名字由爱的平等来衡量，被称为正义。\BibleQuote{祂必审判，}祂说，\BibleQuote{人按他的行为，} \BibleRef{德 16:12}——这是公义的天平，天主藉着耶稣让我们认识正义的面貌，我们藉着祂，如同藉着平衡的天平，认识天主。\BibleBook{}{智慧篇}也明确地说：\BibleQuote{因为怜悯和忿怒都在祂那里，祂是这两者的唯一主，}是宽恕的主，并照祂丰厚的怜悯倾注忿怒。\BibleQuote{祂的责斥也是如此。} \BibleRef{德 16:12}因为怜悯和责斥的目的，是为了被责斥之人的救恩。\par

如今，我们主耶稣基督的天主和圣父是良善的，圣言本身将再次证明：\Quote{因为祂以仁慈对待忘恩和作恶的人；}并且，当祂说：\BibleQuote{你们应当慈悲，如同你们的圣父是慈悲的一样。}\BibleRef{路 6:35}此外，祂还清楚地说：\BibleQuote{除了我在天上的圣父外，没有一个是良善的。}\BibleRef{玛 19:17}除此之外，祂又说：\BibleQuote{我的圣父使太阳照耀所有的人。}\BibleRef{玛 5:45}在此应注意到，祂宣告祂的圣父是良善的，也是造物主。造物主是公义的，这是没有争议的。祂又说：\Quote{我的圣父降雨给义人，也给不义的人。}就降雨而言，祂是水和云的造物主。就降雨给所有人而言，祂公平且正直地持守平衡。而作为良善者，祂同样对待义人和不义的人。\par

那么，我们非常清楚地得出祂是独一同一的天主。因为圣神曾歌唱：\Quote{我要仰观诸天，你手所造的；}以及\Quote{创造诸天的那一位居住在诸天；}和\Quote{天堂是你的宝座。}主在祈祷中说：\BibleQuote{我们的天父，你是在天上的。}\BibleRef{玛 6:9}诸天属于那创造世界的祂。所以，无可争议，主是造物主之子。如果万有之上的造物主被承认为公义者，而主是造物主之子，那么主就是公义者之子。因此保禄也说：\BibleQuote{但如今天主的义德在法律之外已显示出来；}\BibleRef{罗 3:21}并且，为使你更好地认识天主：\Quote{就是天主的义德，因对耶稣基督的信德，赐予所有相信的人，没有分别。}此外，为见证真理，他稍后补充说：\Quote{借着天主的宽容，为显示祂是公义的，并且耶稣是使那出于信德的人成义者。}而他知道公义的就是良善的，这显现在他的话语中：\BibleQuote{所以法律是圣的，诫命也是圣的、公义的、良善的，}\BibleRef{罗 7:12}用这两个名称表示同一权能。但\Quote{没有一个是良善的，}除了祂的圣父。那么，正是祂的这位圣父，虽是独一，却借许多权能显示出来。这话的含义是：\BibleQuote{没有人认识圣父，}\BibleRef{路 10:22}祂在圣子来临之前就是万有。所以清楚明白，万有的天主是独一的良善、公义的造物主，而圣子在圣父内，愿光荣归于祂，直到永远。阿们。但出于关怀的责备与救恩的圣言并不矛盾。因为这是天主对人之圣爱的良药，由此羞赧的红晕浮现，对罪过的羞愧随之而来。若人应受责难，则亦须责备；当时间到来时，要刺伤麻木的灵魂，不是致命地，而是拯救性地，借着轻微的痛苦确保免于永死。\par

祂教导中所显示的智慧是伟大的，祂为了救恩而施行的方式是多样的。因为导师为良善者作证，召唤被召的人走向更善；劝阻那些急于行恶的人尝试，并劝勉他们转向更善的生活。因为当一方被证实，另一方并非没有见证；从见证而来的恩宠是极大的。此外，愤怒的情感（若可以这样称呼祂的告诫），是充满对人之爱的，天主为了人的缘故而屈就于情感；正是为了人，天主的圣言成了人。\par

% structure: p0088 source=h2 target=subsection reason=relative-heading-rank; base=h2

\subsection{第九章 同一权能的恩赐：既施恩又公正惩罚；以及圣言的教导方式}

因此，全能的教导人类者——神圣的圣言，运用一切智慧的资源，致力于拯救子女：劝诫、斥责、责备、训诲、指责、警告、医治、应许、施恩；如同用许多缰绳，遏制人类非理性的冲动。简言之，主对待我们如同我们对待自己的子女。\Quote{你若有了儿子，就该管教他们，}是智慧书的劝导，\BibleQuote{从年轻时就要使他们顺服。你若有了女儿，就该关心她们的身体，不要对她们显露笑容。}\BibleRef{德 7:23}——虽然我们极其爱自己的子女，无论是儿子还是女儿，胜过其他一切。因为那些只为讨人欢心而对人说话的人，对他们的爱很少，因为他们不使人痛苦；而那些为使人得益而说话的人，虽然暂时带来痛苦，却永远使人受益。主所着眼不是眼前的快乐，而是将来的喜乐。\par

现在，让我们借着先知的见证，继续思考祂慈爱训导的方式。\par

劝诫，乃是慈爱的纠正，并产生理解。这就是导师在其劝诫中所行的，正如祂在福音中所说：\Quote{我多次愿意聚集你的子女，有如母鸡把自己的幼雏聚集在翅膀底下，但你偏不愿意！}\BibleRef{玛 23:37} 再者，圣经也劝诫说：\Quote{他们与木头和石头行淫，并向巴耳焚香。}这乃祂大爱的极明证：虽深知那踢跳挣脱之民的恬不知耻，祂仍劝其悔改，并藉厄则克耳说：\Quote{人子，你住在蝎子中间；尽管如此，你要对他们讲话，或许他们会听。}\BibleRef{则 2:6} 又对梅瑟说：\Quote{你去告诉法郎，要他遣散我的百姓；但我知道他不会遣散他们。}\BibleRef{出 3:18} 因为祂同时显示了两者：祂的神性在于预知将要发生的事，祂的爱在于给灵魂的自由决断提供悔改的机会。祂也藉依撒意亚劝诫，因祂关心百姓，说：\Quote{这百姓用嘴唇尊敬我，他们的心却远离我。}接下来是责备的斥责：\Quote{他们徒然敬拜我，将人的规条当作道理教训人。}\BibleRef{依 29:13} 此处祂的慈爱在显出他们的罪的同时，也并列显示了救恩。\par

斥责是因卑劣之事而责备，引导人趋向高尚。这由耶肋米亚显明：\Quote{他们成了发情的雄马，各向邻舍的妻子嘶鸣。难道我不该为这些事报复吗？——上主说：难道我的灵魂不该对这样的民族报仇吗？}\BibleRef{耶 5:8} 祂处处穿插恐惧，因为\Quote{敬畏上主是智慧的开端。}\BibleRef{箴 1:7} 再者，藉欧瑟亚祂说：\Quote{难道我不该报复他们吗？因为他们自己与妓女混杂，并与领过奥秘的人一同献祭；那明悟的民竟拥抱妓女。}祂指出他们的过犯更为明显，表明他们既明白而故意犯罪。明悟是灵魂的眼；因此以色列也意为\Quote{看见天主的人}——即明悟天主的人。\par

抱怨是对被视为藐视或忽略者之责备。祂用此形式时，藉依撒意亚说：\Quote{诸天，请听！大地，请侧耳！因为上主说了：我生育并抚养了儿女，他们却悖逆了我。牛认识自己的主人，驴认识自己主人的槽；但以色列却不认识我。}\BibleRef{依 1:2} 因为若那认识天主的人却不认识上主，我们岂不视为可怕？而牛和驴，愚钝无知的牲畜，尚且认识喂养它们的主人，以色列竟比这些更无理性？又藉耶肋米亚多端抱怨百姓后，祂加上：\Quote{他们离弃了我——上主说。}\par

痛斥是带侮辱的斥责，或斥责性的责备。导师在依撒意亚中使用此方式，说：\Quote{祸哉，你们这些悖逆的儿女！上主这样说：你们策划，却非出于我；你们结盟，却非出于我的圣神。}\BibleRef{依 30:1} 祂在每一情形中都运用恐惧这一极苦之媒染剂抑制百姓，同时又引导他们归向救恩；正如羊毛在染色过程中，常先用媒染剂处理，以预备承受不易褪色的颜色。\par

指摘是将罪过显明出来，摆在人面前。祂采用这种教导形式，因其最为必要，因许多人的信德软弱。因为祂藉依撒意亚说：\Quote{你们离弃了上主，并且激怒了以色列的圣者。}\BibleRef{依 1:4} 祂也藉耶肋米亚说：\Quote{诸天为此震惊，大地极其战栗。因为我的百姓作了两件恶事：他们离弃了我这活水的泉源，却为自己凿出破裂的水池，不能蓄水。}\BibleRef{耶 2:12} 又藉同一先知说：\Quote{耶路撒冷犯了罪，因此成了动荡。凡先前荣耀她的，都因见了她的卑贱而羞辱她。}\BibleRef{哀 1:8} 祂在撒罗满的劝言中使用了指摘那辛辣而刺耳的话语，暗指祂教导中特有的如父母爱子之情：\Quote{我儿，不要轻视上主的惩戒，也不要因祂的责备而灰心；因为上主所爱的，祂必惩戒，祂鞭策祂所收纳的每个儿子；}\BibleRef{箴 3:11} \Quote{因为犯罪的人逃避责备。}\BibleRef{德 32:21} 因此，圣经又说：\Quote{让义人指摘并纠正我，但不要让罪人的油膏我的头。}\par

\OriginalTerm{使醒悟}{\GreekText{φρένωσις}}是使人思考的责备。祂也未避免采用这种教导形式，而藉耶肋米亚说：\Quote{我呼喊要到几时，你们却不听？你们的耳朵未受割损。}\BibleRef{耶 6:10} 有福的忍耐！再者，藉同一先知说：\Quote{所有外邦人都未受割损，但这百姓心未受割损；}\BibleRef{耶 9:26} \Quote{因为这百姓是悖逆的儿女，}祂说，\Quote{心中没有信德。}\BibleRef{依 30:9}\par

谴责是严厉的斥责。祂在福音中使用此类：\Quote{耶路撒冷，耶路撒冷，你常杀害先知，用石头砸死那些奉派到你这里来的人！}重复地名使斥责更有力。因为那认识天主的人，怎能迫害天主的仆人呢？因此祂说：\Quote{你们的家园将成荒场；因为我告诉你们：从今以后，你们不得再见我，直到你们说：奉主名来的，当受赞美。}\BibleRef{玛 23:37} 因为若你们不接受祂的爱，你们将认识祂的大能。\par

痛斥是激烈的言辞。祂以痛斥为药，藉依撒意亚说：\Quote{祸哉，犯罪的民族，悖逆的儿女，满盈罪恶的百姓，邪恶的种子！}\BibleRef{依 1:4} 在若望福音中祂说：\Quote{蛇类，毒蛇的孽子。}\par

控告是对作恶者的谴责。祂藉\Person{达味}使用这种教导方式，说：\Quote{我素不认识的百姓竟事奉了我，他们一听见我的名声就服从了我。外邦人的子孙向我撒谎，他们偏离自己的道路。}又藉\Person{耶肋米亚}说：\Quote{我给了她休书，而背约的犹大却不畏惧。}\BibleRef{耶 3:8}又说：\Quote{以色列家藐视我，犹大家却向主说谎。}\BibleRef{耶 5:11}\par

哀叹命运是一种潜在的谴责，并以巧妙的方式在面纱下施行救恩。祂藉\Person{耶肋米亚}使用这种方式：\Quote{先前满有人民的城，现在何竟独坐！她素来管辖诸地，如今竟如寡妇；她成了进贡的，在夜间痛哭流泪。}\BibleRef{哀 1:1}\par

斥责是责备性的谴责。神圣导师藉\Person{耶肋米亚}使用了这种帮助，说：\Quote{你却有娼妓之额，在众人面前不知羞耻；你也不呼我到你家中——我是你的父亲，是你童贞的主。}\BibleRef{耶 3:3}\Quote{又是一位美丽雅致、精于邪术的娼妓。}\BibleRef{鸿 3:4}在娴熟地以娼妓这一污名称呼童女之后，祂随即又以羞耻充满她，召她回返尊贵的生活。\par

义怒是正当的责备；或是因行为高过正义而责备。祂藉\Person{梅瑟}以此方式教导，说：\Quote{有毛病的儿女，弯曲悖逆的世代，你们就这样报答主吗？这百姓愚昧，无智慧。祂岂不是你的父亲，创造了你？}\BibleRef{申 32:5}祂又藉\Person{依撒意亚}说：\Quote{你的官长是悖逆的，与盗贼同伙，喜爱礼物，追求酬报，不为孤儿伸冤。}\BibleRef{依 1:23}\par

最终，祂所施行的引发敬畏的体系是救恩的根源。而拯救是良善的特权：\Quote{主的仁慈遍及一切血肉，祂如牧人牧放自己的羊群，惩戒、纠正、教导。祂怜悯那些接受祂教导的人，以及那些急切寻求与祂结合的人。}\BibleRef{德 18:13}带着这样的引导，祂护卫了那六十万步卒，他们聚集在内心的顽硬中，祂鞭打、怜悯、击打、医治，以慈悲和管教：\Quote{因为祂的仁慈何其伟大，祂的责备也何其大。}\BibleRef{德 16:12}因为不犯罪固然高尚，但罪人悔改也是好的；正如常保健康最好，但从病中痊愈也是好的。因此祂藉\Person{撒罗满}命令：\Quote{用杖责打你的儿子，好救他的灵魂免于死亡。}\BibleRef{箴 23:14}又说：\Quote{不可不责打你的儿子，用杖管教他，他必不致死亡。}\BibleRef{箴 23:13}\par

因为责备与斥责——正如原词所暗示的——是灵魂的鞭笞，惩罚罪恶，阻止死亡，并引导那些被放纵所掳的人走向自律。同样，\Person{柏拉图}也知道责备是改造的最大力量，是最纯净的净化，符合上面所说，他观察到\Quote{那最不洁的人是无知而卑劣的，因为他从未在他应当最纯洁、最美好——以成为真正幸福之人——的方面受到责备。}\par

因为如果统治者不是善行的恐怖，那么那位本性为善的天主，怎会是不犯罪者的恐怖呢？\Quote{你若作恶，就该惧怕，}\BibleRef{罗 13:3}宗徒说。因此，宗徒本人也处处效法主的榜样，对各教会使用严厉的言辞；他自知胆壮，也深知听者的软弱，便对迦拉达人说了：\Quote{我成了你们的仇敌，因为我把真理告诉你们吗？}\BibleRef{迦 4:16}同样，健康的人不需要医生，强壮时不需要；但患病者需要他的医术。同样，我们这些在生活中患了可耻的贪欲和应受责备的放纵，以及其他激情炎症的人，需要救主。祂不仅施予温和的药物，也施予强烈的药物。因此，恐惧的苦根制止了我们罪恶的溃烂。因此，恐惧即使苦涩，也是有益的。我们是病人，真正需要救主；是迷失者，需要引领者；是盲人，需要引向光明者；是口渴者，\Quote{需要生命之泉，凡饮了这泉的，将永不再渴；}\BibleRef{若 4:13}是死者，需要生命；是羊，需要牧人；是孩子，需要导师。全人类都需要\Person{耶稣}，免得我们继续顽固地犯罪到底，从而堕入审判，而是被从糠秕中分离出来，收进父亲的谷仓。\Quote{因为簸箕在主手中，祂要簸净禾场，把麦子收进仓里，把糠秕用不灭的火烧尽。}\BibleRef{玛 3:12}你若愿意，可以学习至圣牧者和导师——全能而慈父般的圣言——的至高智慧，祂将自身比喻为羊群的牧者。祂也是孩子们的导师。因此祂藉\Person{厄则克耳}说，将话语指向长老们，为他们描绘祂智慧关切的救恩图画：\Quote{我必缠裹瘸腿的，医治有病的，领回迷失的，并在我的圣山上牧养他们。}这就是善牧的应许。\par

喂养我们这些孩子，如同羊群。是的，主啊，以您的义，就是您自己的牧场，充满我们；是的，导师啊，在您圣山——教会——上喂养我们，这教会高耸入云，触及天庭。祂说：\BibleQuote{我将作他们的牧者}\BibleRef{则 34:14}，并要靠近他们，如同衣服贴近皮肤。祂愿意拯救我的肉身，用不朽的长袍包裹它，并傅油了我的身体。祂说：\BibleQuote{他们必呼求我，我要说：我在这里}\BibleRef{依 58:9}。主啊，您比我所预料的更早听到了。主说：\BibleQuote{他们逾越时，必不滑跌}\BibleRef{依 43:2}。因为我们这些前往永生的人不会陷入朽坏，因为祂必扶持我们。祂既如此说，也如此愿意。这就是我们的导师，公义而良善。祂说：\BibleQuote{我来不是为受服事，而是为服事人}\BibleRef{玛 20:28}。因此，福音记载祂\BibleQuote{疲倦}\BibleRef{若 4:6}，因为为我们劳苦，并应许\BibleQuote{为众人交出自己的生命作赎价}\BibleRef{玛 20:28}。唯独这样做的祂，祂才承认是好牧人。因此，那为我们舍弃最伟大的恩赐——祂自己的生命——的祂，是何等慷慨；祂极其仁慈，热爱世人：当祂本可为主时，却愿意成为人的兄弟；祂如此良善，以至为我们而死。\par

此外，祂的公义呼喊说：\BibleQuote{你若正直地到我这里来，我也必正直地到你那里去；但你若行走弯曲，我也必弯曲而行，万军的上主说。}意指弯曲的道路就是罪人的惩罚。因为那条正直而自然的道路，由耶稣名字中的\Emph{\OriginalTerm{伊奥塔}{Iota}}所指示的，就是祂的良善，这良善对那些听信的人坚定而可靠：\BibleQuote{我呼唤，你们却不听从，上主说；你们藐视我的劝告，不理会我的责备。}\BibleRef{箴 1:24}因此上主的责备极其有益。达味也论到他们说：\BibleQuote{一个乖僻悖逆的世代，一个心不正直、心神对上主不忠的世代：他们不遵守天主的盟约，不愿行走在祂的法律中。}\par

这些就是激怒的原因，审判者来惩罚那些不愿选择良善生活的人。因此，后来祂更严厉地攻击他们，为要尽可能将他们从冲向死亡的冲动中拖回。祂便藉达味说出这威胁的最明显原因：\Quote{他们不信祂奇妙的作为。祂击杀他们时，他们才寻求祂，回转并清早寻求天主；他们记起天主是他们的帮助者，至高天主是他们的救赎者。}因此祂知道他们因恐惧而回转，却轻视祂的爱：因为大多数情况下，那始终温和的良善常被藐视；但那以正义之爱的敬畏来劝诫的，却受尊敬。\par

恐惧有两种：一种带有敬畏，如公民对好统治者，我们对于天主，以及正直的子女对于父亲那样。\BibleQuote{一匹未经驯服的马会变得难驾驭，一个任其自便的儿子会变得鲁莽。}\BibleRef{德 30:8}另一种恐惧带有憎恨，如奴隶对严酷的主人，以及希伯来人所感受的，他们使天主成了主人，而非父亲。就虔诚而言，自愿和自发的与强迫的相差甚远，甚至完全不同。经上说：\Quote{因为祂是仁慈的；祂会医治他们的罪，不毁灭他们，并完全转离祂的怒气，不燃起全部的怒火。}看，导师的公义如何以责备彰显，天主的良善如何以怜悯行事。因此达味——即藉着他的圣神——将两者包含在内，歌颂天主本身说：\Quote{正义与审判是祂宝座的根基：仁慈与真理将行在祢面前。}他宣称，审判与行善属于同一权能。因为两者都在于权能，审判将公义与其对立面分开。那真正是天主者既是公义又是良善；祂自己是万有，万有也是祂；因为祂是天主，唯一的天主。\par

因为正如镜子对于丑陋的人并非邪恶，因它向他显示他的本相；又如医生对于病人并非邪恶，因他告诉他发烧——医生并非发烧的原因，只是指出发烧——同样，那责备的人对灵魂患病的人也并非怀有恶意。因为祂不将过犯加在他身上，只是指出那里存在的罪，为使他远离类似的行为。因此，天主因自身是良善的，也因我们而公义，祂之所以公义，是因为祂良善。祂的公义由祂自己的圣言从那里，从上面，圣父所在之处，显示给我们。因为在成为创造者之前，祂就是天主；祂是良善的。因此祂愿意成为创造者和圣父。那一切爱的本质就是正义的源头——也是祂发出太阳并派遣祂自己儿子的原因。祂首先宣告那来自天上的良善正义，当祂说：\BibleQuote{除了圣父，没有人认识子；除了子，也没有人认识圣父。}\BibleRef{路 10:22}这种相互和彼此的认识是原始正义的象征。然后，正义降临到人间，既在文字中，也在身体内，在圣言和法律中，约束人类走向救赎的悔改；因为那是良善的。然而，你不服从天主吗？那就责备你自己吧，是你将审判者引到自己身上。\par

% structure: p0111 source=h2 target=subsection reason=relative-heading-rank; base=h2

\subsection{第十章 同一的天主，藉同一的圣言，以威胁阻止人犯罪，以劝勉拯救人类}

如果，那么，我们已经表明，对人性采取严厉处理的计划是美善且有救恩性的，并且是圣言必然采用的，有利于悔改和防止罪恶；那么我们现在应当依次来看圣言的温和。因为祂已被证明是公义的。祂将祂自己那些召叫我们走向救恩的倾向摆在我们面前；藉着这些，按照圣父的旨意，祂愿意使我们知道美善与有用。请思考这些。\OriginalTerm{美善}{\GreekText{τὸ} \GreekText{καλόν}}属于颂扬式的言说方式，而有用属于劝服式的。因为劝勉和劝戒是劝服式的一种形式，而赞扬和指责是颂扬式的一种形式。\par

因为劝服式句子的一种形式成为劝勉，另一种成为劝戒。同样，颂扬式的一种形式成为指责，另一种成为赞扬。在这些操练中，导师，即公义的那位，祂以我们的益处为目标，是主要关注的。但指责和劝戒的言说方式已经向我们显示；现在我们必须处理劝服和赞扬，并且如同在衡量上平衡公义的天平。对有益的劝勉，导师借由\Person{所罗门}采用，大意如下：\BibleQuote{我劝告你们，人子啊；我向世人发声。听我，因为我要谈论卓越的事；}\BibleRef{箴 8:4}等等。祂劝告那有救恩性的事：因为劝告的目的在于选择或拒绝某一途径；正如祂借由\Person{达味}所说的：\BibleQuote{不随从恶人的计谋，不站罪人的道路，不坐亵慢人的座位，而喜爱主的法律的，这人是有福的。}劝告有三个部分：以过去的事为例的；如希伯来人敬拜金牛犊时所受的苦，以及他们行淫时所受的苦，等等。第二部分，其意义从当前的事而了解，就如藉由感官所领受的；正如对那些问主的人所说的：\BibleQuote{祂是基督吗？还是我们等候别人？你们去告诉\Person{若翰}：瞎子看见，聋子听见，癞病人洁净，死人复活；凡不因我而跌倒的，是有福的。}正如\Person{达味}预言时所说的：\BibleQuote{我们听见了，也看见了。}第三部分劝告包括未来的事，藉此我们被命令防备将要发生的事；正如也说过：\BibleQuote{那些陷于罪恶的，将被丢在外面的黑暗里，那里有哀哭和切齿。}等等。所以从这些事中清楚看出，主遍历所有医疗方法的途径，呼召人类走向救恩。\par

藉着鼓励，祂缓和罪恶，减少情欲，同时激发对救恩的希望。因为祂藉\Person{厄则克耳}说：\Quote{你们若全心归回，并说：父啊，我要垂听你们，如同圣洁的子民。} 祂又说：\BibleQuote{凡劳苦和负重担的，你们都到我这里来，我要使你们得安息。}\BibleRef{玛 11:28} 而接下来所加的，是主亲自说的。祂也藉\Person{撒罗满}非常清楚地呼唤人行善，祂说：\BibleQuote{寻得智慧的人是有福的，获得悟性的人是有福的。}\BibleRef{箴 3:13}\Quote{因为善被寻求它的人寻得，且惯于被寻得它的人看见。} 祂也藉\Person{耶肋米亚}陈明智慧，当他说：\BibleQuote{我们有福了，以色列！因为那中悦天主的事，为我们所知；}\BibleRef{巴 4:4}——并且这藉着圣言而被知晓，藉着圣言我们蒙福且有智慧。因为同一先知提到了智慧和知识，他说：\BibleQuote{以色列，请听生活的诫命，侧耳以领悟智慧。}\BibleRef{巴 3:9} 祂也藉\Person{梅瑟}，因祂对人的爱，向那些奔向救恩的人应许一份恩赐。因为祂说：\BibleQuote{我要领你们进入那佳美之地，就是主向你们祖先所誓许的。}\BibleRef{申 31:20} 并且，\Person{依撒意亚}说：\BibleQuote{我还要领你们上我的圣山，使你们喜乐。}\BibleRef{依 56:7} 还有另一种教导形式是祝福。祂藉\Person{达味}说：\BibleQuote{有福的人，是没有犯过罪的；他要像一棵栽在溪水旁的树，按时结果，叶子也不枯干}（祂藉此暗示复活）；\BibleQuote{凡他所行的，尽都顺利。} 祂愿意我们如此，好使我们蒙福。再者，为了显示公义天秤的另一边，祂说：\BibleQuote{恶人却不是这样——绝不是这样；他们乃像糠秕，被风吹散。} 导师藉着显示罪人的惩罚，以及他们轻易被分散、被风吹走，用惩罚来劝阻犯罪；又藉着展示应得的刑罚，以最巧妙的方式显明祂恩赐的良善，好使我们拥有并享受其祝福。祂也邀请我们追求知识，当祂藉\Person{耶肋米亚}说：\BibleQuote{你若行在天主的道上，你必永远安居于平安中；}\BibleRef{巴 3:13} 因为，祂在那里展示知识的赏报，呼唤智慧人热爱知识。并且，祂宽恕那犯错的人，说：\BibleQuote{转回，转回，如同摘葡萄者转向自己的篮子。}\BibleRef{耶 6:9} 你看见公义的美善了吗？它劝人悔改。而且，祂又藉\Person{耶肋米亚}光照那些犯错的人，使他们认识真理。\BibleQuote{上主这样说：你们应站在路上，细察，询问上主永恒的道路，哪一条是善路，然后行在其上，你们必为你们的灵魂找到洁净。}\BibleRef{耶 6:16} 为了促进我们的救恩，祂引领我们悔改。因此祂说：\BibleQuote{你们若是悔改，上主必洁净你们的心，和你们后裔的心。}\BibleRef{申 30:6} 我们本可以引证那些哲学家作为这个问题的支持者，他们说只有完美的人才值得赞美，而恶人应受责备。但因为有人诽谤真福，说它本身不承担任何劳苦，也不使人劳苦，因此不明白它对人的爱；为了他们，也为了那些不将正义与良善相联的人，加上以下这些话。因为可以合理推论，斥责和指责适合人，既然他们说所有人都是恶的；但唯独天主是智慧的，智慧来自祂，唯独祂是完美的，因此唯独祂配得赞美。但我并不采用这种说法。我要说的是，赞美或责备，或任何类似赞美或责备的东西，对人来说是最必要的药物。有些人不堪医治，如同铁器，需要用火、锤子和铁砧来锻造，也就是用威胁、责备和惩罚；而另一些人，紧贴信德本身，如同自学，并凭自己的自由意志行事，则因赞美而成长：——\par

\Quote{因为被称赞的德行\newline{}像树木一样生长。}\par

据我看，\Person{毕达哥拉斯}理解这一点，便给出训令：——\par

\Quote{当你做了卑劣之事，就斥责\Emph{你自己}；\newline{}但当你做了善事，就欢喜。}\par

斥责也称为\OriginalTerm{劝诫}{\GreekText{νουθέτησις}}；而\OriginalTerm{劝诫}{\GreekText{νουθέτησις}}的词源是\OriginalTerm{心智植入}{\GreekText{νοῦ} \GreekText{ἐνθεματισμός}}，即把理解植入人内；因此，指责就是使人醒悟。\par

然而，尚有无数诫命，其目的在达成善而避免恶。\Quote{恶人必不得平安，主说。}因此，祂借撒罗满命令子女警醒：\BibleQuote{我儿，不要让罪人欺骗你，不要随从他们的道路；如果他们引诱你，说：同我们来，与我们一同分享无辜之血，让我们不义地将义人埋在地里；趁他活着，把他扔进阴间，你不要去。}\BibleRef{箴 1:10}这同样是对主受难的预言。借着厄则克耳，生命赋予诫命：\BibleQuote{犯罪的人，必要死亡；但行义的人必是正义的。他不吃山上的祭物，不举目仰望以色列家的偶像，不玷污邻人的妻子，不接近经期的妇人，不欺压人，归还债务人的抵押，不抢劫，将自己的食物给饥饿的人，给赤身者衣服穿。他不放高利贷，不取利息；他收回手不作恶，在人与人之间施行公正的审判。他遵行我的法律，遵守我的典章，遵行它们。这是义人，他必生存——吾主上主的断语。}\BibleRef{则 18:4}这些话描述了基督徒的行为，是对真福生活的显著劝勉，而这真福生活就是善生的赏报——永生。\par

% structure: p0120 source=h2 target=subsection reason=relative-heading-rank; base=h2

\subsection{第十一章　圣言借法律和先知施教}

祂的爱与教导的方式，我们已尽己所能予以指明。因此，祂亲自极优美地描述自己，将祂比作一粒芥子，\BibleRef{玛 13:31}指出所播种的圣言的精神性、其本性的多产性，以及圣言能力的宏伟与显赫；此外，还暗示惩罚的辛辣与净化之德因其锐利而有益。借着这粒小种子（如其比喻所称），祂丰盛地赐予全人类救恩。蜂蜜虽极甜，却能生胆汁，正如善易生轻慢，而轻慢是犯罪之因。但芥子减少胆汁，即愤怒，并制止炎症，即骄傲。由此圣言产生灵魂的真健康及其永恒的\OriginalTerm{福乐和谐}{\GreekText{εὐκρασία}}。\par

因此，祂古时借梅瑟教导，然后借先知教导。梅瑟也是一位先知。因为法律是训练顽梗孩童的。\BibleQuote{他们吃饱了，}经上说，\BibleQuote{就起来玩耍。}\BibleRef{出 32:6}这种无知的饱食被称为\OriginalTerm{饲料}{\GreekText{χόρτασμα}}，而非\OriginalTerm{食物}{\GreekText{βρῶμα}}。当他们无知地吃饱后，又无知地玩耍；为此，法律赐给了他们，并随之而来的是恐怖，以阻止过犯并促进正当行为，确保注意，从而赢得对真导师（即同一圣言）的服从，并使其符合法律的紧迫要求。因为保禄说，法律被赐为\BibleQuote{启蒙师，引领我们归于基督。}\BibleRef{迦 3:24}由此可见，那唯一、真实、善、义、具有圣父肖像和模样者，即祂的圣子耶稣、天主的圣言，乃是我们的导师；天主将我们托付给祂，如同慈父将子女托付给可敬的导师，清楚地命令我们说：\BibleQuote{这是我的爱子，你们要听从祂。}\BibleRef{玛 17:5}这位神圣的导师是信实的，祂以三种最美妙的装饰为饰：知识、仁爱与言说的权威——知识，因为祂是父的智慧：\Quote{一切智慧皆来自上主，并与祂永远同在；}——言说的权威，因为祂是天主与创造者：\BibleQuote{万物是借着祂受造的，凡受造的，没有一样不是借着祂受造的；}\BibleRef{若 1:3}——仁爱，因为唯有祂将自己献为我们的祭品：\BibleQuote{善牧为羊舍掉自己的性命。}\BibleRef{若 10:11}祂也确实这样做了。仁爱无非是愿意为邻人的益处而行善。\par

% structure: p0123 source=h2 target=subsection reason=relative-heading-rank; base=h2

\subsection{第十二章　导师以父亲般情感的严厉与慈爱为特征}

既然完成了那些事，接下来理应由我们的导师耶稣为我们绘出真生命的典范，并在基督内训练人性。\par

祂所命令的生活的形式与特质，并非十分可怕；由于祂的仁慈，也并非全然容易。祂命令祂的诫命，同时赋予它们可实践的特质。\par

我之见是，祂亲自用尘土塑造了人，又用水重生了人；用祂的圣神使之成长；用祂的圣言训练人，使之获得义子的名分和救恩，以神圣的诫命引导人；为的是，使祂的来临将土生之人转变为圣洁和属天的存在，从而圆满地实现那句神圣的话：\BibleQuote{让我们按我们的肖像，照我们的模样造人。}\BibleRef{创 1:26}事实上，基督成了天主所言说的完美实现；而其余的人类则被视为仅是按祂的肖像受造的。\par

但愿我们——良善圣父的儿女，圣善教育者的乳儿——遵行圣父的旨意，听从圣言，并印上我们救主真正拯救的生命之印记；默想我们得以神化所依照的属天生活方式，让我们以永不失色的不朽喜乐之华——那馨香的膏油——傅抹自己，在主的行为和言谈中，拥有不朽的明确典范，并追随天主的足迹，唯有祂能思虑并看顾人的生命如何能变得更健康的方式与方法。此外，祂为一种自足的生活方式、为简朴、为束紧腰身、以及为我们旅程自由无阻的准备，都做了预备；为获得永恒的福乐，教导我们每个人作自己的仓库。因为祂说：\BibleQuote{不要为明天忧虑}，\BibleRef{玛 6:34}意思是那已献身于基督的人应当自足，作自己的仆人，并且过着每一天为自己预备的生活。因为我们不是在战争中，而是在和平中受训。战争需要大量准备，奢侈渴求丰盛；但和平与爱，这对简朴而安静的姊妹，不需要武器或过多的准备。圣言是她们的滋养。\par

我们在训导和纪律上的监督是圣言的职分，我们从祂学习节俭、谦卑，以及一切与爱真理、爱人和爱美德相关的事。所以，一句话，藉着参与道德的美善而被同化于天主，我们不可退回到疏忽和怠惰中。但要劳苦，不要疲倦。你将成为你所不期望、所不能猜测的。正如哲学家有一种训练方式，演说家和运动员各有不同；同样，有一种高尚的性情，适合那专注于道德之美的选择，这来自基督的训练。而对于那些按照这一影响受训练的人，他们行走的步伐、坐席、饮食、睡眠、就寝、养生之道以及其余生活方式，都获得了一种超卓的尊严。因为圣言所实行的这种训练并不过度紧张，而是具有适当的张力。因此，圣言也被称为救主，因为祂为人类发现了那些产生感官活力与救恩的理性药物；并致力于等候有利时机，谴责邪恶，揭露邪恶情感的根源，斩断非理性情欲的根，指出我们应当戒除什么，并为患病者提供一切救恩的解毒剂。因为天主最伟大、最尊贵的工作就是人类的救恩。病人会因医生不给予他们恢复健康的建议而烦恼。但我们怎能不向那位神圣的教育者致以最高的感恩呢？祂并不缄默，不省略那些指向毁灭的警告，反而揭示它们，并切断迈向它们的冲动；而且祂教导那些导向真正生活方式的训诲？因此，我们必须承认对祂的最深义务。因为我们说，理性受造物——我指的是人——除了默观天主之外，还有什么是当行的呢？我也说，默观人性、照真理的引导而活、以及钦佩教育者及其训令，因为它们彼此合适和谐，是必要的。我们也应当依照这肖像，使我们自己符合教育者，并使言语与行为一致，过一个真实的生活。\par

% structure: p0129 source=h2 target=subsection reason=relative-heading-rank; base=h2

\subsection{第十三章 美德是理性的，罪是非理性的}

一切违背正确理性的事都是罪。因此，哲学家们认为应当这样定义最普遍的情欲：贪欲，是违背理性的欲望；恐惧，是违背理性的软弱；快乐，是违背理性的精神亢奋。那么，如果对理性的不顺从是罪的生成原因，我们怎能避免得出结论：对理性——即圣言——的顺从，我们称之为信德，必然是义务的有效原因？因为美德本身是灵魂的一种状态，由理性使其在整个生命中和谐一致。不仅如此，总而言之，哲学本身被宣称为正确理性的培养；因此，凡因理性错误而做的事便是过犯，并正确地被称之为\OriginalTerm{罪}{\GreekText{ἁμάρτημα}}。既然第一个人犯了罪，违背了天主，经上就说：\BibleQuote{人变成了像野兽一样。}由于他正确地被视为非理性，就被比作野兽。因此智慧说：\BibleQuote{那用于交配的马；好色者和奸淫者已变得像非理性的野兽。}\BibleRef{德 33:6}因此又加上：\BibleQuote{凡骑在它上面的，它嘶叫。}意思是，那人不再说话；因为那违背理性的人不再是理性的，而是一个非理性的动物，被情欲所制服，如同马被骑手所骑乘。\par

但是，那合乎理性而正确完成的事，斯多葛派称为\OriginalTerm{本分}{\GreekText{προσῆκον}}和\OriginalTerm{合宜}{\GreekText{καθῆκον}}，即本分和合宜。合宜即是本分。顺从基于命令。这些命令，因与劝告相同——以真理为指向——培育人达至渴求的终极目标，这目标被认作是\Emph{终向}（\OriginalTerm{终向}{\GreekText{τέλος}}）。而虔敬的终向是在天主内获得永息。永生的开始即是我们的终向。虔敬的正确运作藉着行为成全本分；因此，根据正当的推理，本分在于行动，而非言谈。基督徒的品行乃是理性灵魂依照正确的判断与对真理的渴求而运作，这渴求通过身体——灵魂的伴侣与盟友——达成其预定终向。德行是在生活中与天主和基督契合的意志，正确地导向永生。因为基督徒的生活（我们如今在其中受培育）是合乎理性行动的系统——即由圣言所教导的那些事——一种我们称为信德的不竭动力。这系统就是主的诫命，它们是神圣的规条和属神的劝告，为我们而写，适用于我们和我们的近人。此外，它们又回到我们身上，如同球被反弹回到投掷者那里。因此，本分对神圣的纪律是必要的，因为它们由天主治定，为我们的救恩而预备。既然在必需的事物中，有些只关乎此世的生命，另一些则关乎来世的幸福生命，使我们展翅高飞；那么，以类似的方式，本分中有些是为生命而制定，另一些是为幸福生命而制定。关于自然生命的诫命是向大众颁布的；而那些适合于善生、并从中涌出永生的诫命，我们则要像在草图中那样，当我们从圣经中读到它们时加以思考。\par

\SourceRecord{Translated by William Wilson.；Ante-Nicene Fathers；Vol. 2.；Edited by Alexander Roberts, James Donaldson, and A. Cleveland Coxe.；Buffalo, NY: Christian Literature Publishing Co.,；1885.；<http://www.newadvent.org/fathers/02091.htm>.}

% structure: p0001 source=h1 target=section reason=page-title

\section{\WorkTitle{导师（第二卷）}}

% structure: p0002 source=h2 target=subsection reason=relative-heading-rank; base=h2

\subsection{第一章  论饮食}

因此，我们坚守目标，选取那些关乎生活训练之益处的圣经章节，现在必须简明地描述那被称为基督徒的人，在他整个一生中应当如何。我们当从自身开始，论及我们应如何自律。因此，本着对此书整体的合宜考量，我们须说明每个人应如何对待自己的身体，或更确切地说，如何调理身体本身。因为当一个人被圣言从外在事物及对身体的关注中带离，转向心灵时，他便能清晰地看清人性中合乎自然之事，从而明白自己不应热切关注外在事物，而应关注那属于人且为人所特有的——净化灵魂之眼，并圣化自己的肉身。因为那已摆脱那些仍使他成为尘土之物的人，除了他自己，还有什么更合用的工具，能让他行走在通往认识天主的道路上呢？\par

诚然，有些人活着是为了吃，如同无理性的受造物，\Quote{他们的生命就是肚腹，别无他物。}但导师命令我们为了活而吃。因为食物不是我们的事务，快乐也不是我们的目的；两者都是为了我们在此世的生命，而圣言正将这生命训练至不朽。因此，在食物上也需加以分辨。它应当简单、真正朴素，恰好适合单纯无伪的孩童——为生命服务，而非为奢华。它所导向的生命包含两件事：健康与力量；对此，清淡的饮食最为适宜，既利于消化，又使身体轻健，从而带来成长、健康和正当的力量，而非那种依靠强迫进食所获得的运动员式的错误、危险且可怜的力量。\par

因此，我们必须摒弃那些引发各种祸害的食物种类，诸如败坏的身体状况和胃部失调，因那不妙的烹饪术和制作糕点的无用技艺败坏了味觉。因为人们竟敢将他们沉溺于奢华的行为称为食物，而这奢华正悄然滑入有害的快乐之中。\Person{提洛的安提法内斯}，那位\Place{提洛}医生曾说，食物的多样性是疾病唯一的原因；有些人厌恶真理，因各种荒谬观念而摒弃饮食的节制，自寻许多麻烦去弄来海外的珍馐。\par

就我而言，我为这种疾病感到悲哀，而他们却不羞于赞美自己的美味佳肴，费尽心思去获取西西里海峡的七鳃鳗、迈安德河的鳗鱼、米洛斯的山羊、斯基亚索斯的鲻鱼、佩洛鲁斯的贻贝、阿比多斯的牡蛎，还不忘利帕拉的小鲱鱼和曼提尼亚的芜菁；此外，还有阿斯克拉地区生长的甜菜：他们寻找米提利尼的鸟蛤、阿提卡的比目鱼、达夫尼斯的画眉鸟和红褐色的无花果干——正是为了这些，那位不幸的波斯人曾率领五十万大军侵入希腊。除此之外，他们还购买法西斯的鸟、埃及的沙锥鸟和米底亚的孔雀。贪食者用调味料改变这些食物，觊觎着酱汁。\Quote{无论大地、海洋的深处，还是广阔的天空所出产的一切，}他们都为满足口腹之欲而搜罗。贪食者在贪婪和焦虑中，仿佛用拖网彻底扫遍世界，以满足他们奢侈的口味。这些贪食者被嘶嘶作响的煎锅声包围，一生消磨在杵和臼之间，像火一样紧贴着物质。不仅如此，他们还使简单的食物——即面包——失去精华，滤去谷物的营养部分，以致食物的必要部分成为奢侈的耻辱。人类对美食的追求没有止境。因为这已驱使他们追求蜜饯、蜂蜜蛋糕和糖渍果仁，发明了无数甜点，追逐各种菜肴。在我看来，这样的人除了下巴，什么都不是。圣经说：\BibleQuote{不要贪求富翁的美味；}\BibleRef{箴 23:3}因为它们属于虚假而卑贱的生活。他们分享奢侈的菜肴，这些菜肴片刻之后便化为粪土。但我们这些寻求天上之粮的人，必须管束那低于天地的肚腹，更不用说那些与之相宜的事物，宗徒说，\BibleQuote{天主必要毁坏它们，}\BibleRef{格前 6:13}他公正地咒骂贪食的欲望。因为\BibleQuote{食物是为肚腹，}\BibleRef{格前 6:13}因为真正属肉体和导致毁灭的生活依赖于此；因此，有些人纵口无遮拦，竟敢用\OriginalTerm{爱宴}{\GreekText{ἀγάπη}}这个名称来称呼那些可怜兮兮的、充满香味和酱汁的晚宴。他们侮辱圣言的良善和救赎工作，即神圣的\OriginalTerm{爱宴}{\GreekText{ἀγάπη}}，用锅碗和倒酱汁的方式，以及用饮料、美味和烟雾亵渎那名称，他们在观念上受到了欺骗，以为天主的应许可以用晚宴来购买。为了欢愉而聚集的场合，以及我们自己称为的宴席，我们恰当地称之为晚餐、正餐和宴会，效法主的榜样。但主并没有称这样的宴席为\OriginalTerm{爱宴}{\GreekText{ἀγάπη}}。他因此在某处说：\BibleQuote{你被人请去赴婚宴，不要坐在首席上；但当你被请时，要坐在末位上；}\BibleRef{路 14:8}；在别处说：\Quote{你设筵或吃晚饭时；}又说：\BibleQuote{但你设宴时，要请贫穷的，}\BibleRef{路 14:12}——尤其是为了他们的缘故，理当设宴。此外，\BibleQuote{有一个人设了盛大的筵席，请了许多人。}\BibleRef{路 14:16}但我看出这貌似高雅的筵席之名从何而来：\Quote{从贪食者喉咙和疯狂的筵席之爱而来}——正如那位喜剧诗人所说。因为实际上，\Quote{对许多人来说，许多事情都是为了这顿筵席。}因为他们尚未学到，天主为祂的受造物——我指的是人——提供了食物和饮料，是为滋养，而非为享乐；因为身体从奢侈的菜肴中得不到任何益处。相反，饮食最简朴的人恰恰是最强壮、最健康、最高贵的；如同仆从比主人更健康、更强壮，农夫比地主更是如此；不仅如此，他们还更智慧，如同哲学家比富人更智慧。因为他们没有将心思埋葬在食物之下，也没有用享乐欺骗它。但爱（\OriginalTerm{爱宴}{\GreekText{ἀγάπη}}）实为天上的食物，是理性的盛宴。\BibleQuote{爱凡事包容，凡事忍耐，凡事盼望；爱永不止息。}\BibleRef{格前 13:7}\BibleQuote{凡在天主的国里吃饼的，是有福的。}\BibleRef{路 14:15}然而，最糟糕的情况是，那永不止息的爱德竟从天上坠落尘世，坠入酱汁之中。你以为我想到的是一顿将要废去的晚餐吗？经上记载：\BibleQuote{我若将我所有的财产全部施舍，却没有爱德，我就算不得什么。}\BibleRef{格前 13:3}唯有在这爱上，才依存着法律和圣言；如果\BibleQuote{你爱主你的天主和你的邻人，}这就是天上的节庆。但地上的被称为晚餐，正如圣经所示。因为晚餐是为爱而设，但晚餐本身并非爱（\OriginalTerm{爱宴}{\GreekText{ἀγάπη}}）；它只是彼此和相互善意的证明。宗徒说：\BibleQuote{所以，不要让你们的善被人毁谤；因为天主的国不在于吃喝，而在于义德、平安和在圣神内的喜乐。}\BibleRef{罗 14:16}吃这膳食——最好的膳食——的人，将拥有天主的国，在此定睛于爱的神圣集会，即天上的教会。爱，本是纯洁而相称于天主的，其工作是相通。正如智慧所说：\BibleQuote{而训诫的关怀就是爱，而爱就是遵守法律。}\BibleRef{智 6:17}这些喜乐从公共的营养中获得了爱的灵感，这营养使人习惯于永久的佳肴。因此，爱宴不是晚餐。但让宴席建立在爱之上。因为经上说：\BibleQuote{主啊，愿你所爱的儿女们知道，养育人的不是果实的产物，而是你的话保存了那些信靠你的人。}\BibleRef{智 16:26}\BibleQuote{义人不是单靠饼生活。}\BibleRef{申 8:3}但让我们的饮食清淡易消化，适于保持清醒，不混杂多种多样。这并非超出训导的范围。因为爱是相通的良好滋养者；它以充足为丰盛的供应，这充足以适度的量掌管饮食，以健康的方式对待身体，并从其资源中分给身边的人。但超出充足限度的饮食损害人，败坏其精神，使其身体易患疾病。此外，那些精致的口味，为丰盛的菜肴而烦恼，驱使人走向恶名昭彰的行为——讲究、贪吃、贪婪、饕餮、不知足。适合如此放纵之人的称谓是苍蝇、黄鼠狼、谄媚者、角斗士，以及各种畸形的寄生者——一类放弃理性，另一类放弃友谊，还有一类放弃生命，只为满足肚腹；他们用肚子爬行，是人形的野兽，按他们父亲——那贪婪的野兽——的形象。人们最初称那些被遗弃的人为\OriginalTerm{浪费的人}{\GreekText{ἀσώτους}}，这在我看来似乎是指明了他们的结局，将其理解为那些\OriginalTerm{无救的人}{\GreekText{ἀσώστους}}，省略了字母\Emph{\GreekText{ς}}。因为那些沉迷于锅碗和精心准备的精致调味品的人，难道不显然是卑贱的、土生土长的、过着短暂生活的人，仿佛他们不会（在将来）生活吗？这些人，圣神通过依撒意亚斥责为可悲的，暗暗剥夺了他们爱的名称（\OriginalTerm{爱宴}{\GreekText{ἀγάπη}}），因为他们的宴饮不合乎圣言。\BibleQuote{他们却欢乐，宰杀牛犊，宰杀羊，说：我们吃喝吧，因为明天我们就要死了。}而祂将这种奢侈视为罪，由祂接着所说的显明：\BibleQuote{你们的罪必不得赦免，直到你们死去，}\BibleRef{依 22:13}——这并非意指那剥夺感觉的死亡是罪的赦免，而是指那作为罪之报偿的救恩的死亡。\BibleQuote{别喜欢可憎的美味，}智慧说。\BibleRef{德 18:32}此时，我们也必须提及所谓祭偶像之物，以显示我们如何被命令禁绝它们。在我看来，那些被污染的、可憎的东西，它们的血……\par

\Quote{灵魂来自无生命尸体的幽冥。}\par

\BibleQuote{我不愿你们与魔鬼共融} \BibleRef{格前 10:20} 宗徒说；因为得救者与丧亡者的食物是不同的。因此，我们不可因恐惧而戒绝这些肉食（因为它们本身毫无权能），而应出于我们良心的圣善和对献祭给魔鬼的憎恶而厌恶它们；此外，还因那许多人的软弱，他们看待事物易于跌倒，\BibleQuote{其良心软弱就被玷污了：因为食物并不能使我们得到天主的嘉许。} \BibleRef{格前 8:7} \BibleQuote{不是入口的使人污秽，而是出口的使人污秽。} \BibleRef{玛 15:11} 食物的自然使用原是无关紧要的。\BibleQuote{因为我们若吃，并无更好，} 据说，\BibleQuote{若不吃，也无更坏。} \BibleRef{格前 8:8} 但那些被堪当分享神圣和属神食物的人，竟去参与魔鬼的宴席，这是不合情理的。\Quote{我们难道没有权柄吃喝，} 宗徒说，\Quote{并携带妻子吗？} 但藉着克制快乐，我们阻止了私欲。看，所以你的这权柄切勿\Quote{成了软弱者的绊脚石。}\par

因为若我们效法福音中富翁的儿子\BibleRef{路 15:11}，如浪子般滥用天父的恩赐，这并不合宜；但我们应使用它们，不妄加留恋，如同我们掌管自身一样。因为我们受命要统辖并管理肉食，而非成为它们的奴隶。因此，这是一个令人钦佩的事：举目仰望真实者，依赖上天的神圣食粮，以对真实之存有的无尽默观来满足自己，从而品尝唯一可靠而纯净的喜乐。因为这就是\OriginalTerm{爱宴}{agape}，即来自基督的食粮表明我们应当领受的。但那些属乎尘土的人，像牛一样把自己养肥以喂养给死亡；向下看地，弯腰伏桌；过着饕餮的生活；将今生一切美好埋葬在即将结束的生命中；只贪求暴食，以致厨师比农夫更受尊敬——这是完全无理性、无用、且不人道的。因为我们并未废除社交，而是对风俗的陷阱心存警惕，视其为灾祸。因此，应避免美味，只取用少量而必须的食物。\BibleQuote{若有不信者请你们赴宴，你们决定要去}（因为不与放荡者混杂是好事），宗徒吩咐我们\BibleQuote{只管吃摆上的东西，不要为良心的缘故追问什么。} \BibleRef{格前 10:27} 同样，他也吩咐购买\BibleQuote{市场上所卖的，} 不要好奇追问。\BibleRef{格前 10:25}\par

我们不应完全戒绝各种食物，而只是不可为其所羁。我们应如基督徒所宜，怀着对邀请者的尊敬，以无害且有节制的方式参与社交聚会，享用摆在我们面前的食物；视桌上之丰盛为无关紧要，藐视那些佳肴，因其不久必将朽坏。\BibleQuote{吃的人不可轻视不吃的人；不吃的人也不可判断吃的人。}\BibleRef{罗 14:3} 稍后，他便解释这条命令的理由，说：\BibleQuote{吃的人是为主吃，因为他感谢天主；不吃的人也是为主不吃，他同样感谢天主。}\BibleRef{罗 14:6} 因此，真正的食物是感恩。感恩的人不沉溺于享乐。倘若我们想劝说同席的客人归向德行，就更应为此戒绝那些美味佳肴；如此，我们便显示自身是德行的光辉榜样，如同我们在基督内所拥有的那样。\BibleQuote{所以，若这些食物使弟兄跌倒，我就永远不吃肉，免得使我的弟兄跌倒。}\BibleRef{格前 8:13} 我以少许的节制赢得这人。\BibleQuote{我们难道没有权利吃喝吗？}\BibleRef{格前 9:14} 并且\Quote{我们知道}——他说的是真话——\Quote{偶像在世界中算不得什么；我们只有一位真天主，万物都出于祂，并有一位主耶稣。}他又说：\Quote{然而，因着你的知识，你那软弱的弟兄便丧亡了，基督原是为他死了的；你们这样得罪软弱的弟兄，并伤他们软弱的良心，就是得罪基督。}因此，宗徒出于对我们的关怀，在宴饮的场合加以区分，说道：\BibleQuote{若有人称为弟兄，却是淫乱者、或奸淫者、或拜偶像者，不可与他一同吃饭。}\BibleRef{格前 5:11} 无论在言谈或饮食上，我们都不可与他们交往，怀疑从此而来的玷污，如同在魔鬼的筵席上。\BibleQuote{因此，不吃肉、不喝酒，}\BibleRef{罗 14:21} 这是他和毕达哥拉斯学派都认可的。因为这更是野兽的特性；而且它们升起的浓密烟雾，会使灵魂昏暗。若有人享用它们，他并不犯罪。只是要节制地享用，不可依赖它们，也不可贪求美食。因为会有声音低语对他说：\BibleQuote{不可为食物毁坏天主的事业。}\BibleRef{罗 14:20} 因为在圣言中已有丰盛的筵席之后，反而对俗宴所呈现之物惊异迷惑，是愚昧心智的标志；而让自己的眼睛成为美味的奴隶，以至于贪欲仿佛被仆人牵着走，更是愚不可及。众人从卧榻上起身，几乎将脸埋入盘中，从卧榻上探出身来，如同从巢中伸出，正如俗语所说：\Quote{好吸入飘散的蒸气！}这是何等愚蠢！用手涂抹调料，不停地伸向酱汁，无节制且无耻地塞满自己，不像品尝的人，而是贪婪地夺取，这是何等无理性！因为你可以见到这样的人，与其说是人，不如说更像猪狗般贪食，急于吃饱，以致双颚同时塞满，面部青筋暴起，而且汗流浃背，因贪得无厌而绷紧，因过量而气喘；食物被带着不合群的热切推入胃中，仿佛将食物储存起来作为旅行的补给，而非为了消化。过度——在一切事上都是恶——在饮食方面尤其应受谴责。贪食（希腊文称为\OriginalTerm{贪食}{\GreekText{ὀψοφαγία}}）无非是过量使用\OriginalTerm{调味品}{\GreekText{ὄψον}}；\OriginalTerm{咽喉癫狂}{\GreekText{λαιμαργία}}是对咽喉的癫狂；\OriginalTerm{腹癫狂}{\GreekText{γαστριμαργία}}是对食物的过度——如名称所示，对肚腹的癫狂；因为\OriginalTerm{狂人}{\GreekText{μάργος}}就是疯人。宗徒责备那些在宴饮中行为失当的人，说：\BibleQuote{因为在你们聚会时，各人先吃自己的晚餐；有的饥饿，有的醉酒。难道你们没有家可以吃喝吗？或是你们轻视天主的教会，并使那些没有的人羞愧呢？}\BibleRef{格前 11:21} 而在那些有家的人中，那些无耻地吃且贪得无厌的人，羞辱了自己。两者都行为不当：前者使没有的人痛苦，后者在有的人面前暴露自己的贪婪。因此，宗徒必然继续以不悦之声斥责那些丢掉羞耻、毫无节制地滥用宴饮的人，那些贪得无厌、永不满足的人：\BibleQuote{所以，我的弟兄们，当你们聚集吃饭时，要彼此等待。若有人饿了，可以在家里先吃，免得你们聚集受审判。}\BibleRef{格前 11:33}\par

我们必须戒除一切奴性的习气和过度，以端正的方式取用摆在面前的食物；保持手、躺椅和下巴洁净；保持面容的优雅不受干扰，在吞咽时不做任何不雅之举；而是有秩序地间歇伸出手。我们必须谨防在进食时说话：因为当嘴巴塞满时，声音变得难听且含混；舌头被食物压迫，阻碍了其自然功能，发出压抑的声音。同时进食和饮水也不适宜。因为混淆用途相悖的时间，实为无节制的极致。并且\BibleQuote{你们或吃或喝，一切都要为光荣天主而做，} \BibleRef{格前 10:31} 追求真正的节俭，主在祝福饼和熟鱼以宴请门徒时，似乎也暗示了这一点，为我们引入了简单食物的美好榜样。那主命令伯多禄所捕的鱼，便指向易于消化、天主所赐且适中的食物。并且通过那些从水中上钩以得正义的人，祂告诫我们要除去奢侈和贪婪，就像从鱼中取出的钱币；为了除去虚荣；并通过将\OriginalTerm{斯塔特钱币}{stater}交给税吏，且\BibleQuote{凯撒的归凯撒，} 以保存\BibleQuote{天主的归天主。} \BibleRef{玛 22:21} \OriginalTerm{斯塔特钱币}{stater}还可以有其他解释，为我们所熟知，但当前不适合讨论它们。我们为当前目的所做的提及就足够了，因为它并非不适合\OriginalTerm{圣言}{Word}的花朵；我们常常这样做，将最有益的泉源引向问题的紧要之处，以浇灌那些由圣言所栽种的人。\BibleQuote{凡事我都可行，但不都有益处。} \BibleRef{格前 10:23} 因为那些做一切可行之事的人，很快便会堕入非法之事。正如正义不因贪婪而得，节制不因过度而得；同样，基督徒的生活方式也不因放纵而形成；因为真理的餐桌远离淫佚的美味。因为虽然万物主要是为人而造，但并非一切事物都适合使用，也并非任何时候都适合。因为场合、时间、方式和意向，对于一个正确受教者而言，在很大程度上决定了何为有用；这些是适宜的，并且有助于制止饕餮的生活，而财富容易选择这种生活，并非那看得很清楚的财富，而是那使人在贪食方面瞎眼的富足。在必需品方面，没有人是贫穷的，人也绝不会被忽视。因为有一位天主喂养飞鸟和鱼类，总之，喂养一切无理性的受造物；它们一无所缺，虽然\BibleQuote{它们不为食物忧虑。} \BibleRef{格前 10:23} 我们比它们优越，是它们的主人，并因更智慧而与天主更亲近；我们受造，不是为吃喝，而是为献身于认识天主。\BibleQuote{义人吃饱，必使灵魂满足；恶人的肚腹，必要缺乏。} \BibleRef{箴 13:5} 充满无止尽贪食的欲望。奢侈的花费不仅适于享乐，也适于社交。因此我们必须提防那些使我们在不饿时仍想吃东西的食物，它们诱惑食欲。因为在节制的简朴中，难道没有各种健康的食物吗？球茎、橄榄、某些草药、奶、奶酪、水果、各种不加酱汁的熟食；如果想要肉类，可摆上烤的而不是煮的。\BibleQuote{你们这里有什么吃的没有？} \BibleRef{路 24:41} 主在复活后对门徒说；他们因受主教诲而实行节俭，\Quote{便给了祂一片烤鱼。} 路加说，祂在他们面前吃了以后，便对他们说了所说的话。此外，不可忽视的是，那些按照圣言进食的人，并不被排除在蜜脾这样的美味之外。因为在食物中，最适宜的是那些不需要火即可立即食用的，因为它们最方便；其次是我们之前说过的最简单的。但是那些围着刺激食欲的餐桌，滋养自己疾病的人，被一个最贪馋的魔鬼统治，我将毫不羞愧地称它为\OriginalTerm{肚腹之魔}{Belly-demon}，是最坏、最堕落的魔鬼。因此，它完全像那被称为\OriginalTerm{腹语之魔}{Ventriloquist-demon}的。快乐远胜于有魔鬼与我们同住。而快乐在于修德。因此，宗徒玛窦只吃种子、坚果和蔬菜，不吃肉。而若望，将节制发挥到极致，\Quote{吃蝗虫和野蜜。} 伯多禄戒避猪肉；\Quote{但他神魂超拔，} 正如宗徒大事录所记载的，\BibleQuote{他看见天开了，一个器皿降下，好像一块大布，系着四角，缒到地上，里面有各样四足兽、爬虫和天上的飞鸟。有声音对他说：起来，宰了吃！伯多禄却说：主，绝对不可！因为我从来没有吃过凡俗和不洁之物。第二次又有声音对他说：天主所洁净的，你不可称为凡俗！} \BibleRef{宗 10:10} 因此，它们的食用对我们而言是无所谓的。\BibleQuote{不是入口的污秽人，} \BibleRef{玛 15:11} 而是关于不洁的虚妄意见。因为天主在创造人时说：\BibleQuote{凡地上的一切都赐给你们作食物。} \BibleRef{创 9:2} \BibleQuote{有情吃蔬菜，胜于无情吃肥牛。} \BibleRef{箴 15:17} 这很好地提醒了我们上面所说的：蔬菜不是爱，而是我们的进餐应带着爱；在这些事上，中庸是好的。的确，在一切事上都是如此，尤其在为宴会的准备上，因为极端危险，而中庸良好。不缺乏必需品便是中庸。因为符合本性的欲望以满足为限度。犹太人被法律以最系统的方式命令实行节俭。因为导师通过梅瑟，禁止他们使用无数事物，并附加理由——属灵的理由是隐藏的；属肉的理由是明显的，他们确实信靠这些；论到某些动物，因为不分蹄；另一些因为不反刍；还有一些因为在水产中唯独没有鳞；以致总共只剩下少数适合他们食用。在他允许他们接触的动物中，他禁止了那些死去的、祭过偶像的或勒死的；因为接触这些是不合法的。因为享用美味的人不可能戒除不享用它们，所以他指定了相反的生活方式，直到他打破因习惯而产生的放纵倾向。快乐常在人类中产生伤害和痛苦；饱足在灵魂中生出不安、遗忘和愚昧。据说，儿童的身体在长高时，因营养不足而会长得端正。因为那时，为了生长而遍及身体的精神，不会被大量食物阻碍其自由进程而受到抑制。因此，那位追求真理的哲学家柏拉图，在谴责奢侈生活时，扇起了希伯来哲学的火花，说：\Quote{我来到此处，这里所谓快乐的生活，充满了意大利和叙拉古的餐桌，丝毫不能使我满意，[它在于]每日两次吃饱，夜里从未独自安眠，以及伴随这种生活方式的任何其他附属物。因为普天之下，若有人自幼在这样的实践中长大，无论天生禀赋多么卓越，也绝不会成为智者。} 因为柏拉图并非不熟悉达味，达味\Quote{将约柜安放在自己的城中，在会幕中间；} 并命令所有臣民在主前欢乐，\Quote{向以色列全军，无论男女，每人分给一个饼，一块烤肉，一个油饼。}\par

这就是以色列子民足够的食粮。但外邦人的食粮却过于丰盛。凡使用它的人绝不会努力成为有节制的，因为他把心智埋葬在肚腹里，很像一种名为鲇鱼的鱼，\Person{亚里士多德}说，它是所有生物中唯一心脏长在胃里的。这种鱼，喜剧诗人\Person{埃庇卡摩斯}称之为\Quote{大肚腹}。\par

这样的人信靠自己的肚腹，\BibleQuote{他们的天主是肚腹，以羞辱为光荣，只思念地上的事}。宗徒对他们所说的不吉之言，\BibleQuote{他们的结局是丧亡}。\BibleRef{斐 3:19}\par

% structure: p0014 source=h2 target=subsection reason=relative-heading-rank; base=h2

\subsection{第二章　论饮酒}

\BibleQuote{稍微用点酒}，宗徒对饮水多的弟茂德说，\BibleQuote{为了你的胃}，\BibleRef{弟前 5:23}十分恰当地将酒的助益用作滋补的强壮剂，适合于因水湿体液而虚弱多病的身体；并指明\Quote{稍微}，以免这药因用量不当而不被察觉地产生需要另加治疗的情况。\par

因此，口渴时天然的、有节制的、必需的饮料是水。这是恬淡寡欲的简单饮料，曾由主从击打的磐石中流出，供给古希伯来人。\EditorialNote{出 17；户 20}在他们的漂泊中，节制是极为必要的。\par

此后，神圣的葡萄树结出了预言的葡萄串。这对他们是一个标记，当他们从漂泊中被训练到安息时，代表了那大葡萄串——为我们将要压碎的\OriginalTerm{圣言}{Word}。因为葡萄的血——即圣言——渴望与水调和，正如祂的血与救恩调和。\par

主的血是双重的。有祂肉体的血，藉此我们由腐朽中被救赎；也有属灵的，藉此我们被傅油。饮耶稣的血，就是有分于主的不朽；\OriginalTerm{圣神}{Spirit}是圣言的能力本原，正如血是肉体的本原。\par

因此，正如酒与水调和，\OriginalTerm{圣神}{Spirit}也与人调和。一方面，酒与水的混合物滋养信德；另一方面，圣神引领人走向不朽。\par

两者的调和——水与圣言——被称为\OriginalTerm{圣体圣事}{Eucharist}，那是备受赞誉的荣耀\OriginalTerm{恩宠}{grace}；凡藉信德领受它的人，灵魂和身体都得以圣化。因为那神圣的调和——人，圣父的意愿通过\OriginalTerm{圣神}{Spirit}和\OriginalTerm{圣言}{Word}神秘地组合而成。因为，事实上，圣神与灵魂结合，灵魂由祂所感召；而肉体——因圣言成为血肉的缘故——与圣言结合。\par

因此，我钦佩那些采纳了苦修生活、喜爱饮水（节制的良药）并尽可能远避酒的人，视酒如火险而躲避。所以，少年男女应尽可能远离这药。因为将最热的饮料——酒——倒入生命燃烧的季节，犹如火上浇油。由此激起狂野冲动、燃烧情欲和火爆习性；少年人内里燃烧，倾向于放纵恶性癖好；以致身体出现损伤迹象，欲念的肢体过早成熟。乳房和生殖器官因酒而发炎，可耻地膨胀肿胀，已预先显出淫乱的形象；身体迫使灵魂的伤口发炎，无耻的悸动接踵而至，刺激行为端正的人行越轨之事；于是青年的纵欲越过了谦逊的界限。我们必须尽可能扑灭青年的冲动，移开那危险威胁的酒神燃料；并倾注消炎的解毒剂，以压制燃烧的灵魂，抑制肿胀的肢体，平息已在骚动的欲望。至于成年人，那些适合的人可有时进餐，只品尝面包，完全戒绝饮料；以便多余的水分被干燥食物的摄入吸收和吸干。因为不断吐痰、擦汗和急于排泄，是过度之象，是由过量供应给身体的液体无节制使用所致。若口渴来临，应以少量水满足食欲。因为不应供应过量的水，以免食物被淹死，而应磨碎以助消化；这发生在食物凝聚成团时，只排出少量。\par

此外，神圣的研习也不宜醉酒。\Quote{因为纯酒远不能使人明智，更不用说节制，}正如喜剧诗人所说。但傍晚、晚餐前后，当不再从事较严肃的阅读时，可用酒。这时空气也比白天更冷，以致衰弱的天然热量需要通过引入热来滋养。但即使如此，也只能用少量酒；因为不可走向无度的豪饮。年事已高的人可以更畅快地饮酒，以这无害的葡萄药物来温暖岁月衰败所致的老年寒意。因为老年人的情欲通常不会激动到驱使他们堕入醉酒海难的地步。他们藉理性和时间如同锚泊，更能忍受从无节制涌来的情欲风暴。在宴会上，他们也可以被允许享受玩笑。但对于他们，饮酒的限度应达到保持理性稳定、记忆活跃、身体不被酒动摇不晃的程度。在这种状态中的人，在懂行的人口中被称为\OriginalTerm{\GreekText{ἀκροθώρακες}}{acrothorakes}。因此，及早停杯是好的，以免失足。\par

一位名为\Person{阿尔托里乌斯}的人在所著\WorkTitle{论长寿}一书中（据我所忆）认为，饮酒仅宜至食物湿润为止，以求延年益寿。因此，有人将酒用作药物，纯为健康；另有人则为之放松与享乐。初饮酒后，人较前更和善，对酒伴更可亲，对家人更慈祥，对朋友更愉悦。但醉后却变为暴戾。因为酒性温热，适量调合时甘甜，其温暖能化解不洁的排泄物，并将辛辣卑劣的体液与宜人芳香混合。\par

因此话说得好：\Quote{酒初造时，是灵魂和心灵的喜乐，适量饮用。}\BibleRef{德 31:27}最好尽可能多兑水，不可视酒如水，以致沉溺醉乡，也不可因嗜酒而将酒像水一样灌下。因为二者皆是天主的化工；水与酒之调和，共助健康，因为生命由必需与有益之物构成。水为生命所需，当大量使用，亦当与有益之物相调合。\par

饮酒过量，则舌结唇弛，眼球乱转，视线因水汽过多而如浮游，被迫欺骗，以为万物皆在绕己旋转，远物不能一一计数。\Quote{真的，我以为看见两个太阳。}底比斯老人举杯时说。因为视线被酒的热力所扰，常将一物之实体幻想为多。移动眼睛与被视之物并无区别。两者对视线的效果相同，因晃动而不能准确感知对象。双足如在洪水中从人下被冲走，随之而打嗝、呕吐和醉醺醺的胡言乱语；\Quote{每个醉汉，}据悲剧所言，——\par

\Quote{被愤怒征服，失去理智，\newline{}喜欢倾吐许多蠢话；\newline{}常不愿倾听\newline{}自己甘愿说出的恶言。}\par

在悲剧之前，智慧早已呼喊：\Quote{多饮烈酒，易生愤怒与种种错误。}\BibleRef{德 31:29}因此多数人认为，宴饮时当放松，将严肃事务推迟到早晨。但我认为，此时尤应引入理性介入宴席，充当饮酒的\OriginalTerm{导师}{pædagogue}，以免社交不知不觉沦为酗酒。因为正如无人认为睡前有必要闭眼，同样无人可正当地希望理性缺席宴席，或蓄意使之沉睡直至事务开始。然而圣言永不离开属于祂的人，即使我们睡眠时亦然；因为我们也当邀请祂进入睡眠。因为完美的智慧，即对天主与人事物之知识，涵盖所有关于照管人类羊群之事，在生命方面成为艺术；如此，当我们活着时，它始终与我们同在，不断完成其本职工作，其成果便是善生。\par

然而那些可悲之徒，将节制逐出宴席，认为纵饮为最幸福的生活；他们的生活不过是狂欢、放荡、沐浴、过量、便溺、怠惰、酗酒。你可见其中一些人半醉蹒跚，颈挂花环如酒瓮，以友爱为名相互呕吐酒水；另一些人饱受纵欲之苦，肮脏、面色苍白、铁青，在昨日的酗酒之上又添一夜之饮直至黎明。朋友们，最好最好与这幅景象保持最远距离，并塑造自己趋向更善，唯恐我们也成为他人所见之笑柄。\par

恰如其言：\Quote{如炉火在淬炼中试用钢刃，酒亦考验骄傲之人的心。}\BibleRef{德 31:26}纵酒是过量用酒，醉态是此种用酒导致的紊乱；\OriginalTerm{宿醉}{\GreekText{κραιπ}\GreekText{αλη}}是纵酒后之不适与恶心；其名源于\OriginalTerm{摇头}{\GreekText{κ}\GreekText{αρα} \GreekText{π}\GreekText{αλλειν}}。\par

如是生活（若可称为生活，其实耗费在怠惰、对感官放纵的焦躁及纵欲的幻觉中），神圣的智慧嗤之以鼻，并命令其子女：\Quote{不要嗜酒，不要花钱买肉；因为凡酗酒和行淫者必致贫穷，凡懒惰者必衣不蔽体。}\BibleRef{箴 23:20}因为凡不警醒于智慧而沉溺于酒者，便是懒惰者。他又说：\Quote{醉汉，}他说，\Quote{必衣衫褴褛，在众人面前为自己的醉态羞惭。}\BibleRef{箴 23:21}因为罪人的创伤是肉身衣裳的破洞，是情欲所撕裂的孔洞，透过它们，内在灵魂的羞耻——即罪——暴露无遗，因这罪，那衣裳难以保全，它已被四面撕破，在多欲中朽坏，从救恩中撕裂。\par

因此，祂加上这些最劝诫的话。\Quote{谁有祸患？谁有争吵？谁有纷争？谁有可厌的胡言？谁有无用的悔恨？} \BibleRef{箴 23:29} 你看，那爱好酒的人，穿着破烂衣服，轻视圣言本身，并已放弃自己，沉溺于酗酒。你看圣经发出何等威吓的话反对他。并且在其威吓之上又加上：\Quote{谁的双眼发红？岂不是那些长时泡在酒中，寻找宴饮之处的人吗？} 祂在此指出，那爱喝酒的人已对圣言死了，因为提到布满血丝的眼睛——这是出现在尸体上的标记，向他宣告在主内的死亡。因为忘记导向真生命的事物，便倾向于毁灭。因此，导师关心我们的救恩，有理由禁止我们：\Quote{不要醉酒。} 为什么？你会问。因为，祂说：\Quote{你的口那时要说出乖谬的话，你躺卧如同在海中心，又像大浪中船的舵手。} 因此，诗歌也来帮助我们，说道：\par

\Quote{愿有火一般力量的酒来到人前。\newline{}它将搅动他们，如北风或南风搅动利比亚的海浪。}\par

并且进一步说：——\par

\Quote{酒使言语游移，泄露一切秘密。\newline{}欺骗灵魂的酒是饮者之祸根。}\par

如此等等。\par

你看见覆舟的危险。心在过多的酒中沉没。酗酒的过度被比作海的危险，在其中身体一旦像船一样沉没，便下降到污秽的深处，被大酒浪淹没；而舵手，即人的理智，在酗酒的巨浪上颠簸，这巨浪高高涌起；并被埋在波谷中，被暴风雨的黑暗所蒙蔽，漂离真理的港湾，直至撞在海中的礁石上，毁灭了，被自己驱向肉欲的放纵。\par

因此，宗徒有理由吩咐说：\Quote{不要醉酒，其中有许多放纵；} 使用\Quote{放纵}一词（\OriginalTerm{放纵}{\GreekText{ἀσωτία}}）暗示酗酒与救恩的\OriginalTerm{不一致}{\GreekText{τὸ} \GreekText{ἄσωστον}}。因为祂虽在婚宴上把水变成酒，但并没有准许人醉酒。祂赋予了法律意义的水要素以生命，用祂的血充满了那属亚当的执行者，即整个世界；用真理之藤的酒供给虔诚，这酒是旧法律与新圣言的混合，以应验预定的时期。因此，圣经将酒命名为圣血的象征；但责备那些卑劣地喝剩酒渣的人说：\Quote{酒是放纵，酗酒是傲慢。} \BibleRef{箴 20:1} 因此，根据正确的理性，我们可以因冬天的寒冷而饮酒，直到麻木从那些易感的人身上驱散；在其他情况下，则作为肠道的药物。因为，我们既应当用食物充饥，也应用饮料解渴，并极其谨慎地防备过失：\Quote{因为酒引入是危险的。} 这样，我们的灵魂将是纯洁、干燥和光明的；灵魂本身干燥时最为智慧、最为美好。同样，它也适于默观，不会因酒气而潮湿，形成一团云状的雾气。因此，我们不必费心去获取希俄斯酒（若无），或阿里乌西安酒（若不在手边）。因为渴是一种缺乏感，渴求适合满足缺乏的手段，而非奢华的酒。从海外输入酒是为了那些因过度而衰弱的食欲，其中灵魂甚至在醉酒之前就已对其欲望疯狂。因为有芬芳的塔索斯酒，气息怡人的莱斯博斯酒，甜美的克里特酒，甜美的叙拉古酒，孟杜西安酒（一种埃及酒），以及岛屿纳克索斯酒，\Quote{香气浓郁、风味独特}的酒，还有意大利地方的另一种酒。这些是许多名称。对于节制的人，一种酒就够了，那就是独一天主栽培的产物。因为为何自己本地的酒就不能满足人的欲望，除非他们也像愚蠢的波斯王一样输入水？科阿斯佩斯河，一条印度河流，曾是最佳饮用水——科阿斯皮安水——的来源。正如酒一旦被饮用，就使人爱上它，水也一样。圣神通过亚毛斯发声，谴责富人因奢侈而受祸：\BibleRef{亚 6:4} \Quote{那些喝过滤酒，躺卧在象牙床上的人，} 他说；以及类似的谴责。\par

尤其要注意端庄（就如神话所代表的雅典娜，不管她是谁，出于端庄，放弃了吹笛的快乐，因为那景象不雅）：因此我们要饮酒时不扭曲面容，不贪婪地抓住酒杯，也不在饮前以不雅的动作转动眼睛；也不因不节制而一口喝干杯子；也不溅湿下巴，不在一口气吞下所有酒液时溅湿衣服——脸几乎填满碗口，淹没其中。因为酒液猛烈冲入，伴随大量吸气，仿佛倒入陶器时发出的咕噜声，同时喉咙因快速吞咽而发出声响，这是不节制之可耻不雅的景象。此外，急于饮酒对饮用者有害。我的朋友，不要急于作恶。你的饮料不会被夺走。它被赐予你，并等候你。不要急于张开喉咙喝下而胀破。即使你喝得慢些，遵守礼节，有秩序地一小口一小口喝，你的口渴也得到满足。因为不节制贪婪夺取的，不会因放慢速度而被夺去。\par

\Quote{不要饮酒逞强，}他说，\Quote{因为酒已击败了许多人。}\BibleRef{德 31:25}斯基泰人、克尔特人、伊比利亚人和色雷斯人，全都是好战的民族，极其沉迷于醉酒，并认为这是一件光荣而愉快的追求。但我们，和平之民，为合法享乐而欢宴，非为放荡，我们喝下清醒的友谊之杯，使我们的友谊能以真正名副其实的方式表现出来。\par

你认为主为了我们而降生成人时是怎样饮的？像我们一样无耻吗？岂不是带着规矩合宜？岂不是审慎的吗？因为你要确信，祂自己也喝了酒；因为祂也是人。祂祝福了酒，说：\Quote{你们取饮：这是我的血}——即葡萄树的血。祂以比喻方式称圣言为\Quote{为多人流出的，为赦免罪过}——即喜乐的圣流。祂清楚地通过宴席上的教导表明，饮酒者应当节制。因为祂教导时并未受酒的影响。祂又表明所祝福的正是酒，当祂对门徒说：\Quote{我不再喝这葡萄树的果实，直到我与你们在我父的王国里同饮。}主所饮的是酒，祂又亲自论及自己，责备犹太人心硬时告诉我们：\Quote{人子来了，他们却说：看，一个贪吃好酒的人，税吏的朋友。}\BibleRef{玛 11:19}我们应当牢记这一点，以对抗那些被称为\OriginalTerm{苦行派}{Encratites}的人。\par

女人们，诚然声称追求优雅，为了不让嘴唇因张大在宽阔的杯口而撕裂，从而撑大嘴，便用过于窄小的\OriginalTerm{细颈瓶}{alabastra}不体面地饮酒：她们仰头向后，不雅地裸露脖颈，在我看来；吞咽时扩张喉咙，大口灌下酒液，仿佛要向酒伴尽量袒露一切；并且像男人，更确切地说像奴隶一样打出嗝来，纵情于奢靡狂欢。因为没有任何可耻之事适合有理性的人，对于女人更是如此，她甚至连思考自己的本性都应有羞耻心。\par

经上说：\Quote{醉酒的妇人是极大的愤怒，}\BibleRef{德 26:8}仿佛醉酒的女人是天主的愤怒。为何？\Quote{因为她不会隐藏自己的羞耻。}\BibleRef{德 26:8}因为女人一旦把选择放在享乐上，就很快被引向放荡。我们并没有禁止用细颈瓶饮酒；但我们禁止一味追求专用细颈瓶，视为傲慢；劝告女人随遇而安地使用手边的东西，并根除她们心中危险的情欲。至于那些翻涌上来导致打嗝的气流，让它无声地排出吧。\par

然而绝不可容许女人裸露或展示身体的任何部分，以免双方都跌倒——男人因观看而被激动，女人因招引男人的目光。\par

我们总应当时刻如在主面前行事，免得祂对我们说，正如宗徒愤怒地对格林多人所说：\Quote{你们聚集在一起，不是吃主的晚餐。}\par

在我看来，数学家所称的\OriginalTerm{无头星}{Acephalus}排在游星之前，头垂于胸前，似乎是贪食者、纵欲者和酗酒者的写照。因为在这些人的头部，理性官能并不居于正位，反而屈从于肠腹的情欲，被贪欲和怒气所奴役。正如\Person{厄尔珀诺尔}因醉酒折断了脖子，同样，被酒醉眩晕的头脑从高处重重摔下，落入肝脏和心脏，即落入纵欲和怒气中——正如诗人之子们说\Person{赫淮斯托斯}被\Person{宙斯}从天上抛到地上。经上说：\Quote{失眠的烦恼、胆汁和绞痛，总伴随着不知足的人。}\BibleRef{德 31:20}\par

因此\Person{诺亚}的醉酒也被记载下来，让我们在眼前清晰地读到他的过犯时，全力防备醉酒。为此，那些遮盖他醉酒之耻的人被主祝福。圣经于是以极其全面的纲要，用一句话概括了一切：\Quote{对有训诲的人，酒已足矣，他将在床上安息。}\par

% structure: p0047 source=h2 target=subsection reason=relative-heading-rank; base=h2

\subsection{第三章 论贵重器皿}

因此，使用金银制成的杯子以及其他镶有宝石的器皿，是不合适的，它们只是视觉上的欺骗。因为如果你把任何热的液体倒入其中，器皿会变热，触摸起来很痛苦。反之，如果你倒入冷的液体，材料会改变其性质，损害混合液，使得这富贵的饮品有害。那么，抛弃\OriginalTerm{特里克里安杯}{Thericleian cups}和\OriginalTerm{安提戈尼得斯杯}{Antigonides}、\OriginalTerm{坎塔罗斯杯}{Canthari}、高脚杯、\OriginalTerm{勒帕斯泰杯}{Lepastæ}，以及无尽形状的饮器、酒冷却器和酒注壶吧。因为总的来说，金银，无论是公共还是私人，当超过必需时，是一种令人嫉妒的财产，难得获取，难以保管，且不适合使用。还有玻璃雕刻的精巧虚荣，因为工艺更易破碎，教导我们在饮用时恐惧，也应从我们良好的秩序中驱逐。还有银制的卧榻、锅、醋碟、木盘和碗；除此之外，还有银和金制的器皿，一些用于盛放食物，另一些用于其他我羞于提及的用途，使用易裂的雪松木、松木、乌木、象牙制的三脚架、银脚镶象牙的卧榻、镶金并饰以玳瑁的床扉，以及紫色和其他难以制作的颜色的床单，这些都是无品味奢侈的证明，嫉妒和柔弱的狡诈发明——都应放弃，因为它们根本不值得我们费力。\BibleQuote{因为时间是短暂的。}\BibleRef{格前 7:29}宗徒如此说。那么这剩下的就是我们不要成为可笑的人物，如同有人在公共表演中，外表涂油引人注目而内里悲惨。他更清楚地解释，补充说：\BibleQuote{今后有妻子的，要像没有一样；有买卖的，要像没有一样。}\BibleRef{格前 7:29}如果他在婚姻上这样说——关于婚姻，天主说：\BibleQuote{你们要繁殖增多。}——那么你岂不认为无意义的炫耀应按主的权柄被驱逐吗？因此主也说：\BibleQuote{变卖你所有的，施舍给穷人；然后来跟随我。}\BibleRef{玛 19:21}\par

跟随天主，脱去骄傲，脱去消逝的炫耀，拥有属于你的那美好的、唯一不能夺走的——对天主的信德，对那位受苦者的宣认，对人的仁爱，这是最珍贵的财产。就我而言，我赞同\Person{柏拉图}，他明确立下法则：人不应为金银财富而劳苦，也不应拥有无用的器皿，除非是为了某种必要的和适度的目的；这样同一件物品可以用于多种用途，而多样物品的拥有可以消除。因此，神圣的\WorkTitle{圣经}美妙地针对夸口者和自爱者说：\BibleQuote{列国的统治者在哪里？地上野兽的主宰在哪里？他们在天上的飞鸟中戏耍，他们积蓄了金银，人信赖他们，他们的财富无穷无尽，他们制造了金银，充满忧虑？找不到他们的作品。他们已经消逝，下到阴府了。}这就是炫耀的报应。因为我们当中耕种土地的人需要锄头犁头，但没有人会用银子做镐头或用金子做镰刀，而是使用对农业有用的材料，而非昂贵的。什么能阻止那些有能力思考类似事物的人，对家用器皿持同样的态度，以使用而非花费为尺度？告诉我，餐刀非要有银饰和象牙柄才能切割吗？或者我们必须锻造印度钢来切肉，如同在战斗中要武器一样？如果盆子是陶制的，它就不能接收手上的污垢吗？洗脚盆就不能接收脚上的污垢吗？镶象牙脚的桌子会因承载一个半便士的面包而愤怒吗？灯不会因为它是陶匠的作品而非金匠的作品而发光吗？我断言，矮床并不比象牙卧榻提供更差的休息；而山羊皮被单足够铺在床上，不需要紫色或猩红色的床罩。然而，由于奢侈的愚蠢——恶的源头——而谴责节俭，这是何等巨大的错误，何等无知的自负！看。主用一个普通的碗吃饭，使门徒们躺在地上草地上，用一条毛巾束腰洗他们的脚——祂，谦卑的天主，宇宙的主宰。祂没有从天上带下一个银洗脚盆。祂向撒玛黎雅妇人要求喝水，她从井里用陶器打水，并非寻求皇家的金子，而是教导我们如何轻易解渴。因为祂以使用为目的，而非奢华。祂在宴会上吃喝，不是从地下挖出金属，也不是使用金子和银子的器皿，即那些散发铁锈气味的器皿——冒烟金属的铁锈所发出的那种气味。\par

因为简而言之，在食物、衣服、器皿以及所有其他属于家的事物上，我全面地说，人必须遵循基督信徒的习俗，正如对个人、年龄、职业、生命阶段有用且适合的一样。因为作为同一天主仆人的那些人，其财产和家具应当展现同一美好生活的标记；并且每个人应个别地显现在信德中，这信德显示没有区别，实践所有其他与这统一生活方式相符合并与此单一计划和谐的事物。\par

我们轻松获得、容易使用的东西，我们赞美它、轻易保持它、自由地分享它。有用的东西更可取，因而廉价的东西优于昂贵的东西。总之，财富若管理不当，便是邪恶的堡垒；许多人盯着它，就永远无法到达天国，因为他们贪恋世俗之物，并因奢侈而傲慢地生活。但那些认真追求救恩的人，必须预先在心中定下这一点：\Quote{我们所拥有的一切都是赐给我们使用的，而使用是为了足够，这一点通过少量东西就能达到。}因为那些出于贪婪而对自己囤积之物沾沾自喜的人是愚昧的。经上说：\Quote{那聚敛工钱的，是把工钱装入破漏的口袋。}\BibleRef{盖 1:6} 那聚敛粮食并封存起来的人就是这样；那什么也不给别人的人，变得更穷了。\par

有一种可笑的滑稽场面，叫人忍俊不禁：人们带进银制的尿壶和水晶的\Emph{夜壶}，同时请来他们的顾问；愚蠢的富婆则制作金质的粪便容器；这样，她们即使富有，也不能用高雅的方式解手。我但愿她们整个一生都认为金子适合做粪土。\par

但如今，贪财被证明是邪恶的堡垒，宗徒说：\Quote{贪财是万恶的根源；有些人因贪慕钱财而偏离了信德，并用许多痛苦刺透了自己。}\BibleRef{弟前 6:10}\par

但最好的财富是欲望的贫穷；真正的宽宏大度不是以财富为傲，而是藐视财富。夸耀自己的器皿是极其卑鄙的。因为过分在意任何人都可以在市场上买到的东西，显然是错误的。但智慧并非用世上的钱币购得，也不是在市场出售，而是在天堂。它以真正的钱币——不朽的圣言、王者的黄金——出售。\par

% structure: p0055 source=h2 target=subsection reason=relative-heading-rank; base=h2

\subsection{第四章 如何在宴会上行为举止}

让狂欢远离我们理性的宴饮，也让那些在放纵中狂欢的愚蠢熬夜远离。因为狂欢是一支令人沉醉的笛子，是爱欲之桥的锁链，即悲伤的锁链。愿爱、醉酒和愚昧的情欲从我们的队列中被除去。粗俗的歌唱是醉酒的同伙。通宵饮酒招致醉酒，激起淫欲，并在可耻的行为中胆大妄为。因为如果人们把时间花在笛子、瑟琴、合唱、舞蹈、埃及式的拍手以及诸如此类无序的轻浮事上，他们就变得非常不庄重、难以管束，敲钹击鼓，在迷惑的乐器上发出噪音；因为显然，这样的宴席在我看来就是醉酒的剧场。因为宗徒规定：\Quote{脱去黑暗的行为，穿上光明的武器，行事为人要端正，好像在白昼，不可荒宴醉酒，不可好色邪荡。}\BibleRef{罗 13:12} 让笛子留给牧人，让笛子留给那些沉溺于偶像崇拜的迷信之人。因为实际上，这类乐器应被排除在节制的宴席之外，它们更适合野兽而非人类，更适合人类中较不理性的部分。因为我们听说，牡鹿会被笛子迷住，被音乐诱入猎人的网罗。当母马交配时，会吹奏一首笛曲——一首婚礼之歌。总之，一切不当的视听，以及每一种放荡的可耻感觉——事实上，这是感觉的缺失——都必须排除；我们必须提防任何取悦耳目、使人软弱的快感。因为卡里亚缪斯的破碎曲调和哀婉音调的各种魔力败坏了人的道德，通过放荡而有害的音乐艺术，将人引向心灵的扰乱。\par

圣神将神圣的崇拜与这种狂欢区分开来，歌唱道：\Quote{要用号角的声音赞美祂；}因为祂要用号角声唤醒死者。\Quote{要用瑟琴赞美祂；}因为舌头是主的瑟琴。\Quote{要用竖琴赞美祂。}竖琴指的是被圣神击打的嘴，好像被拨子拨动一样。\Quote{要用鼓和舞蹈赞美祂，}指的是教会默想那在响亮的皮上死者的复活。\Quote{要用弦乐和风琴赞美祂。}祂称我们的身体为风琴，其神经为琴弦，身体由此获得和谐的张力，当被圣神击打时，就发出人的声音。\Quote{要用响亮的钹赞美祂。}祂称舌头为口的钹，因嘴唇的搏动而发出响声。因此祂向人类呼喊：\Quote{凡有气息的，都要赞美主！}因为祂关心祂所造的每一个有气息之物。因为人确实是一种和平的乐器；而其他乐器，你若探究，会发现它们是好战的，煽动情欲，或点燃爱情，或激起愤怒。\par

因此，在战争中，埃特鲁里亚人用喇叭，阿卡迪亚人用笛子，西西里人用\OriginalTerm{佩克提斯}{pectides}，克里特人用竖琴，拉刻代蒙人用长笛，色雷斯人用号角，埃及人用鼓，阿拉伯人用铙钹。我们所用的唯一的和平乐器就是圣言，我们用它来荣耀天主。我们不再使用古代的瑟琴、喇叭、鼓和笛子，那些擅长战争并藐视敬畏天主的人过去也在他们节庆集会的合唱中使用它们，企图用那样的曲调振奋他们沮丧的心灵。但让我们的饮酒中的愉悅按法律成为双重的。因为\Quote{你若爱主你的天主，}然后\Quote{爱你的邻舍，}就让它首先向天主表达感恩和诗篇，其次向邻舍表达得体的团契。因为宗徒说：\Quote{让主的圣言丰丰富富地住在你们心里。}\BibleRef{哥 3:16} 而这圣言使自己适应时节、人物和地方。\par

当前，祂是我们中的一位宾客。因为宗徒又补充说：\Quote{在一切智慧中，用圣咏、赞美诗和属神的歌曲彼此教导和劝勉，怀着恩宠在心中向主歌唱。}又说：\Quote{无论你们在言语或行为上做什么，都要因主耶稣的名而行，藉着祂感谢天主圣父。}这就是我们感恩的欢宴。即使你愿意弹奏竖琴或里拉琴歌唱，也无可指责。你要效法那位正直的希伯来君王的感恩赞颂。\Quote{义人哪，你们应当靠主喜乐；正直人的赞美是合宜的，}预言说。\Quote{要弹琴称谢主，用十弦瑟歌颂祂；向祂唱新歌。}那十弦瑟岂不指示那由十的元素所彰显的圣言耶稣吗？在领受食物之前，理应祝福万物的创造者；同样，在饮酒时，领受祂的受造物而赞美祂也是适宜的。因为圣咏是一种和谐而清醒的祝福。宗徒称圣咏为\Quote{属神的歌曲}。见：\BibleRef{弗 5:19}\par

最后，在睡眠之前，感谢天主是一种神圣的义务，因为我们已经享受了祂的恩宠和爱，然后就径直入睡。他说：\Quote{要在唇舌的歌曲中向祂称颂，因为在祂的命令中，祂的一切美意都得以成就，祂的救恩毫无欠缺。}\par

此外，在古希腊人当中，在他们充满酒杯的宴会上，他们唱一首被称为\OriginalTerm{饮酒歌}{skolion}的歌曲，仿照希伯来圣咏的方式，大家一起高声唱\OriginalTerm{赞歌}{pæan}，有时也在轮流饮酒时轮流唱歌；而那些更有音乐天赋的人则弹奏里拉琴歌唱。但要远避情歌，让我们的歌曲成为赞美天主的赞美诗。他说：\Quote{让他们赞美，}经上说，\Quote{在舞蹈中赞美祂的名，用铃鼓和瑟弹奏歌颂祂。}那弹奏的合唱团是什么？圣神会指示你：\Quote{愿赞美在圣者的集会（教会）中归于祂；愿他们因他们的君王而欢乐。}他又补充说：\Quote{主必喜悦祂的子民。}温和的曲调可以接纳；但我们要尽可能从我们刚毅的心灵中远避那些柔媚的曲调，它们通过音调中有害的技巧，使人变得柔弱和粗俗。庄重而节制的旋律告别了醉酒的骚乱。因此，半音曲调应留给不节的狂欢和浮华媚俗的音乐。\par

% structure: p0062 source=h2 target=subsection reason=relative-heading-rank; base=h2

\subsection{第五章 论欢笑}

那些模仿可笑感觉——或者更准确地说，模仿应受嘲笑之事——的人，应当被逐出我们的社会。\par

因为既然一切形式的言语都源于心灵和品行，那么可笑的话语若不是源于可笑的行为，就不会被说出。因为这句话：\Quote{不是好树结坏果子，也不是坏树结好果子}\BibleRef{玛 7:18}正适用于此。因为言语是心灵的果子。因此，如果滑稽之人要被逐出我们的团体，我们自己绝不可被视为挑动欢笑的人。因为如果我们被禁止聆听的事物，却发现自己在模仿它们，那是荒谬的；而人若使自己成为笑柄——即成为侮辱和讥讽的对象——则更加荒谬。如果我们不能忍受像有些人在游行中那样做出滑稽的体态，我们又如何能正当地容忍内在之人被当面取笑呢？因此，我们绝不应主动表现出滑稽的性格。那么，我们怎能致力于使自己在谈话中显得可笑，从而歪曲言语——这人类最宝贵的禀赋呢？所以，这样做是可耻的；因为这类滑稽之人的谈话不适合我们的耳朵，因为那些表达方式使我们熟悉可耻的行为。\par

戏谑是可以允许的，但不可打诨。此外，即使欢笑也须加以节制；因为以正当方式表达欢笑显示得体，而若以不同方式发出则显示缺乏自制。\par

总之，凡是人性中自然的事物，我们不可根除，而应加以限制并定下适当的时候。因为人不是在任何情形下都该发笑，因为他是会笑的动物，正如马不是在任何情形下都该嘶鸣，因为它是会嘶鸣的动物。但作为理性受造物，我们应适度地调节自己，和谐地放松我们严肃追求中的严峻和过分紧张，而非不和谐地将之一并打破。\par

因为面容和谐得体的松弛——如同乐器一般——被称为微笑。在自律的人脸上，欢笑亦如此称呼。但在女人身上，不和谐的松弛被称为咯咯笑，是色情的笑；在男人身上，称为哄笑，是野蛮而侮辱的笑。经上说：\Quote{愚人在欢笑中提高他的声音}\BibleRef{德 21:20}；而聪明的人几乎不为人察觉地微笑。这里他称聪明人为智者，因为他和愚人不同。但另一方面，人不必阴郁，只需庄重。因为我确实更喜欢一个面容严厉的人微笑，而非相反；因为这样他的笑就不太会成为嘲笑的对象。\par

微笑本身需要加以训练。如果微笑针对可耻之事，我们应当脸红而非微笑，免得我们因同情而显得以此为乐；如果针对痛苦之事，更应当显得忧伤而非高兴。因为前者是理性思考的标记；后者则暗示残忍之嫌。\par

我们不可不停地笑，因为那会越界；也不可在长者或应受尊敬的人面前笑，除非他们为使我们高兴而说笑。我们也不可在所有人面前、在所有地方、对所有人、为所有事而笑。因为对于儿童和妇女而言，笑尤其容易成为引人绊倒的原因。甚至显出严肃的表情也能使周围的人保持距离。因为庄重仅凭一瞥就能阻止放荡的接近。总之，所有无理智的人，简而言之，酒\par

\Quote{命令人纵情大笑和跳舞，}\par

将柔弱的举止变为淫逸。我们也必须考虑，随后言论自由如何将不当行为引向污言秽语。\par

\Quote{他说了一句不该说的话。}\par

因此，尤其在酒中，狡猾之人的品性常常被看穿，因为他们借醉酒的可耻放纵剥下了面具；藉此，理性在灵魂本身中被醉意压垮而沉睡，无序的情欲被唤醒，制服了心灵的软弱。\par

% structure: p0074 source=h2 target=subsection reason=relative-heading-rank; base=h2

\subsection{第六章 论污言秽语}

我们自身必须完全戒绝污言秽语，并以严肃的眼神、转过脸去，以及我们所谓\Quote{嘲弄}来堵住从事此事之人的口；也常用更严厉的言辞。祂说：\BibleQuote{凡从口里出来的，} \BibleRef{玛 15:18} \BibleQuote{污秽人，}——显示出他是不洁的、外邦的、未受训练的、放荡的、不蒙拣选的、不合适的、不体面的、不节制的。\par

鉴于在听闻和观看猥亵之事上适用类似的规则，神圣的训导者循着二者的同一途径，以贞洁的言语作为护耳装备那些在斗争中奋斗的子女，以免淫念的冲击穿透到灵魂的损伤；祂又引导眼睛观看体面之事，说：失足于脚比失足于眼更好。宗徒击退这种污言秽语，说：\BibleQuote{不可出口污秽的言语，只说善言。} \BibleRef{弗 4:29} 又说：\BibleQuote{正如圣徒所应有的，淫乱、一切污秽、并贪婪的话，在你们中间连提都不可，相宜的乃是感谢。} \BibleRef{弗 5:3} 如果\BibleQuote{凡称兄弟为傻子的，难免受裁判，}那么对于说愚妄话的人，我们该如何评论呢？关于这样的人，不是写着：\BibleQuote{凡说闲话的，在审判之日要交账于主吗？} 又说：\BibleQuote{凭你的言语，你将被成义，}祂说，\BibleQuote{凭你的言语，你将被定罪。} \BibleRef{玛 12:37} 那么，什么是救命的护耳，什么是滑溜眼睛的规则？乃是与义人的交谈，预先占据并武装耳朵，以抵抗那些引人离开真理的人。\par

\Quote{恶友败坏善行，}\par

诗歌如是说。宗徒更高贵地说：\BibleQuote{恶要厌恶，善要亲近。} \BibleRef{罗 12:9} 因为与圣徒交往的人必将被圣化。我们必须完全戒绝那些从耳朵听到的可耻之事，以及言语和景象。更必须保持纯洁，远离可耻的行为：一方面，不展示和暴露我们不应暴露的身体部分；另一方面，不看禁止观看的事物。因为孝顺的儿子不忍观看义人羞耻的暴露；而羞愧遮掩了醉酒所暴露的——那无知之罪的景象。\BibleRef{创 9:23} 同样，我们应当远离诽谤的传闻，信了基督之人的耳朵应当对之关闭。\par

正因如此，在我看来，训导者不容许我们说任何不体面的话，在早期就加强我们抵抗放荡。因为祂始终擅长斩断罪根，例如，藉着\Quote{不可奸淫}，祂说\Quote{不可贪恋。}因为奸淫是贪恋的果实，而贪恋是邪恶的根。同样，在此事例中，训导者也指责名称上的放荡，从而剪断了放荡的交往。因为名称上的放荡产生行为不端的欲望；而名称上的节制是对抵抗淫荡的训练。我们在更详尽的论著中已表明，无论在名称上还是在那些用非常用名称称呼的肢体上，都没有真正猥亵的指称。\par

因为膝和腿，以及类似的肢体，以及它们的名称和活动，都不是猥亵的。就连\OriginalTerm{私处}{pudenda}也应视为提示端庄的对象，而非羞耻之物。只是它们不合法的活动才是可耻的，应受侮辱、责备和惩罚。因为实际上唯一可耻的事是邪恶，以及因邪恶而行的事。根据这些言论，谈论邪恶的行为，如谈及奸淫、男色之类，就适当地被称为污秽之谈。轻浮的唠叨也该被制止。因为经上说：\BibleQuote{多言多语难免不犯罪。} \BibleRef{箴 10:19} \Quote{因此，口舌的罪必受惩罚。} \BibleQuote{有人沉默却被视为智慧；有人因多言而为人憎恶。} \BibleRef{德 20:5} 但更甚的是，喋喋不休的人使自己成为厌恶的对象。\BibleQuote{因为多言的人，憎恶自己的灵魂。} \BibleRef{德 20:8}\par

% structure: p0081 source=h2 target=subsection reason=relative-heading-rank; base=h2

\subsection{第七章 同居者的指引}

让我们远离嘲弄，它是侮辱的起源；由它爆发纷争、争论和敌意。我们说过，侮辱是醉酒的仆人。判断一个人，不仅凭他的行为，也凭他的言语。\BibleQuote{在宴会上，}据说，\BibleQuote{不要斥责你的邻人，也不要对他说一句指责的话。}\BibleRef{德 31:31}因为我们特别被命令要与圣人来往，嘲弄圣人便是罪：\BibleQuote{因为从愚昧人的口中，}圣经说，\BibleQuote{是侮辱的棍杖，}\BibleRef{箴 14:3}——以棍杖意指侮辱的支撑，侮辱倚靠它并安息其上。因此，我钦佩那位宗徒，他为此劝诫我们不要说\BibleQuote{猥亵或不恰当的话。}\BibleRef{弗 5:4}因为如果节日的聚会是因友爱而举行，宴会的目的是对与会者的友善，而饮食伴随着友爱，那么对话岂不应以理性方式进行，并且出于友爱避免以问题困惑他人吗？因为如果我们相聚是为了增进彼此的好意，为何要借嘲弄挑起敌意呢？沉默胜于反驳，从而在无知上增加罪过。\BibleQuote{有福的，}真实地说，\BibleQuote{是那没有用口失言，也没有被罪的痛苦刺透的人；}\BibleRef{德 14:1}或是为他所说错话而忏悔，或说话不伤害任何人。总之，让青年男女完全远离这样的宴会，以免他们在不合适的事上失足。因为那些他们耳朵不习惯的事物和不宜的视线，点燃他们的心，而内在的信德仍摇摆不定；他们年龄的不稳定性合谋使他们轻易被情欲冲昏头脑。有时他们也是他人跌倒的原因，因展示其生命时期危险的魅力。因为智慧似乎良好地训诫：\BibleQuote{绝对不要与有夫之妇同坐，也不要用肘与她一同斜倚；}就是说，不要常与她一起晚餐或吃饭。因此她补充说：\BibleQuote{也不要与她在酒中作伴，免得你的心倾向她，且因你的血滑向灭亡。}\BibleRef{德 9:9}因为醉酒的放纵是危险的，且易于玷污。她称她为\Quote{有夫之妇}，因为那试图破坏婚姻之锁链的人危险更大。\par

然而若有什么必要的事，要求已婚妇女出席，就让她们衣着得体——外表以衣服，内心以端庄。至于未婚女子，她们出现在男人的宴席上，尤其是那些有醉意的人当中，是极大的丑闻。并且让男人们只以耳朵在场，眼睛盯着卧榻，双肘支撑着不动；如果他们坐着，不要交叉双脚，不要把一条腿搭在另一条腿上，也不要把手托在下巴上。因为不自持是粗俗的，对年轻人来说是一个缺点。不停地移动和改变姿势是轻浮的表现。一个节制的人在饮食上也应当少量取用，从容不迫，不要急切，无论是在开始还是在上菜过程中，并且及时停止，以此表明他的超然态度。\Quote{吃}，经上说，\Quote{像一个男子，吃摆在你面前的。为了养生的缘故，要首先停止；如果坐在许多人中间，不要在他们之前伸手。} \BibleRef{德 31:16} 你绝不可在暴食的冲动下猛冲上前；也不可虽渴望，却迟迟不伸手，因为贪婪表现出不受控制的胃口。在宴席中，你也不可像野兽一样紧抱着食物，也不可取用过多的酱汁，因为人天生不是吃酱汁的，而是吃面包的。一个节制的人还必须在众人之前起身，安静地从宴会退场。\Quote{因为起身的时候，}经上说，\Quote{不要最后；赶快回家。} \BibleRef{德 32:11} 十二位宗徒召集了众多门徒，说：\Quote{我们撇下天主圣言而去伺候饭食，是不相宜的。} \BibleRef{宗 6:2} 如果他们回避了这事，那么他们更会避免贪吃。宗徒们自己也写信给安提约基雅、叙利亚和基里基雅的弟兄们说：\Quote{圣神和我们决定，不将别的重担加在你们身上，惟有这几件事是必要的：就是禁戒祭偶像的物和血，并勒死的牲畜和奸淫；你们若保守自己不沾染这些，就必做得好。} 但我们必须防备醉酒如同防备毒芹；因为两者都导致死亡。我们也要控制过度的欢笑和不能节制的眼泪。因为常常在酒的影响下，人们过度大笑之后，不知怎地，由于某种醉意的冲动而流泪；因为软弱和粗暴都与圣言不合。年长的人视年轻人为孩子，可以（虽然很少）与他们玩耍，开玩笑以训练他们良好的举止。例如，在一个害羞寡言的青年面前，人们可以这样打趣说：\Quote{我这个儿子（我指的是那个沉默的）总是在说话。} 这样的玩笑增强了青年的谦逊，通过戏谑地展示他的优点，批评那些他不具备的缺点。因为这个方法是有教益的，它通过缺失的东西来确认现有的东西。确实，说一个只喝水、清醒的人也会喝醉，正是此意。但如果有人喜欢嘲笑别人，我们必须保持沉默，避免多余的话，就像满杯的酒一样。因为这样的玩笑是危险的。\Quote{急躁人的口招致破碎。} \BibleRef{箴 10:14} \Quote{你不可接受愚昧的报告，也不可附和恶人做不义的见证，} \BibleRef{箴 24:28} 无论是在诽谤中，还是在伤害性的言语中，更不必说恶行了。我也认为应该对尚能自持的人说话加以限制，他们被促使向与他们谈话的人说话。\Quote{因为沉默是妇女的卓越，是青年的安全奖赏；但善于言谈是有经验、成熟的年龄的特征。老人，在宴席上说话吧，因为这对你是合宜的。但要毫无困窘地、准确地说出知识。青年，智慧也命令你。如果必须说话，要犹豫，在被问第二次之后才说，用几句话总结你的言论。} 但让双方说话者都按照适当的分寸调节他们的言辞。因为声音响亮是极愚蠢的；而声音微弱则是愚昧人的特征，因为人们听不见：前者是怯懦的标志，后者是傲慢。为无谓的得胜而争辩应当被抛弃，因为我们的目标是免于烦扰。这就是\Quote{愿你们平安}这句话的含义。在你听到之前不要回答一个字。软弱的声音是娇弱的标志。而声音的调节是智者的特征，他使自己的发声避免响亮、拖沓、急促、冗长。因为我们不应该说得太长太多，也不应该说得轻浮。我们也不能说得又快又草率。因为声音本身，可以说，应当得到公正的对待；那些吵嚷喧哗的人应当被制止。因此，智慧的尤利西斯用鞭子惩罚了忒耳西忒斯：—\par

\Quote{只有忒耳西忒斯，用无节制的言词，\newline{}他有很多，来责骂领袖们，\newline{}并不很得体，但他认为\newline{}能引动群众发笑，便大声喧嚷。}\par

\Quote{多言的人在毁灭中是可怕的。} \BibleRef{德 9:18} 对待轻浮之人，正如对待旧鞋一样：其余部分都被邪恶磨尽了，只剩下舌头用于毁灭。因此，智慧给了这些最有益处的劝诫：\Quote{不要在长老聚会中谈论琐事。} 此外，根除轻浮，从天主开始，它为我们定下律法，大致如此：\Quote{在祈祷中不要重复你的话。} \BibleRef{德 9:15} 叽喳声、口哨声以及用手指发出的用来召唤仆人的声音，都是无理性的信号，理性的人应当放弃。频繁吐痰、用力清嗓子以及在宴会上擤鼻子，都应避免。因为肯定要尊重宾客，以免他们厌恶这样的污秽，这表现了缺乏节制。因为我们不应效仿牛和驴，它们的食槽和粪堆在一起。因为许多人甚至在用餐时擤鼻子和吐痰。\par

如果有人打喷嚏，就像打嗝的情况一样，他不可用爆裂声惊吓旁边的人，从而证明他的教养不好；但打嗝应随着呼气安静地排出，嘴巴要端庄地闭合，不要像悲剧面具那样张着口打哈欠。因此，可以通过轻柔地呼吸来避免打嗝的干扰；因为这样一来，那团气的威胁症状就会以最得体的方式消散，通过控制其排出，同时也能隐藏被强行排出的空气可能带出的任何东西。想要增加噪音而非减少它，是傲慢和混乱的标志。此外，那些刮牙齿导致出血的人，既令自己不适，也为邻人所厌恶。搔耳朵和打喷嚏的刺激是猪圈的瘙痒，伴随着放纵的淫乱。无论是可耻的景象还是关于它们的可耻谈话，都应避免。眼神要稳定，颈部的转动和移动，以及谈话中的手势，都要得体。总之，基督徒以沉着、宁静、平和与平安为特征。\par

% structure: p0087 source=h2 target=subsection reason=relative-heading-rank; base=h2

\subsection{第八章　论香膏与冠冕的使用}

冠冕和香膏的使用对我们而言并非必要；因为它驱使人们追求享乐和放纵，尤其在夜间临近时。我知道那位女子在神圣的晚宴上带来了\Quote{一个玉瓶的香膏}，并抹了主的脚，使祂舒畅；我也知道古代的希伯来诸王以黄金和宝石加冕。但那女子尚未领受圣言（因为她仍是一个罪人），她用自己认为最宝贵的东西——香膏——来荣耀主；并用自己的装饰——她的头发——擦去多余的香膏，同时向主流下忏悔的眼泪：\Quote{因此，她的罪被赦免了。} \BibleRef{路 7:47}\par

这可能是主教导和受难的象征。因为用芬芳的香膏抹脚，意味着神圣的教导带着声誉传到地极。\Quote{因为他们声音传到了地极。} 如果我不显得过于坚持，那么被抹香膏的主的脚就是宗徒们，他们按照预言领受了圣神的馨香傅油。因此，那些走遍世界并宣讲福音的人，被比喻地称为主的脚，圣神也在圣咏中预言了他们：\Quote{让我们在祂脚站立的地方朝拜，}也就是说，宗徒们，祂的脚，所到达的地方；因为他们宣讲了祂，祂就来到了地极。眼泪是忏悔；散开的头发宣告了从对装饰的喜好中得释放，以及为了主而忍受宣讲所带来的痛苦，旧的虚荣因着新的信德而被除去。\par

此外，它还显示主的受难，如果你这样神秘地理解：\OriginalTerm{油}{\GreekText{ἔλαιον}}就是主自己，从祂而来的\OriginalTerm{怜悯}{\GreekText{ἔλεος}}临到我们。但香膏，是掺杂的油，就是叛徒犹达斯，主通过他被抹了脚，从而从他在世上的寄居中被释放。因为死者是被傅油的。眼泪就是我们这些悔改的罪人，我们相信了祂，祂赦免了我们的罪。散乱的头发是哀悼的耶路撒冷，被遗弃的，曾为它发出了预言性的哀歌。主自己会教导我们，这指的是诡诈的犹达斯：\Quote{同我蘸手在盘子里的，就是他要出卖我。} \BibleRef{玛 26:23} 你看到这个背信弃义的宾客，正是这个犹达斯用亲吻出卖了师傅。因为他是个伪君子，给了背信的一吻，效法了古时的另一个伪君子。祂责备那些人民，关于他们曾说过：\Quote{这人民用嘴唇尊敬我，但他们的心远离我。} \BibleRef{依 29:13} 因此，很可能是，祂用油指那个蒙受怜悯的门徒，而用污损和毒害的油指叛徒。\par

因此，受傅的脚所预言的，正是犹达斯的背叛，当时主正走向祂的苦难。救主亲自为门徒们洗脚，\BibleRef{若 13:5}并派遣他们去行善事，指明了他们为外邦人利益而行的朝圣之旅，藉祂的能力预先使他们纯洁美好。随后香膏向他们散发芬芳，那馨香的工作传至四方，被宣告出来；因为主的苦难使我们充满甜蜜的香气，而希伯来人则充满罪责。宗徒最清楚地表明了这一点，他说：\Quote{感谢天主，祂常使我们因基督而凯旋，并藉我们在各处显扬认识祂的芬芳。因为我们在天主面前，无论在得救的人或丧亡的人中，都是基督的馨香：对后者是致命的死味，对前者是生命的活味。}\BibleRef{格后 2:14}犹太人的君王使用黄金、宝石和斑斓的冠冕，那受傅者在头上象征性地戴着基督，却不自觉地被主的头装饰了。宝石、珍珠或翡翠，都指代圣言本身。黄金也是不朽的圣言，不容腐化的毒害。因此，贤士们在他诞生时献上黄金，作为王权的象征。而这冠冕，按照主的肖像，永不凋谢如花。\par

我也知道昔兰尼的亚里士提波的话。亚里士提波是一个奢华的人。他以如下言词要求对一道诡辩命题作答：\Quote{一匹马抹了香膏，在它作为马的卓越性上并无损害；一条狗抹了香膏，在它作为狗的卓越性上也无损害；人也是如此，}他补充道，然后结束。但狗和马并不在意香膏，而对于那些知觉更理性的人，将女性的香气涂在自己身上，其使用更应受谴责。这些香膏有无数种类，如布伦香、梅塔利香、王香；埃及的普朗贡香和普萨格迪香。西摩尼得斯毫无羞耻地在抑扬格诗句中说：\par

\Quote{我被香膏和香水涂抹，\newline{}还有甘松油脂。}\par

因为有一个商人当时在场。他们还使用从百合花和柏树制成的香膏。甘松在他们当中很受重视，还有由玫瑰和其他妇女常用的香膏，包括湿的和干的，用于涂抹和熏蒸的香料；因为日复一日，他们的思想都指向满足无尽的欲望，指向无穷无尽的香气种类。因此，他们也散发着过度的奢华气息。他们熏蒸和喷洒衣服、床铺和房屋。奢华几乎迫使最卑贱的器皿也散发着香水味。\par

有些人对此感到厌烦，在我看来，他们因香水使人变得女性化而反感是正确的，以至于他们禁止在治理良好的城邦中配制和售卖香水，也禁止染制花彩羊毛的染工。因为诱人的衣物和香膏不应被允许进入真理之城；但对于属于我们的人，极为必要的是，他们散发出来的不是香膏的气味，而是高贵与良善的气味。让妇女散发真正的君王香膏的气味，即基督的气味，而不是香膏和香粉的气味；让她永远沐浴在端庄的神圣香膏中，并在圣神这圣洁的香膏中喜乐。基督为祂的门徒预备了这种芬芳的香膏，以天上的芳香成分调制。\par

因此，主自己也用香膏受傅，正如达味所说：\Quote{因此，天主，你的天主，以喜乐之油傅了你，胜过你的同伴；你的衣服有没药、沉香和肉桂。}但让我们不要无意识地厌恶香膏，如同秃鹫或甲虫（据说它们被涂上香膏就会死去）；妇女应选用少量香膏，那些不会令丈夫反感的。因为过分涂抹香膏带着葬礼的气味，而非婚姻生活的气味。然而，油本身对蜜蜂和昆虫有害；它有益于一些人，却召唤另一些人去战斗；那些先前是朋友的人，一旦被涂上油，就会转为致命的争斗。\par

香膏是光滑的油，难道你不认为它能使高尚的举止变得女性化吗？确实如此。正如我们放弃了味觉上的奢侈，同样我们也放弃了视觉和嗅觉上的纵欲；免得通过感官，如同通过无人看守的门户，我们无意中将已驱走的过度放纵放进灵魂。因此，如果我们说主，伟大的大司祭，向天主献上芬芳的馨香，让我们不要以为这是祭献和馨香之香；而是让我们明白，主将爱的可悦祭献，即属灵的芬芳，放在祭台上。\par

总之：油本身足以润滑皮肤、松弛神经，并除去身体上的任何浓重气味，如果我们为此目的使用油的话。但对甜香的关注是引我们进入感官欲望的诱饵。因为放荡的人处处被牵引，通过他的食物、床榻、谈话，通过他的眼睛、耳朵、下颚，甚至通过他的鼻孔。正如牛被环和绳索牵引，纵欲者也被熏香、香膏和花冠的甜香牵引。但由于我们不给快乐以任何位置，除非它与生活的某种有用性相连，那么让我们在这里也加以区分，选择有用的东西。因为有些甜香既不会使头沉重，也不会挑动情欲，不散发拥抱和放荡陪伴的气味，而是带着适度，有益健康，滋养病中的大脑，强化胃部。因此，当一个人希望柔韧神经时，不应以花朵来凉爽自己。因为它们的用途并非完全不可抛弃，但香膏应作为药物和帮助来使用，以恢复衰弱的力量，并用于治疗感冒、伤风和倦怠，正如喜剧诗人所说：——\par

\Quote{抹油在鼻孔上，乃是至关重要的事，因为用良好的香气充满大脑有益于健康。}\newline{}\par

用温热或清凉膏油的油脂搓脚，因其有益效果而实行；因此，在那些如此浸透的人身上，会从头部向较低部位发生吸引和流动。但那些无益处的享乐，则令人怀疑有娼妓的习气，是激发情欲的药物。用油膏涂抹自己与用油膏敷抹自己完全不同。前者是柔弱的，而敷抹油膏在某些情况下是有益的。哲学家\Person{亚里斯提卜}曾在被抹油膏时说：\Quote{那些可悲的淫荡者应得悲惨地死去，因为他们把油膏的效用弄臭了。}\BibleQuote{应因医生的有用而尊敬他，}圣经说，\BibleQuote{因为至高者造了他；医术来自主。}然后他补充道：\BibleQuote{调制药膏者将配制混合物，}因为油膏显然是为使用而赐予的，并非为纵欲。我们绝不应关心油膏的刺激性，而应选择其中有益之物，因为天主允许生产油以减轻人的痛苦。\par

愚昧的妇女染白发、抹头发，却因所用干性的香料而更快变白。因此，那些涂抹油膏的人也变得更干，干燥使她们头发更白。如果白发是头发的干燥，或热量的缺乏，那么干燥吸走了作为头发天然滋养的湿气，使之变白，我们如何还能保持对油膏的喜爱呢？女人们通过它试图逃避白发，却反而变白。正如嗅觉灵敏的狗通过气味追踪野兽，有节制的人通过多余的油膏香气追踪放荡的人。\par

花冠的使用也已堕落为狂欢和酗酒的场面。不要用花冠环绕我的头，因为春天在露湿的草地上消磨时光是令人愉快的，那时柔软多彩的花朵盛开，像蜜蜂一样享受天然纯净的芬芳。但用\Quote{从新鲜草地编成的花冠}来装饰自己，并在家中佩戴，对一个有节制的人是不合适的。因为用玫瑰叶、紫罗兰、百合或其他花朵填满放纵的头发，剥夺草地上的花朵，是不适宜的。因为花冠环绕头部，因其湿气和凉爽而使头发冷却。因此，医生根据生理学断定大脑是寒冷的，赞成在胸部和鼻孔尖端抹油，以便温热的蒸气轻轻通过，有益地温暖寒冷。人不应因此用花朵使自己凉爽。此外，那些戴花冠的人破坏了花朵的乐趣：因为他们既不能享受视觉，因为花冠戴在眼睛上方；也不能享受香气，因为他们把花朵放在呼吸器官之上。因为香气上升并自然散发，呼吸器官被剥夺了享受，香气被带走。正如美一样，花朵在观看时令人愉悦；通过欣赏美好事物的视觉来荣耀造物主是适宜的。使用它们是有害的，迅速消逝，随后是懊悔。很快它们的短暂性就得到证明：两者都凋谢，花朵和美。此外，触摸花朵的人被前者冷却，被后者点燃。总之，除了视觉之外享受它们是一种罪行，而非奢侈。我们这些真正追随圣经的人，应当有节制地享受，如同在乐园中。我们必须认为女人的花冠是她的丈夫，丈夫的花冠是婚姻；而婚姻的花朵是双方的子女，神圣的农夫从肉体的草地上采摘他们。\BibleQuote{子孙是老人的冠冕。}\BibleRef{箴 17:6}据说，儿女的荣耀是他们的父亲；我们的荣耀是万有的父；整个教会的冠冕是基督。如同根和植物，花朵也有各自的特性，一些有益，一些有害，一些甚至危险。常春藤是凉的；坚果发出令人麻痹的气味，正如词源所示。水仙花是一种香气浓重的花；名字证明了这一点，它使神经陷入\OriginalTerm{麻木}{\GreekText{νάρκην}}。玫瑰和紫罗兰的香气微微凉爽，缓解并预防头痛。但我们不仅不被允许与别人共饮至醉，甚至不能多喝酒，因此不需要番红花或柏树花来引导我们安睡。它们中的许多还通过其气味温暖天然寒冷的大脑，使头部的渗出物挥发。因此据说玫瑰得名（\OriginalTerm{玫瑰}{\GreekText{ῥόδον}}），因为它散发出大量（\OriginalTerm{流}{\GreekText{ῥεῦμα}}）的\OriginalTerm{气味}{\GreekText{ὀδωδή}}。因此它也迅速凋谢。\par

但在古希腊人中，花冠的使用根本不存在；因为无论是求婚者还是奢侈的费阿刻斯人都没有使用它们。但在竞技中，最初是给运动员的礼物；第二是起立鼓掌；第三是撒叶；最后是花冠，希腊在米底战争之后沉溺于奢侈。\par

因此，受圣言训练的人被禁止使用花冠；并且不要认为这居于脑中的圣言应当被束缚，不是因为花冠是狂欢放纵的象征，而是因为它已被奉献给偶像。\Person{索福克勒斯}因此称水仙花为\Quote{伟大诸神的古老冠冕，}指地生神祇；\Person{萨福}用玫瑰为缪斯加冕：——\par

\Quote{因为你们不分享从皮埃里亚来的玫瑰。}\par

他们还说，\Person{赫拉}喜爱百合，\Person{阿耳忒弥斯}喜爱桃金娘。因为如果花卉是特别为人所造，而无知的人不是取来作恰当而感恩的用途，反而滥用它们去事奉忘恩的魔鬼，我们就因良心的缘故必须远离它们。花冠是安宁的象征。为此，他们给死者加冠，也给偶像加冠，以此证实他们是死的。因为狂欢者没有花冠就不举行他们的狂欢；一旦被花环绕，他们就越发激动。我们不可与魔鬼共融。也不可按着死偶像的方式，给天主的活像加冠。因为那美丽的\OriginalTerm{不凋花}{amaranth}花冠是为善生者存留的。这花大地不能承载，唯有上天才能产生。再者，我们既听说主被戴上了荆棘冠冕\BibleRef{玛 27:29}，却用花来给自己加冠，如此侮辱主的圣受难，这是不合理的。因为主的冠冕预言性地指向我们，我们曾经荒芜，但藉着教会（祂是元首）被安置在祂周围。但它也是信德的预象：就木料的实质而言，是生命的预象；就冠冕的名称而言，是喜乐的预象；就荆棘而言，是危险的预象，因为不流血就不能亲近圣言。但这编织的冠冕会褪色，邪恶的编织会解开，花会枯萎。因为那些不信主的人的荣耀会消退。他们给高举的耶稣加冠，为自己的无知作证。因为他们心硬，不明白他们所谓的主的羞辱，正是一明智宣告的预言：\BibleQuote{人民不认识上主}\BibleRef{依 1:3}——这人民犯了错误，理智未受割损，黑暗未被光照，不认识天主，否认上主，丧失了真正以色列的地位，迫害天主，希望使圣言蒙羞；而他们作为罪犯钉死的这位，却加冕为君王。因此，他们所不信的那人，他们将要认识祂是仁爱的天主、上主、公义者。他们激怒祂显明自己是上主，当祂被高举时，他们用常青的荆棘作为正义的冠冕环绕祂——那超乎万名之上的名——来为祂作证。这冠冕对图谋反对祂的人是敌对的，约束他们；对组成教会的人是友善的，保护他们。这冠冕是信从受光荣者之人的花朵，却用血覆盖并惩罚那些不信的人。它也是主胜利工作的象征，祂将我们所有的过犯——那些刺伤我们的过犯——都担在祂头上，即祂身体的首部。因为祂以自己的受难解救我们脱离过犯、罪恶及诸如此类的荆棘；毁灭了魔鬼后，祂胜利地、恰当地说：\BibleQuote{死亡，你的刺在哪里？}\BibleRef{格前 15:55}我们从荆棘采葡萄，从蒺藜采无花果；而对那些祂伸出双手的人——那悖逆不结果实的百姓——祂却使他们伤痕累累。我还能向你们显示其中另一个神秘意义。因为当宇宙的全能上主开始藉着圣言立法，并愿意向梅瑟显明祂的权能时，一个神性的、取了形状的光之异象在燃烧的荆棘（荆棘是一种多刺植物）中向他显示；但当圣言结束颁布法律及与人同在时，主又神秘地戴上荆棘冠冕。当祂离开这个世界，回到祂所来的地方时，祂重复了祂古老下降的开端，为的是那起初在荆棘中显现，后来被荆棘冠冕所高举的圣言，能显示这一切都是同一权能的工作，祂自身是独一的，圣父的圣子，祂确实是独一的，时间的起始与终结。\par

但是，我已经偏离了训育者的说话风格，引入了教导的风格。因此我回到主题。\par

那么，总结一下：我们已经表明，在医学领域，为了治愈，有时也为了适度的娱乐，从花中获得的愉悦以及从香膏和香水中获得的益处，是不应被忽视的。如果有人问：那么，对于不使用花的人，花有什么乐趣呢？让他们知道，香膏是由花制成的，而且非常有用。\OriginalTerm{苏西尼安香膏}{Susinian ointment}由多种百合制成；它有温热、通利、吸引、湿润、净化、细微、消胆、软化等功效。\OriginalTerm{纳西苏斯香膏}{Narcissinian}由水仙制成，与苏西尼安香膏同样有益。\OriginalTerm{桃金娘香膏}{Myrsinian}由桃金娘和桃金娘果制成，是一种收敛剂，能止身体的溢液；而玫瑰香膏则是清凉的。总之，这些也都是为我们的使用而造的。经上说：\BibleQuote{请听我，要像溪水旁种植的玫瑰一样生长，要像乳香一样散发芬芳，并为祂的作为赞美上主。}\BibleRef{德 39:13}关于花卉和香气是为必要的用途而非奢侈的过分而造，我们本有许多话可说。如果必须让步，人们享受花的芬芳就够了，但不要用它们来给自己加冠。因为圣父深切关怀人，唯独将祂自己的技艺赐给人。因此圣经说：\BibleQuote{水、火、铁、奶、细面、蜜、葡萄汁、油和衣服——这一切都是为敬虔人的益处。}\BibleRef{德 39:26}\par

% structure: p0109 source=h2 target=subsection reason=relative-heading-rank; base=h2

\subsection{第九章 论睡眠}

现在我们必须说明，在适当的时候，我们应如何就寝，以纪念节制的规诫。因为饭食之后，为分享的乐趣和愉快度过的一天感谢天主之后，我们的谈话必须转向睡眠。华丽的床单、金线绣花的地毯、平滑的金线地毯、长长的紫色细袍、昂贵的毛绒斗篷、加工过的紫色地毯、厚绒外套，以及比睡眠还柔软的卧榻，都应摒弃。\par

因为，除了对肉欲的谴责之外，睡在羽绒床上是有害的，由于床垫的柔软，我们身体会像掉进一个张开的凹陷一样坠落。\par

因为它们不适合睡眠者在其中翻身，因为床在身体两侧隆起成山丘。它们也不适合消化食物，反而会烧灼食物，从而破坏营养。但伸展身体在平坦的卧榻上，为睡眠提供一种天然的锻炼场所，有助于消化食物。而那些能在其他床上滚动的人，以此作为睡眠天然的锻炼，更容易消化食物，使自己更能适应紧急情况。此外，银脚床榻显得极为铺张；床上的象牙，当身体离开灵魂时，对圣洁之人是不允许的，这是一种懒散的休息方式。\par

我们不应为这些事费心，因为拥有它们的人并不禁止使用它们；但过分操心是被禁止的，因为幸福并不在于它们。另一方面，像\Person{狄俄墨得}那样行事，带有犬儒式的虚荣——\par

\Quote{他伸开身体躺在一张野牛皮下面，}\par

除非环境所迫。\par

\Person{奥德修斯}用一块石头修正了婚床的不平坦。这种节俭和自助不仅由私人实行，而且由古希腊的首领实行。但何必说这些？\Person{雅各伯}睡在地上，一块石头作他的枕头；然后他配得看见那异象——那是超乎常人的。合乎理性地，我们使用的床必须简单节俭，这样构造：避免极端（过度放纵和过度忍耐），使其舒适：如果天热，则保护我们；如果天冷，则温暖我们。但床榻不要精致，要有光滑的床脚；因为精致的车削有时会为爬行动物形成路径，它们缠绕在作品的凹槽中，不会滑落。\par

尤其是，床榻适度的柔软适合男子气概；因为睡眠不应是身体的完全衰弱，而是它的放松。因此我说，睡眠不应为了放纵而降临我们，而是为了从行动中休息。因此我们必须睡得容易醒来。因为经上说：\BibleQuote{你们腰间要束上带，灯也要点着；你们好像仆人等候主人，从婚宴回来时，他来叩门，就立刻给他开门。主人来了，看见仆人警醒，那仆人就有福了。}因为一个睡着的人毫无用处，如同死人一样。因此我们应当常常夜间起来赞美天主。因为那些为他警醒的人是有福的，从而使自己像天使一样，我们称天使为\Quote{警醒者}。但一个睡着的人毫无价值，与未活者无异。\par

但那有光的人警醒，\BibleQuote{黑暗不能胜过他，}\BibleRef{若 1:5}睡眠也不能，因为黑暗不能。那被光照的人因此向天主警醒；这样的人活着。\BibleQuote{凡受造的，在他内有生命。}\BibleRef{若 1:3}\BibleQuote{有福的人，}智慧说，\BibleQuote{那听从我，每日在我门前警醒，在我门框旁守候的人。}\BibleRef{箴 8:34}\BibleQuote{所以我们不要像别人一样睡觉，而要警醒，}经上说，\BibleQuote{并要清醒。因为睡觉的人是在夜间睡，醉酒的人是在夜间醉，}即在无知的黑暗中。\BibleQuote{但我们既然是白昼的，就该清醒。因为你们都是光明之子，白昼之子；我们不属黑夜，也不属黑暗。}\BibleRef{得前 5:5}但我们中谁最关心活出真正的生命并持有高尚情感，他就会尽可能长时间地保持清醒，只在这件事上为自己保留有利于健康的部分；而这并不常见。\par

但热忱于活动在劳作之后产生永久的警醒。不要让食物压垮我们，而要使我们轻盈；好让我们尽可能少受睡眠的损伤，如同挂着重物游泳的人被拖累。相反，让节制将我们从深渊之下提升到警醒的努力中。因为睡眠的压抑如同死亡，它迫使我们陷入无知觉，通过闭合眼皮切断光明。因此，我们这些真光的儿子，不要关闭这光的大门；而是转向内在，照亮隐藏之人的眼睛，凝视真理本身，接受它的水流，让我们清晰而可理解地揭示那些真实的梦。\par

但那些灌满酒的人的打嗝，那些塞满食物的人的鼾声，在床上被褥里翻滚的鼾声，以及疼痛胃部的咕噜声，遮盖了灵魂的明见之眼，因为用无数幻想充满了心智。原因是食物过多，它将人的理性部分拖入愚钝状态。因为过多的睡眠对我们的身体和灵魂都没有益处；也完全不适合以真理为目标的过程，虽然它合乎自然。\par

现在，义人\Person{罗特}（我暂且略过重生的计划）不会被他的女儿们灌醉并受睡眠压制，否则他就不会陷入那不堪的交往。因此，如果我们切断导致嗜睡的原因，我们就会更清醒地睡眠。因为那有不眠的\OriginalTerm{圣言}{Word}住在我们里面的人，不应整夜睡觉；而应在夜间起来，特别是在白日将尽的时候，有人专心于文学，有人开始他的手艺，妇女们拿起纺锤，我们所有人都应以温和而渐进的方式与睡眠斗争，以便通过警醒，我们能在更长的时间内分享生命。\par

我们既将夜间最好的部分用于醒寤，便绝不可在白日睡眠；而毫无用处的睡意、打盹、伸懒腰和打哈欠，都是灵魂轻浮不安的表现。此外，我们还必须知道，睡眠的需要不在于灵魂，因为灵魂永不停息地活动。但身体通过向休息投降而得到舒缓，而灵魂那时不通过身体行动，而是在自身内运用理智。同样地，在正确思考的人看来，那些真实的梦境，乃是清醒灵魂的思绪，暂时不受身体情感的分心，并以最佳方式与自身商谈。因为灵魂停止自身内的活动，对它而言就是毁灭。因此，灵魂常默观天主，并藉与天主的持续交谈将醒寤注入身体，从而将人提升至与天使恩宠同等的高度，并从醒寤的实践中把握生命的永恒。\par

% structure: p0123 source=h2 target=subsection reason=relative-heading-rank; base=h2

\subsection{第十章——关于生养子女应如何处理}

至于结合的适当时间，只应留给那些已婚者来考虑；已婚者的目的和宗旨是生育子女；其结局则是子女成为好人；正如农夫撒种的原因是为了照顾获得营养，农耕的目的则是收获果实。但耕种有灵土地的那农夫远为更好：前者为食物而寻求时间，后者则为使整体长存而操劳，尽农夫之责；前者为自己栽种，后者则为天主栽种和撒种。因为祂说：\BibleQuote{你们要繁衍增多}\BibleRef{创 1:27}；在此应理解为：\Quote{藉此方式，人成为天主的肖像，因为人合作以生成人。}并非每块土地都适合接受种子；即使每块都适合，也不一定适合同一位农夫。也不可撒种在石头上，也不可凌辱那作为生殖之引导和元首的种子，它乃是同时具有内在自然理性的实体。将那些合乎自然的理性交付给不合理的通道，加以凌辱，是极其不虔敬的。因此，请看最富智慧的梅瑟如何以象征方式拒绝了不结果实的播种，他说：\BibleQuote{不可吃兔子，也不可吃鬣狗。}\BibleRef{申 14:7}祂不愿人分享那些动物的特性，也不愿人尝到与它们同等的淫欲；因为这些动物以某种疯狂的冲动追求性交。据说兔子每年因其活过的年数而有多重肛门；以此方式，禁止吃兔子表示祂劝诫避免男童之爱。而鬣狗则每年将雄性变为雌性；这表示那远离鬣狗的人不应犯奸淫。\par

好，我也同意，极富智慧的梅瑟显然藉上述禁令表示，我们不可模仿这些动物；但我不同意对那些象征性言词的这种解释。因为本性永不能被强迫改变。一旦本性被印上什么，就不能被私欲转变为相反的东西。因为私欲不是本性，私欲常会毁坏形式，而不是将其塑造成新的形状。虽然据说许多鸟类随季节变化，在颜色和声音上都有改变，例如\OriginalTerm{乌鸫}{\GreekText{κόσσυφος}}从黑色变为黄色，从歌唱鸟变为饶舌鸟。同样地，夜莺也交替改变其颜色和叫声。但它们并不改变本性本身，以致在转变中从雄性变为雌性。而是新长出的羽毛，如同新衣服，产生某种羽毛的着色，随后在严寒的冬天消逝，如同花朵褪色。同样，声音本身因寒冷受损而变弱。因为，由于外层皮肤被周围空气增厚，颈部的动脉被压迫并充满，紧紧压在气息上；这气息被严重堵塞，发出窒息的声响。而当气息与周围空气同化并在春天放松时，它便摆脱了堵塞状态，通过扩张的（此前闭塞的）动脉流动，不再唱出衰微的曲调，而是发出尖锐的鸣声；声音四处流淌，春天便成为鸟儿歌声之季。\par

因此绝不可相信鬣狗会改变本性：因为同一动物不会同时具有雌雄两性的生殖器官，正如一些人荒谬地认为的那样，他们怪异地捏造了阴阳人，并在这两种性别之间创新了一种雄雌同体的本性。但他们大错特错，因为他们没有注意到，万物之母和孕育者自然多么爱子女：因为这种动物，我说的是鬣狗，极其淫荡，在尾巴下方、排泄口之前，长有一块肉瘤，形状与女性的生殖器非常相似。但这块肉状物没有任何通道通向任何有用的部分，我说的是子宫或直肠：它只有一个大空洞，用来容纳空虚的淫欲，当用于受孕的通道被阻断时。而这一点，雄性和雌性鬣狗都长有，因为它们显著地喜欢被动性交：因为雄性交替地既主动又被动：因此雌性鬣狗极少能被找到：因为这种动物不经常怀孕，因为违反自然的满足感在它们体内大量泛滥。由于这个原因，我认为柏拉图在\WorkTitle{费德罗篇}中，在拒斥对男童的爱时，称这种爱为野兽，因为那些放纵情欲、淫荡的人咬住嚼子，像四足动物一样交合，并试图播下子女。不虔敬的人，\Quote{天主任由他们}，正如宗徒所说，\BibleRef{罗 1:26}\Quote{陷入可耻的情欲：他们的女人把自然的用途变为违反自然的用途；同样，男人也放弃了与女人的自然用途，彼此间欲火中烧，男人与男人行可耻的事，并在自己身上受到他们错谬应得的报应。}但自然甚至不允许最淫荡的动物将精液射入排泄通道：因为尿液排入膀胱，湿润的食物进入胃，眼泪进入眼睛，血液进入血管，污垢进入耳朵，鼻涕进入鼻孔：而直肠的末端连接着肛门，通过它排出粪便。因此，只有鬣狗多变的自然构想出这个多余的器官用于多余的性交，因此它有一定的凹陷，以服务于发痒的部位，但从那里凹陷就堵塞了：因为它不是被造来用于生殖的。由此对我们来说显而易见且公认的是，应当避免与男性交合、不结果实的播种、颠倒的性行为，以及那些自然无法结合的阴阳人的结合，跟随自然的引导，自然通过结构安排禁止了这一点，因为它造男性不是用来接收精液，而是用来射出精液。而\Person{耶肋米亚}，即通过他说话的圣神，当他说：\Quote{我的家成了鬣狗的洞穴}，他憎恶那由死尸组成的食物，用智慧的比喻谴责了对偶像的崇拜：因为活天主的家确实必须远离偶像。再次，\Person{梅瑟}也禁止吃兔子。因为兔子在所有时间交配，跳跃，雌兔蹲着，它从背后攻击：因为它属于从后面跳跃的动物。它每个月受孕，并超胎；它交配，又生产；生产后，立即被任何兔子交配（因为不满足于一个配偶），又在哺乳时受孕：因为它的子宫有两个腔，而不是只有一个空腔，这为交配提供了足够的空间（因为一切空虚都渴望被填满）；但发生的是，当它们怀孕时，子宫的另一部分被欲望抓住，并被情欲激动；因此它们发生超胎。因此，远离强烈的欲望、相互交合、与孕妇结合、交替交配、奸淫男童、通奸和情欲，这谜语的禁令劝诫了我们。因此\Person{梅瑟}公开地，而非通过谜语，禁止说：\Quote{不可奸淫；不可通奸；不可奸淫男童。}\BibleRef{出 20:14}因此我们必须全力遵守律法的规定，决不可做任何违法的事，也不可削弱诫命。因为邪恶的欲望被称为\OriginalTerm{放肆}{\GreekText{ὕβρις}}，即\Quote{无礼}；而\Person{柏拉图}将欲望之马称为\Quote{放肆的}，因为他读到：\Quote{你们对我成了发情的母马。}\BibleRef{耶 5:8}而前往索多玛的天使使我们认识情欲的惩罚。他们将想要侮辱他们的人连同整个城市一起烧毁，用这明显的迹象描绘了情欲的果实——火。因为古代所发生的事，正如我们之前所说，是为了警戒我们而写的，免得我们陷入同样的恶习，并谨慎避免落入类似的惩罚。应当视男孩为儿子；应将别人的妻子看作自己的女儿：因为控制快乐、支配肚腹和肚腹以下的部分，是最大的统治。如果智慧的理由不允许智者随意移动手指，如斯多葛派所承认的，那么对那些追求智慧的人来说，岂不是更应当控制用于交配的部位吗？正是由于这个原因，它似乎被命名为\Quote{羞部}，因为身体这一部分比其他任何部分更应当以羞耻来使用；自然就像允许我们适量地使用食物一样，也允许我们合宜地使用合法的婚姻，并允许我们追求生育子女。而那些追求过度的人，在违反法律的交合中伤害自己，堕入违反自然的事。首先，最重要的是，我们决不应像与女性一样与少年进行性行为。因此，受过\Person{梅瑟}教育的哲学家说，\Quote{不应在岩石和石头上播种}，\Quote{因为种子永远不会扎根而获得生殖的本性。}因此，圣言通过\Person{梅瑟}非常明确地命令：\Quote{不可与男人同寝，如同与女人同寝：这是可憎恶的事。}\BibleRef{肋 18:22}此外，\Quote{应当远离一切不属于自己的女性田地}，除了自己的妻子，这些是从\WorkTitle{圣经}中归纳的，著名的\Person{柏拉图}从那里接受了法律：\Quote{不可与你邻舍的妻子交合，在她身上玷污自己。}\BibleRef{肋 18:20}而妾的种子是无效的和奸淫的。不要在你不想让播种的东西生长的地方播种。除了你自己的妻子外，不要碰任何女人，只从她那里你才可获取肉体的快乐，以接受合法的继承。因为只有这些在圣言看来是合法的。确实，那些参与天主赋予生产工作恩赐的人，不应该抛弃种子，也不应该侮辱它，也不应该像把种子交给角一样播种。因此，\Person{梅瑟}本人也禁止与经期的妻子交合。因为用身体的污秽玷污生殖的种子和即将成为人的东西是不适当的，也不应该用肮脏的排泄物和污秽覆盖它；因为种子会退化，变得无用，如果它被剥夺了母牛的垄沟。他从不允许任何一个古代希伯来人与其怀孕的妻子交合。因为即使在婚姻中，仅仅为了快乐而交合也是违法、不义和违背理性的。此外，\Person{梅瑟}使男人远离孕妇，直到她们生产。因为子宫位于膀胱下方，位于直肠之上，它在膀胱的肩部之间伸出颈部；而颈部开口在精液进入时被填满而闭合，当分娩后子宫又变空：果实排出后，然后接受种子。对我们来说，提及胎儿受孕的器官并不羞耻，因为天主在创造它们时并不羞耻。因此，渴望生育的子宫接受种子，并拒绝可耻和应受谴责的交合，在播种后关闭口部，完全排除情欲。而它的欲望，先前专注于友好的拥抱，转而向内，忙于生育后代，与造物主一同工作。因此，在自然已经工作时再烦扰它是不对的，因为过度地爆发到任性的情欲中。而\Quote{放肆}有很多名称和很多种类，当它转向这种色情的不节制时，被称为\OriginalTerm{好色}{\GreekText{λαγνεία}}，即\Quote{好色}；这个词表示好淫的、公开的、乱伦的性欲倾向：当它增长时，同时聚集了一大群疾病，如对佳肴的渴望、酗酒和对女人的爱慕；还有奢侈，以及同时对所有快乐的追求：在这些事情上，欲望实行暴政。与这些相关的无数情感增长，从中不节制者的品行被推向极端。\WorkTitle{圣经}说：\Quote{为不节制者准备了鞭子，刑罚加在愚昧人的背上。}\BibleRef{箴 19:29}它将不节制的力量和持续的忍耐称为\Quote{愚昧人的背}。因此，\Quote{从你的仆人身上移去空虚和可耻的希望}，他说，\Quote{使不当的欲望远离我。愿食欲和性欲不要抓住我。}\par

因此，必须远远避开那些奸诈之徒的种种恶行；因为不仅\OriginalTerm{克拉特斯的袋子}{Cratetis Peram}，连我们的城邦也不会有愚蠢的食客、耽于后庭之乐的纵欲好色者、诡诈的娼妓，以及其他任何这类享乐的禽兽航行到此。因此，在我们整个生命中，应当播撒良善且正直的忙碌。总之，要么结婚，要么完全保持独身；因为问题就在于这一点，这在我们\WorkTitle{论节欲}一书中已经阐明。但如果连\Quote{是否该娶妻}这问题本身都成为思考的对象，那么，怎能允许人像使用饮食一样随意使用性交，仿佛那是必需品呢？由此可见，神经纤维会像经线一样被拉伸，并在激烈的交合中崩断。不仅如此，它还会遮蔽感官，削弱精力。这在无理性的动物和从事锻炼的身体上都很明显：那些禁欲者能在竞赛中战胜对手；而动物被拉离交配后，会被拖拽，几乎被牵走，所有力量和冲劲最终都仿佛被削弱了。那位\OriginalTerm{阿布德拉的诡辩家}{sophista Abderites}称\Quote{性交}为\Quote{小癫痫}，认为它是一种不治之症。难道那归因于排空虚耗的虚弱不会随之而来吗？\Quote{人由人而生，并从人中被拔出。}请看损失的巨大：整个人通过性交的排空而被抽出。因为经上说：\BibleQuote{这真是我骨中的骨，肉中的肉。}\BibleRef{创 2:23}因此，人因精子而排空的程度，正如他身体所显示的那样；因为所离开的就是生成的开始；不仅如此，物质的沸腾还会扰乱、动摇和撼动身体的构造。因此，当有人问他\Quote{对性事感觉如何}时，他机智地回答：\Quote{别提了，我很乐意从那里逃开，如同从一个野蛮疯狂的主人那里。}然而，婚姻确实可以被允许和接纳：因为主愿意人类繁衍；但他并没有说：\Quote{你们要纵欲；}他也没有意愿你们好像是为性交而生的，沉溺于享乐。让我们的导师以羞耻感来激励我们，通过\Person{厄则克耳}呼喊：\Quote{割除你们的淫乱！}无理性的动物也有适宜的播种季节。以非生育为目的的性交，是对本性的伤害；我们应当以本性为导师，谨慎观察它明智引入的适宜时机：我指的是老年和幼年。因为它尚未允许这些年龄，而它也不愿那些老人再娶妻。但主不愿人总是致力于婚姻。婚姻是生儿育女的渴望，而非无序的精液排泄，那是违法的、不合理的。我们整个生命若从起初就能克制欲望，并且不通过邪恶诡诈的技艺来消灭那因天主上智而生的人类，那么它就会合乎自然地前进；因为那些技艺，为了隐藏淫乱，使用了致命的药物，这些药物导致彻底毁灭，与胎儿一同消灭了整个的人性。此外，对于那些被允许娶妻的人，他们需要导师，以便不在白天举行本性的神秘仪式，也不要在早上从教会或市场回来时像公鸡一样交合，因为那时是祈祷、阅读以及白天应做的工作的时间。到了晚上，应在宴饮后安歇，并在为了所领受的美好事物向天主谢恩之后。本性并不总是给婚姻交合留出时间；结合越持久就越值得渴望。也不应在夜间如同在黑暗中那样无节制地行事，而应将羞耻心——它是理性之光——纳入灵魂。如果我们白天编织节制的教义，晚上上床时却将其解开，那么我们就与织布的\Person{佩内洛普}毫无区别了。因为如果必须实践正直，那么更应向你的妻子展示正直，避免不正直的结合；而且，你与近人交往的贞洁，应有家庭内可信的见证。因为在她眼中，若正直不是由那些强烈的欢愉作为确定的见证所证实，就不可能被认为是正直。然而，对交合的冲动所生发的善意，只能短暂繁荣，并随着身体衰老；有时甚至会提前衰老，当婚姻的节制被娼妓般的淫欲败坏时。因为爱人的心是轻浮的，爱情的刺激常被悔恨扑灭；当满足感感到谴责时，爱常常变为恨。至于不洁的言语、可耻的姿态、娼妓般的亲吻和诸如此类的放荡行为，连名字都不该提起，因为我们要跟随有福的宗徒\Person{保禄}，他曾公开说：\BibleQuote{至于淫乱和一切不洁之事，或贪求财物的事，在你们中连提都不该提，正如圣者所当行的。}\BibleRef{厄 5:3}因此，某人似乎说得对：\Quote{性交对任何人都无益，只有不伤害者才算正确。}因为即使是合法的性交，除了为了生育，也是危险的。至于违法的性交，圣经说：\Quote{淫妇如同母猪；凡受丈夫管束的，对使用者是死亡的堡垒。}他将娼妓的性情比作公山羊或公猪。而\Quote{死亡}他指的是，在受看守的娼妓身上所犯的奸淫。\Quote{房屋}和\Quote{城}则是指他们放纵无度的地方。此外，你们那里的诗歌也以某种方式谴责此事，写道：——\par

与你一起就是奸淫，与你一起就是可憎的性交，\newline{}肮脏、女里女气，最坏的城市，完全污秽。\par

相反地，他称赞这些贞洁的人：\par

那些被欲望控制的人，从未行过可耻的他人床笫之事，\newline{}也未曾带着憎恶行过可憎的淫行，\newline{}男人们从未如此。\par

因为许多人认为只有违反本性的事才是快乐，例如他们自己的这些罪恶。而那些比他们好的人，知道这些是罪，却被快乐所胜，黑暗遮蔽了他们邪恶的行径。因为凡奸淫般地使用婚姻、以娼妓方式行事的，就是犯奸淫，且不听导师的呼声：\Quote{那登上自己床榻，心里说：\InnerQuote{谁看见我？黑暗围绕着我，墙壁遮蔽了我，无人看见我的罪过。我何必害怕至高者会记念呢？}}\BibleRef{德 23:18}这样的人最为可怜，只畏惧人眼，以为能逃避天主的鉴察。因为圣经说：\Quote{他不知道至高者的眼睛比太阳光亮万倍，它们注视人的一切行径，并窥探隐秘之处。}因此导师又藉依撒意亚警告他们说：\Quote{祸哉，那些暗中策划，说：\InnerQuote{谁看见我们？}的人。}\BibleRef{依 29:15}因为人或许能逃避感官的光明，却无法逃避理智的光明。\Person{赫拉克利特}说：怎能躲避那永不沉没的呢？我们绝不可用黑暗遮蔽自己，因为光明住在我们内。经上说：\Quote{黑暗不能胜过光。}\BibleRef{若 1:5}连黑夜本身也被节制的理性照亮。圣经称善人的思想为\Quote{不眠的灯}；即使人企图掩饰自己的行为，也是公然犯罪。凡犯罪的，若是犯奸淫，直接伤害的与其说是邻人，不如说是自己，因为他犯了奸淫，使自己变得更糟、更被轻视。因为犯罪的人按犯罪的幅度变得更糟、更不被看重；被卑劣快乐所胜的人，现在放荡完全附着他。因此行淫的人完全死于天主，被圣言遗弃如同尸体被灵魂遗弃。因为圣洁之物按正义厌恶被玷污。但纯洁者常常可以接触纯洁者。我恳求你，不要脱衣服时也脱去廉耻；因为义人绝不可脱去节制。因为，看，这必朽的将穿上不朽；当那冲向放荡的贪欲，经过自律的锻炼，摆脱了腐败之爱，便将人交于永恒的贞洁。\Quote{因为在这世界，人娶嫁。}\BibleRef{玛 22:30}但当我们完成了肉身的作为，穿上了不朽，肉身本身也洁净了，我们就追求那合乎天使尺度的事。\par

因此，曾受教于异邦哲学的\Person{柏拉图}在\WorkTitle{费勒布篇}中神秘地称那些尽其所能毁灭和玷污内在之神（即圣言）的人为无神论者，他们因与恶习交往而如此。因此，献身于天主的人绝不可活于必死之中（\OriginalTerm{必死}{\GreekText{θνητῶς}}）。正如\Person{保禄}所说：\Quote{不合宜使基督的肢体成为娼妓的肢体；也不可使天主的殿成为卑劣情感的殿。}\BibleRef{格前 6:15}要记念那因行淫而被弃绝的两万四千人。但我已说过，那些行淫者的遭遇是纠正我们情欲的预像。再者，导师最明确地警告说：\Quote{不要随从你的情欲，戒避你的贪欲；\BibleRef{德 18:30}因为酒色能夺去智者的心；亲近娼妓的必更胆大。腐朽和虫子要继承他，他必被举为更可耻的公众例证。}又——因祂不厌倦行善——\Quote{那转离快乐目光的人，为生命加冕。}\par

因此，被情欲之事胜过、愚昧地贪恋淫欲、被与理性相悖的欲望驱使、渴望被玷污，都不是正当的。唯有那娶了妻子的人，如同农夫，被允许播种，当季节容许播种之时。对于其他方面的不节制，最好的药物是理性。饱足之匮乏也带来帮助，因为饱足激发的情欲冲向快乐。\par

% structure: p0134 source=h2 target=subsection reason=relative-heading-rank; base=h2

\subsection{第十一章——论服饰}

因此，我们也不该为自己预备昂贵的衣服，如同预备多样食物。主亲自将祂的训令分为关乎身体、灵魂以及第三类外物，劝告我们为身体的缘故预备外物；又藉着\OriginalTerm{灵魂}{\GreekText{ψυκή}}管理身体，训导灵魂说：\Quote{不要为你们的\OriginalTerm{生命}{\GreekText{ψυκῆ}}忧虑吃什么，也不要为身体忧虑穿什么；因为生命贵于食物，身体贵于衣服。}\BibleRef{路 12:22}祂又加上一个浅显的教导实例：\Quote{你们看乌鸦：它们不播种也不收割，没有仓也没有库，天主尚且养活它们。}\BibleRef{路 12:24}\Quote{你们不比飞鸟更贵重吗？}\BibleRef{路 12:24}祂同样，祂关于属于第三类外物的服饰也吩咐说：\Quote{你们看百合花，它们不纺也不织。但我告诉你们：连撒罗满在他极盛的荣华时所披戴的，也不如这些花中的一朵。}\BibleRef{路 12:27}而\Person{撒罗满}王曾因他的财富极其自夸。\par

请问，有什么比花朵更雅致、更绚丽的呢？有什么比百合或玫瑰更令人愉悦的呢？\BibleQuote{如果天主这样装饰田地里的野草，它今天还在，明天就投在炉中，何况你们呢，小信德的人哪！}\BibleRef{路 12:28} 这里，\Emph{\OriginalTerm{什么}{\GreekText{τί}}}这个词杜绝了食物上的花样。因为圣经上这样指示：\BibleQuote{不要为吃什么、喝什么而忧虑。}因为为这些事忧虑，表明贪婪和奢侈。现在，仅从本身来看，进食是需要的标志；饱足，如我们所说，是缺乏的标志。超出此范围，就是多余的标志。而多余的事，圣经宣布是属魔鬼的。接下来的话使意义清晰。因为说了\BibleQuote{不要寻求吃什么、喝什么}，祂又加上\BibleQuote{也不要心怀摇动（或：高傲）。}骄傲和奢侈使人从真理上犹豫不决（或：高举自己）；耽于逸乐、纵情多余之事，引人偏离真理。因此祂说得非常美妙：\BibleQuote{这一切都是世界上的外邦人所寻求的。}\BibleRef{玛 6:32} 外邦人就是放荡和愚昧的人。祂所指的这些东西是什么？奢侈、纵欲、丰盛烹调、美味佳肴、贪食。这些都是\Quote{什么？}至于简单的干粮和水饮的维持，作为必需品，祂说：\BibleQuote{你们的天父知道你们需要这一切。}简言之，如果我们天生就有寻求的倾向，那就不要通过导向奢侈来破坏寻求的能力，而应激发它去发现真理。因为祂说：\BibleQuote{寻求天主的国，这一切自会加给你们。}\par

如果主尚且除去对衣、食和一般多余之物的焦虑，视为不必要，那么关于爱好装饰、染羊毛、多色花样、对宝石的苛求、精巧的金工，更有甚者，假发和卷曲的发辫，以及染眼、拔毛、涂胭脂和铅白、染发，以及这些欺骗中所用的邪恶技艺，我们该当怎样说呢？我们岂不很有理由怀疑，上文关于青草所引的话，是针对那些不朴实而爱装饰的人说的？因为田野是世界，我们这些蒙受天主恩宠滋润的人就是青草；虽被割下，却又重新生长，正如在\WorkTitle{论复活}一书中将要更详细说明的。但干草象征粗俗的民众，依附于短暂的快乐，一时繁茂，爱装饰，爱赞美，除爱真理之外什么都爱，除了被火烧之外一无所用。主叙述说：\Quote{有一个富人，}主叙述说，\Quote{他身穿紫红袍，天天奢华宴乐。}这就是干草。\Quote{又有一个穷人，名叫拉匝禄，浑身生疮，被人放在富人家门口，渴望以富人桌上掉下的碎屑充饥。}这就是青草。那富人在\OriginalTerm{阴府}{Hades}中受罚，与火有分；而另一个却在圣父的怀中重新繁茂。我钦佩那古老的\Place{拉栖代梦}城，它只容许娼妓穿花衣、戴金饰，禁止良家妇女爱好装饰，只准妓女打扮自己。另一方面，\Place{雅典}的执政官们，他们崇尚精致的生活方式，忘记了自己的男子气概，穿着垂到脚踝的长袍，并梳着\OriginalTerm{克洛布卢斯}{crobulus}——一种发髻——饰以金蚱蜢的扣子，以此显示他们出自土地，在放荡的炫耀中自命不凡。这些执政官的竞相效尤也扩展到其他\Place{爱奥尼亚}人，荷马为了显示他们的柔弱，称他们为\Quote{长袍者}。因此，那些沉迷于美的形象（即爱慕虚荣）而非美本身的人，那些以美名重行偶像崇拜的人，应被远远逐出真理，因为他们凭借意见而非知识，梦想美的本质；如此，此世的生活对他们而言只是无知的沉睡；我们应当从中醒来，奔向真正美好而合宜的事物，渴望唯独把握它，将地上的饰物留给世界，在完全沉睡之前向它们告别。我认为，人需要衣服，无非是为了遮盖身体，防御过度的寒冷和酷热，以免气候的严酷伤害我们。如果衣服的目的是这个，那么就要注意不要给男人指定一种，给女人另一种。因为遮盖身体对男女都是共同的，如同吃喝一样。既然需要是共同的，我们就判断供应也应当相似。因为如同双方都需要遮盖之物，所以他们的遮盖也应当相似；尽管这种遮盖应取适合遮蔽女人眼睛的方式。因为如果女性因其软弱而想要更多，我们应责怪那不良训练的习惯，正是由于这个习惯，许多在恶习中长大的男人比女人更柔弱。但这不可让步。如果要做一些调整，可以允许她们使用更柔软的衣物，只要她们除去那些愚蠢薄的织物和奇特的编织花样；告别金绣、印度丝绸和精致的\OriginalTerm{蚕（博姆比克斯）}{Bombyces}（蚕丝），它起初是虫，然后变成毛虫，之后又经历第三次变形，称为蛹，从中产生长丝，如同蜘蛛的丝。因为这些多余而透明的材料是软弱心灵的证明，它们用薄纱遮盖身体的羞耻。因为奢华的衣服，若不能遮掩身体形状，便不再是遮盖。因为这样的衣服紧贴身体，更容易取得形状，如附着在肉体上，接受其形状，向旁观者勾勒出女人的形体，虽未看见身体本身。\par

染衣服也当摒弃。因为它既远离需要，也远离真理，此外还从中生出行为上的非难。因为颜色的使用并无益处，它们对防寒无用，也不比别的衣服更能遮盖，除了招致羞辱。颜色的悦目刺激贪婪的眼睛，使之盲目。对于内心洁白无瑕的人，最适宜穿白色素净的衣服。因此，达尼尔先知明明地说：\Quote{\BibleQuote{安置了宝座，上面坐着一位好似万古常存者，他的衣裳洁白如雪。}\BibleRef{达 7:9}}默示录也说，主自己穿着这样的袍子显现。又说：\Quote{\BibleQuote{我看见那些作证者的灵魂，在祭台下面，每人领受了一件白衣。}\BibleRef{默 6:9}}如果必要寻找别的颜色，真理的自然颜色已足够。但像花朵一样的衣服应留给酒神节的愚行和入教仪式，连同紫色和银器，如喜剧诗人所说：\par

\Quote{对悲剧演员有用，而非为生活。}\par

我们的生命应当不是一场游行。因此，\Place{萨迪斯}的染料、橄榄色、绿色、玫瑰色、猩红色以及成千上万其他染料，都为邪恶的纵欲而被费力发明出来。这样的衣服是为了观赏，而非为了遮盖。我们也当告别那些镀金的、紫色的、以野兽图案命名的奢侈品、浸染番红花色的衣袍、以及那些昂贵的、色彩斑斓的薄纱衣衫，连同其技艺本身。喜剧说：\Quote{这些女人做了什么明智的事？} \Quote{她们坐在那里，满身鲜花，穿着番红花色的衣服，涂脂抹粉？}\par

导师明确告诫说：\Quote{不要因你的衣服而夸耀，也不要因任何光荣而得意，因为这是不合法的。} \BibleRef{德 11:4}\par

因此，祂嘲笑那些身穿奢华衣服的人，在福音中说：\Quote{看，那些穿着华丽衣服、生活奢侈的人，是在王宫里。} \BibleRef{路 7:25}祂说的是易朽坏的王宫，那里有炫耀、爱慕名声、谄媚和欺骗。但那些在天庭上围绕万王之王的，是在圣神不朽的衣裳里，即肉体中，被圣化，从而穿上了不朽。\par

因此，正如未婚的女子只献身于天主，她的挂虑不分；而贞洁的已婚妇女在生活上分心于天主和她的丈夫；至于那些另有心意的则完全投身于婚姻，即情欲：同样，我认为贞洁的妻子，当她将自己献给丈夫时，真诚地事奉天主；但当她喜爱华丽服饰时，她就远离了天主和贞洁的婚姻，以她的丈夫换取世界，仿效那位阿戈斯的妓女，即\Person{厄里费勒}——\par

\Quote{她接受了比亲爱丈夫更珍贵的金子。}\par

因此，我钦佩那位刻奥斯岛的智者，他描绘了相称和合适的德行与恶习的形象，将前者，即德行，描绘成简单站立，穿白衣、纯洁，仅以谦逊为装饰（因为真正的妻子应当如此，被赋予谦逊）。而另一方，即恶习，相反，他介绍她穿着多余的衣裳，用不属于她自己的颜色粉饰；她的步态和神态被描绘成刻意取悦，勾勒出放荡女人的轮廓。\par

但跟随圣言的人不会沉溺于任何卑劣的快乐；因此，在衣着方面，应优先选择有用的。如果圣言藉达味论及主时歌唱说：\Quote{君王们的女儿因尊荣使你喜悦；王后站在你右边，身穿金线衣，佩戴金穗，} 祂所指的并非奢华的衣服；而是指示了那些获得怜悯的人——即教会——的不朽装饰，由信德编织而成；在其中，无邪的耶稣如金子般闪耀，蒙选者则是金穗。如果必须为妇女编织这样的衣物，让我们编织触感愉悦柔软的衣服，而不是像图画一样花哨以悦目。因为图画会随时间褪色，洗涤和浸泡在染料的药汁中会磨损羊毛，使布料脆弱；这不利于节约。为\OriginalTerm{佩普洛斯}{peploi}、\OriginalTerm{克西斯提德斯}{xystides}、\OriginalTerm{厄法普提德斯}{ephaptides}以及\Quote{斗篷}、束腰上衣和荷马所说的\Quote{遮羞之物}而焦虑，是愚蠢炫耀的极致。因为，说实话，当我看到如此多的财富被挥霍在遮蔽裸体上时，我感到羞愧。因为最初的人在乐园里用树枝和树叶为自己遮羞；而现在，既然羊是为我们创造的，让我们不要像羊一样愚蠢，而是受圣言训练，谴责衣着的奢华，说：\Quote{你是羊毛。} 尽管\Place{米利都}自夸，\Place{意大利}受赞美，许多人疯狂的羊毛被保护在皮下，我们仍不应迷恋它。\par

真福若翰，鄙视羊毛卷（因其散发出奢侈的意味），选择了\Quote{骆驼毛}，并穿在身上，使自己成为节俭和简朴生活的榜样。因为他也\Quote{吃蝗虫和野蜜}，\BibleRef{谷 1:6}那是甜美而属神的食物；正如他所做的，预备主的谦逊和贞洁的道路。因为他怎么可能身穿紫袍，而远离城市的繁华，隐退到旷野的孤独中，与天主宁静地生活，远离一切虚浮的追求——一切虚假的良善——一切卑鄙呢？\Person{厄里亚}使用羊皮外衣，并用毛织的腰带束紧。\BibleRef{列下 1:8}另一位先知\Person{依撒意亚}赤身露体，\BibleRef{依 20:2}常常身穿麻衣，即谦逊的服装。如果你提及\Person{耶肋米亚}，他只有\Quote{一条麻布腰带}。\BibleRef{耶 13:1}\par

因为正如营养良好的身体，在赤裸时更明显地显示出活力，同样，品格的美丽也显示出它的宽宏大量，当它不卷入炫耀的愚行时。但是拖着衣服，让它们垂到脚底，是十足的纨绔作风，妨碍行走的活动，衣服像扫帚一样扫过地面的尘土；因为甚至那些被阉割的舞者（他们把无声的无耻放荡转移到舞台上），也不鄙弃这种流荡到如此屈辱的服装；他们奇异的服装、流苏的装饰和精心制作的动作，显示出肮脏的柔弱。\par

如果有人提出主的衣裳垂到脚面，那件彩衣显示了智慧的花朵，即多样而不朽的圣经，主的圣言，闪耀着真理的光辉。圣神藉达味给主披上了另一件这样的袍子，当他这样歌唱时：\Quote{你以赞颂和优雅为衣，披上光明如外氅。}\par

因此，如同在制作衣服上，我们必须避免一切怪异，同样在使用上，也必须当心奢华。因为衣服既不宜短过膝盖，如传说中斯巴达少女那样；也不宜让女子的任何部分暴露在外。尽管你可以恰当地用以下话语来回答对你说\Quote{你的手臂很美；是的，但它不是给公众看的。你的大腿很美；但是回答是：只给我的丈夫。你的面容秀丽；是的，但只给娶我的人。}的人。但我不愿贞洁的妇女为那些藉赞美追逐谴责把柄的人提供如此赞美的机会；这不仅是由于禁止露踝，还因为命令要将头蒙住、脸遮住；因为美貌成为男人的陷阱是邪恶的。女子也不宜以紫色面纱引人注目。但愿能在服饰中废除紫色，以免观者的目光转向穿着者的面容！但妇女在制造其余所有衣物时，将一切都做成紫色，从而煽动情欲。事实上，那些痴迷于这些愚蠢而奢华的紫色——按诗句所说——\Quote{紫色（黑暗）死亡已攫住}的妇女。于是，因这紫色，\Place{提洛}、\Place{漆冬}以及\Place{斯巴达海}附近都备受渴求；其染匠、紫贝渔夫以及紫贝本身——因血产紫——都备受推崇。但狡黠的女子和柔弱的男子，将这些欺骗性的染料与精致的织物混合，将其疯狂欲望推向极端，不再从\Place{埃及}出口细麻布，而是从希伯来人和基里基雅人之地出口其他种类的亚麻布。我且不提\Place{阿莫尔戈斯}和\Place{比索斯}所产的亚麻布。奢华已超越命名。\par

在我看来，覆盖物应当显示被覆盖者比本身更好，如同圣像优于圣殿，灵魂优于肉身，肉身优于衣服。但现在却恰恰相反：这些女士的身体，若出售，值不了一千阿提卡德拉克马。而她们买一件连衣裙就花一万塔冷通，证明自己比布料更无用、更低贱。你究竟为何要追求罕见而昂贵的东西，胜过随手可得又廉价的东西？这是因为你不知道什么是真正美丽的、真正好的，反而从那些如同疯子般把黑当成白的愚妄之人那里急切地追求表象而非实在。\par

% structure: p0152 source=h2 target=subsection reason=relative-heading-rank; base=h2

\subsection{第十二章  论鞋履}

好炫耀的妇女在鞋上也同样行事，在这方面同样表现出极大的奢侈。那些钉有金饰的凉鞋实属卑劣；但人们却认为值得在鞋底上钉上螺旋排列的钉子。许多人还在鞋上刻上爱抚的拥抱，仿佛她们要通过行走将和谐的运动传达给大地，并将其精神的放荡印刻其上。因此，必须告别那些镀金镶宝的\Place{阿提卡}和\Place{西库翁}的半靴、\Place{波斯}和\Place{第勒尼安}的高筒靴，这些邪恶的发明；我们要以正确的目标为指引，正如我们的真理所习惯的，我们必须选择合乎自然的事物。\par

鞋的使用，部分是为了遮盖，部分是为了防备绊倒物体，并保护脚底免受崎岖山路之粗糙。\par

妇女应被允许穿白色鞋，除非在旅途中，则必须使用上油的鞋。在旅途中，她们需要钉鞋。此外，她们大多数时候应当穿鞋；因为赤足露脚不适宜；而且，女子是娇嫩之物，容易受伤。但对于男子，赤足是相当合适的，除非在服役时。因为\Quote{穿鞋近于被捆绑。}\par

赤足最适于锻炼，最有益于健康与舒适，除非有必需阻止。但若我们不在旅途，且不能忍受赤足，可使用拖鞋或白色鞋；雅典人称它们为\Quote{尘足}，我认为是因为它们使脚接近尘土。作证鞋履简朴的，有\Person{若翰}即可，他宣称\BibleQuote{他不配解主的鞋带}。因为那向希伯来人展示真哲学类型的人并未穿精巧的鞋。这还可能意味着什么，将在别处说明。\par

% structure: p0157 source=h2 target=subsection reason=relative-heading-rank; base=h2

\subsection{第十三章  反对过度喜好珠宝金饰}

过分赞赏深色或绿色石头、被大海抛到异国海岸的贝壳、地球的微粒，是幼稚之举。因为追求透明、颜色奇特、染色的玻璃，只是被炫目之物吸引的愚昧人的特征。孩子看到火，被其光亮吸引而冲向它，因无知而不懂触碰的危险。那些愚昧妇女佩戴在链条上和项链中的石头也是如此：紫晶、雷石、碧玉、黄玉和米利都的\par

\Quote{最珍贵的翡翠。}\par

而那备受珍视的珍珠已过度侵入妇女的闺房。它产自一种类似贻贝的牡蛎，大小约与大鱼的眼睛相当。这些可怜的女子不以为羞，竟在这种小牡蛎上费尽心思，而她们本可用那神圣的宝石——天主圣言——来装饰自己，圣经在某处称之为珍珠：纯洁而透明的\Person{耶稣}，在肉身中守望的眼睛——透明的圣言，由这圣言，那藉水重生的肉身成为宝贵的。因为水中那牡蛎将肉身完全包裹，从中生出珍珠。\par

我们也听说，天上的耶路撒冷是用圣石砌成的城墙；我们允许，天城的十二道门因被造成类似宝石，表明宗徒之声的超凡恩宠。因为颜色嵌在宝石上，这些颜色是珍贵的；而其余部分仍属泥土材质。圣人的城，属灵建造的，合宜地以这些象征性的石墙围起。因此，宝石的光辉意指圣神的无与伦比的光辉、存在的永生与圣洁。但这些女子不明白圣经的象征意义，拼命追求珠宝，提出令人惊愕的辩解：\Quote{为什么我不能使用天主所展示的东西？}以及\Quote{我拥有它，为什么不能享受它？}以及\Quote{那么这些东西是为谁造的，如果不是为我们？}这些是那些完全不知天主旨意之人的言语。因为首先的必需品，如水与空气，祂免费供应给所有人；而不必要的东西祂藏在地里和水中。因此蚂蚁挖掘，狮鹫守护黄金，海洋隐藏珍珠。但你们忙乎于你们不需要的东西。看啊，整个天空被照亮，你们却不寻求天主；但被隐藏的金子和珠宝，却被我们中间被判死刑的人挖出。\par

但你们也反对圣经，因圣经明明白白地呼喊：\Quote{你们先该寻求天主的国和它的义德，这一切自会加给你们。}\BibleRef{玛 6:33}但如果一切已赐予你们，一切都允许你们，并且\Quote{凡事都可行，但不全有益，}\BibleRef{格前 10:23}宗徒说。天主起初把自己的圣言赐给众人，使万物成为众人共享，从而将我们的族类带入共融。因此万物是共有的，而非富人独占不应有的份额。因此，那句话：\Quote{我拥有，且大量拥有：为何我不该享受？}既不适合人，也不适合社会。更可爱的是这句：\Quote{我拥有：为何我不该给予需要的人？}因为这样的人——遵行\Quote{你应当爱近人如你自己}命令的——是成全的。因为这是真正的奢侈——珍藏的财富。而浪费在愚昧情欲上的，应算作浪费，而非花费。因为天主确实赐予我们使用的自由，但仅限于必要；且祂规定使用应是共用的。一人奢侈而多人匮乏，是怪异的。向多人行善比奢华生活光荣多少！将钱花在人身上，比花在珠宝和黄金上智慧多少！获得有德的朋友比无生命的装饰有用多少！土地曾像施恩那样造福于人吗？因此，我们仍需消除这句辩辞：那么，如果所有人都选择更简朴之物，谁将拥有更奢华之物？我要说，如果人们公平无分别地使用它们，那么是人。但若所有人都行节制不可能，那么以必要之物为用，我们必须寻求最易获取的，向这些多余之物长久告别。\par

总之，她们必须因此彻底抛弃装饰，如同女孩的小玩意，完全拒绝装饰本身。因为她们应当在内里装饰，显出内在女子的美丽。因为美与丑只显现在灵魂中。故此，惟有有德者才是真正美善的。且有一条原则：惟有美的是善的。惟有德行透过美丽的身体显现，在肉身中开花，展现自制之可爱可亲，当品格如光柱在形体中闪烁时。因为每棵植物和动物的美在于其个体的卓越。人的卓越就是正义、节制、勇敢和虔敬。因此，美丽的人是正义、节制，总之是善的，而非富有的。但如今连士兵也希望用金饰装扮，他们没读过那诗句：\par

\Quote{带着幼稚的愚行他来到战场，\newline{}满载金银。}\par

但爱装饰，远不关心德行，却为自己攫取身体，当爱美变为虚浮炫耀时，必须彻底驱逐。因为将不适宜之物加于身体，仿佛它们是适宜的，产生说谎的习惯和虚假的习性；它不显示端庄、朴素和真正孩童般的，而是浮华、奢侈和柔弱的。但这些女子以金色遮盖，遮蔽了真正的美。她们不知道自己的过犯何等之大，在自己周围系上千万条贵重链条；正如他们说，在外邦人中，罪犯被金链捆绑。这些女子在我看来效仿这些富有的囚犯。因为金项链岂非项圈？他们称为\OriginalTerm{导管}{catheters}的颈饰岂非代替锁链？事实上在阿提卡方言中它们正以此名相称。这些围绕女子双足的不雅之物，\Person{菲勒蒙}在\WorkTitle{Synephebus}中称之为脚镣：\par

\Quote{显眼的衣裳和一种金脚镣。}\par

那么，女士们，你们所渴慕的这装饰，除了展示你们被捆绑，还能是什么？因为如果材料消除责备，那么忍受[你们的镣铐]是无所谓的。在我看来，那些自愿戴上镣铐的人，是在富裕的灾祸中自夸。\par

也许诗意的寓言说，这些锁链曾投在\Person{阿佛洛狄忒}通奸时，意指装饰品无非是通奸的标志。因为\Person{荷马}也称那些为金链。但新潮女子不羞于佩戴这最明显的恶者标记。因为正如蛇欺骗了厄娃，金饰也引诱其他女子走向恶行，以蛇的形状为饵，并制作七鳃鳗和蛇作为装饰。因此喜剧诗人\Person{尼科斯特拉托斯}说：\Quote{锁链、项圈、戒指、手镯、蛇形饰品、脚链、耳环。}\par

因此，阿里斯托芬在\WorkTitle{地母节妇女}中以最严厉的谴责，展示了女性装饰品的全部清单：\par

\Quote{发网、束发带、泡碱、钢刀；\newline{}浮石、头带、背带，\newline{}后纱、脂粉、项链，\newline{}眼影、软衣、发网，\newline{}腰带、披肩、紫边，\newline{}长袍、内衣、深坑、圆衫。}\par

但我尚未提及其中最主要的。那么是什么？\par

\Quote{耳坠、珠宝、耳环；\newline{}锦葵色的串状脚链；\newline{}扣环、搭扣、项饰，\newline{}脚镣、印章、链子、戒指、香粉，\newline{}圆凸饰、带子、仿阳具、玛瑙，\newline{}扇子、螺旋饰。}\par

我厌倦并烦恼于列举这么多装饰品；我不得不惊讶，那些背负如此重担的人竟没有被折磨至死。愚蠢的烦恼！愚蠢的炫耀狂热！她们像妓女般将财富浪费在可耻之物上；因热爱炫耀而损毁天主的恩赐，效仿那恶者的技艺。那将财富囤积在仓里，对自己说：\BibleQuote{你存有许多财物，足够多年之用；你休息罢，吃喝快乐罢！}的富翁，主在福音中清楚称他为\Quote{愚人}。\BibleQuote{今夜就要索回你的灵魂；你所预备的一切，将归谁呢？}\BibleRef{路 12:19}\par

画家阿佩莱斯见一名学生以大量金色颜料描绘海伦，便对他说：\Quote{小子，你无法将她画得美丽，便把她画得富有了。}\par

当今的淑女便是这样的海伦，并非真正美丽，而是装扮得富丽堂皇。圣神藉索福尼亚先知预言她们：\BibleQuote{他们的金银不能拯救他们在主忿怒的日子。}\BibleRef{索 1:18}\par

但对于那些受基督培养的女子，她们适宜不以金子装饰，乃以圣言装饰，因唯有藉着圣言，金子才得以显现。\par

古时的希伯来人若能抛弃或仅仅熔掉她们的女性饰物，便是有福的；但他们将金子铸成牛犊，并以偶像崇拜之，因此他们的技艺和企图均无益处。但他们极其明确地教导我们的女子远离饰物。那与金子行淫的贪欲成为偶像，且受火试验；奢靡被保留作为偶像而非现实。因此圣言藉先知斥责希伯来人：\BibleQuote{他们为巴耳制造金银之物}，即饰物。且极清晰地警告说：\BibleQuote{我要惩罚她为巴力诸日的日子，因为在那些日子她向他们献祭，佩戴耳环和项链。}\BibleRef{欧 2:8}接着祂附加了装饰的原因：\BibleQuote{她追随她的情人，却忘记了我——主说。}\BibleRef{欧 2:13}\par

因此，让我们将这些小玩意留给那邪恶的狡诈大师本身，不要参与这种娼妓式的装饰，也不要藉着似是而非的借口行偶像崇拜。因此，有福的伯多禄极卓越地说：\Quote{同样，愿妇女们以端庄的服饰，以廉耻和节制装饰自己，不以编发、黄金、珍珠或昂贵的衣服，而应以善工——这相称于自称敬畏天主的妇女。}因为祂有理由命令她们远离自我打扮。因为纵使她们美丽，本性已足够。不要让技艺与本性相争；即不要让虚伪与真理相争。若她们本性丑陋，则她们所涂抹之物证明她们缺乏美貌。因此，事奉基督的女子宜于简朴。因为简朴确实为圣洁提供条件，它通过将多余之物削减为平等，并利用现有之物来提供从多余中寻求的享乐。简朴，正如其名所示，并不显眼，不在任何事上膨胀或自大，而是全然平坦、温柔、平等、无过度，因而足够。足够是达致其适当终点而无过度或不足的状态。其母是正义，其保姆是\Quote{独立}；这是满足于必要之物，并自行提供有助于幸福生活的状态。\par

因此，愿你们手上的果实有神圣的秩序、慷慨的施与和节俭的行为。\BibleQuote{因为怜悯穷人的，就是借贷给主。}\BibleRef{箴 19:17}\BibleQuote{手勤的，必得丰裕。}\BibleRef{箴 10:4}祂称那些藐视财富、慷慨施予之人为\Quote{勤勉的}。愿你们脚上显出行善的活跃和朝向正义的旅程。谦逊和贞洁是项圈和项链；这些是天主所铸造的链条。\BibleQuote{找到智慧的人是有福的，获得明哲的人是有福的，}圣神藉撒罗满说：\BibleQuote{因为赚得她比赚得金银更好；她比珍珠更宝贵。}\BibleRef{箴 3:13}因为她是真正的装饰。\par

不要违反自然穿耳孔，以挂上耳环和耳坠。因为强迫自然违反其意愿是不对的。对于耳朵，没有比真正的教导更好的装饰，它自然地进入听觉通道。被圣言膏抹的眼睛和为领悟而穿孔的耳朵，使人成为神圣和圣洁之事的听众和默观者，圣言真正展示那真美，\BibleQuote{是眼睛未曾看见，耳朵未曾听见的。}\BibleRef{格前 2:9}\par

\SourceRecord{Translated by William Wilson.；Ante-Nicene Fathers；Vol. 2.；Edited by Alexander Roberts, James Donaldson, and A. Cleveland Coxe.；Buffalo, NY: Christian Literature Publishing Co.,；1885.；<http://www.newadvent.org/fathers/02092.htm>.}

% structure: p0001 source=h1 target=section reason=page-title

\section{\WorkTitle{导师（第三卷）}}

% structure: p0002 source=h2 target=subsection reason=relative-heading-rank; base=h2

\subsection{第一章　论真美}

由此可见，认识自己是所有功课中最大的功课。因为人若认识自己，就会认识天主；认识天主，就会肖似天主，不是靠穿戴黄金或长袍，而是靠行善并尽可能减少所需。\par

唯独天主一无所缺，并且当我们以智慧的装饰发光时，祂最喜悦；祂也喜悦那穿着贞洁——这身体的神圣长袍——的人。既然灵魂分为三部分：理智，即被称为推理官能，是内里的人，是这可见之人的统治者。而那一部分，在另一方面，由天主引导。但易怒的部分，是兽性的，近乎疯狂。而欲望，是第三部分，其形态多变，甚于\OriginalTerm{普罗透斯}{Proteus}，这变化多端的海神，时而变成这种形状，时而变成那种；它引诱人犯奸淫、纵欲、勾引。\par

\Quote{起初他是一只狮子，胡须浓密。}\par

当他仍保留着装饰时，下巴上的毛发表明他是个男人。\par

\Quote{但随后是一条蛇、一头豹、或一头大母猪。}\par

对装饰的热爱已堕落为放荡。人不再像一头强壮的野兽，\par

\Quote{但他化作了湿润的水，和一棵高枝的树。}\par

激情爆发，快感泛滥；美貌凋谢，比秋日来临前，当淫欲的风暴吹拂它时，叶子落到地上更快，被毁灭所枯萎。因为情欲制造并伪装一切，想要欺骗，以便隐藏真人。但那与圣言同住的人，不改变自己，不打扮自己：他拥有来自圣言的形态；他肖似天主；他是美丽的；他不装饰自己：他的美是真美，因为它是天主；那人成为天主，因为天主如此愿意。因此，\OriginalTerm{赫拉克利特}{Heraclitus}说得对：\Quote{人是神，神是人。}因为圣言本身就是显明的奥秘：天主在人内，人成为天主。而中保实现父的旨意；因为中保是圣言，是两者所共有的——天主之子，人的救主；祂的仆人，我们的导师。既然肉体是奴隶，如保禄所见证的，人怎么能有理地像拉皮条者那样装饰这婢女呢？因为属于肉体的，具有奴隶的形态。保禄论及主时说：\BibleQuote{祂空虛了自己，取了奴仆的形体，}\BibleRef{斐 2:7}称外在的人为奴仆，在主成为奴仆、穿上肉体之前。但慈悲的天主亲自解放了肉体，使它脱离朽坏和痛苦致命的束缚，赐予它不朽，将不朽——这永恒的神圣装饰——披在肉体上。\par

人还有另一种美——爱。按照宗徒的话：\BibleQuote{爱是含忍的，是慈祥的；爱不嫉妒，不夸张，不自大。}\BibleRef{格前 13:4}因为打扮自己——带着多余和无益的外观——正是自大。因此他接着说：\BibleQuote{不作无礼的事：}因为不属于自己、违反自然的形象是无礼的；而人工的不是自己的，正如清楚说明的：经上说：\BibleQuote{不求自己的利益。}因为真理称属于它自己的为自己的；但爱打扮寻求不属于自己的，离开了天主、圣言和爱。\par

主自身在面貌上是不俊美的，圣神通过依撒意亚作证：\BibleQuote{我们看见祂，祂没有威仪，也没有美貌；但祂的容貌是卑微的，逊于常人。}然而有谁比主更可敬佩呢？但祂所展示的，不是肉眼可见的肉体之美，而是灵魂和肉体的真美；前者表现为仁爱，后者——即肉体——表现为不死。\par

% structure: p0013 source=h2 target=subsection reason=relative-heading-rank; base=h2

\subsection{第二章　反对装饰身体}

因此，所要装饰的不是外在之人的面貌，而是灵魂，以善之饰；我们也可说，肉体以节制的装饰。但那些美化外表的妇女，不知不觉在内心全然空虚，正如埃及人的装饰品：他们精心建造带有柱廊和前庭的庙宇，以及附带的圣林和圣田；大厅环绕着许多柱子；墙壁因外来的石头而闪耀，艺术绘画也不匮乏；庙宇因黄金、白银、琥珀而闪耀，并因来自印度和埃塞俄比亚的彩色宝石而闪烁；神龛用金绣帷幔遮掩。\par

但如果你进入围场的内部，急于看到更好的东西，寻找那居于庙中的神像，如果有任何献祭的祭司在那里，庄重地观看，并用埃及语唱一首赞歌，稍微掀开帷幔以显示神祇，他将会让你对那崇拜的对象捧腹大笑。因为所寻找的神祇——你所冲向的那位——不会在里面找到，而是一只猫、一条鳄鱼、或一条本地蛇，或某种这样的野兽，不配这庙宇，却完全配得上一个洞穴、一个坑或污泥。埃及人的神表现为一只在紫色床榻上打滚的野兽。\par

同样，那些佩戴黄金、忙于卷曲头发、涂脂抹粉、画眼染发、并操练奢侈的其他有害技艺、装饰肉体这层覆盖物的妇女——实际上模仿了埃及人，以便吸引她们迷恋的情人。\par

但若有人揭开圣殿的帷幔——我指的是头饰、染料、衣服、金饰、颜料、化妆品，即由这些编织而成的帷幔——为要寻找内在的真正美丽，我知道，他必会感到厌恶。因为他将找不到天主住在里面的肖像，如理当所见；取而代之的是，一个淫乱者和淫妇占据了灵魂的殿堂。于是那真正的野兽将被揭露——一只涂着白漆的猿猴。而那欺骗的蛇，借着虚荣吞噬人的理智部分，以灵魂为洞，充满致命毒液；且注入其欺骗的毒汁，这作恶的龙把女人变成了妓女。因为爱炫耀不是淑女所为，而是妓女的行为。这样的女人很少在乎与丈夫待在家中；反倒松开丈夫的钱包，将供给挥霍在情欲上，为有许多见证人见证她们看似美丽的容貌；终日修饰打扮，与买来的奴隶消磨时光。因此她们像有害的酱汁一样调味肉；白日关在房间里梳妆，以免被人发现正在装扮。到晚上，这伪造的美爬出来，到烛光下，如同从洞里爬出；因为醉酒和昏暗的灯光助长了她们的装扮。染黄头发的女人，喜剧诗人\Person{墨南德}驱逐出家门：——\par

\Quote{现在离开这房子，\newline{}因为没有贞洁的女人应当染黄头发}\par

我还会补充，没有染色脸颊，没有绘画眼睛。可怜的她们无意中通过引入赝品摧毁了自己的美丽。在黎明时分，她们折磨、拉扯、涂抹自己以某些混合物，冷却皮肤，用毒药划破肉身，并用精心准备的洗剂，从而摧残自己的美丽。因此她们因使用化妆品而显得蜡黄，且易患病，其肉身因毒药覆盖而处于融化状态。因此她们羞辱造物主，仿佛他所赐的美貌毫无价值。可想而知，她们在家务上变得懒惰，坐在那里像画出来的东西供人观看，而不是为家庭经济所造。因此在那位喜剧诗人笔下，理智的妇女说：\Quote{我们妇女能做什么聪明或杰出的事情？我们坐在这里染着黄头发，冒犯了淑女的品格；导致家庭的倾覆，婚姻的毁灭，以及来自子女的指责？}同样地，喜剧诗人\Person{安提法奈斯}在\WorkTitle{玛尔塔卡}中，以适用于所有女人的话语嘲笑女人的卖弄风情，并针对她们涂抹化妆品的行为而编成嘲笑语，说道：——\par

\Quote{她来，\newline{}她回去，她走近，她回去。\newline{}她已来，她在此，她洗身，她前进，\newline{}她擦肥皂，她梳头，她外出，被揉擦，\newline{}她洗身，照镜子，穿衣，\newline{}涂油，装饰，涂抹；\newline{}若有何不妥，便（因烦恼）窒息。}\par

我再说，不是一次，而是他们应当死三次，那些使用鳄鱼粪便、以腐烂体液之沫涂身、用煤烟染眉、用铅白擦颊的人。\par

这些女人，其风尚甚至令外教诗人厌恶，又怎能不被真理所弃绝？因此另一位喜剧诗人\Person{亚里克西斯}斥责她们。我将引用他的话，这些话语以过分的言辞羞辱了她们厚颜的固执。因为他并未离谱太远。我因羞愧而无法助佑被喜剧如此嘲笑的女子。\par

然后她毁掉了自己的丈夫。\par

\Quote{首先，与获利和邻人的损害相比，\newline{}其余一切在他们眼中都是多余的。}\newline{}\Quote{其中一人矮小？她把软木缝进鞋底。\newline{}一人高大？她穿薄底鞋，\newline{}出门时低头靠肩：\newline{}这消减了她的身高。一人无腰？\newline{}她缝上东西，使旁观者\newline{}赞叹她背后的优美身段。一人肚子突出？\newline{}添加补物使之直挺，如同我们在喜剧诗人中见到的奶妈，\newline{}她仿若用这些杆子，将前方的肚腩拉回。\newline{}一人黄眉毛？她用煤烟染色。\newline{}若是黑色？她涂铅白。\newline{}一人皮肤很白？她擦胭脂。\newline{}一人身体某部位美丽？她裸露之。\newline{}她有美齿？她必然笑，\newline{}使在场者看见她多美的嘴；\newline{}但若她无心发笑，便整日待在室内，\newline{}双唇间夹着一根纤细的桃金娘枝，\newline{}如同厨师们常有手边的山羊头去卖，\newline{}如此她必须一直保持双唇分开，不论是否情愿。}\par

我将这些引自喜剧诗人的段落摆在你们面前，因为圣言最热切地要拯救我们。稍后我将以圣经来加强它们。因为不被察觉的人惯于戒除罪恶，由于斥责的羞耻。正如涂灰的手和涂油的眼从外貌显出病人可疑的迹象，同样地，化妆品和染料表明灵魂已深染疾患。\par

神圣的导师命令我们不要接近他人的河流，借比喻的言词\Quote{他人的河流}意指\Quote{他人的妻子}；那放荡的女子向众人流淌，并因纵欲将自己交付于所有人与卖淫般的享乐。祂说：\Quote{你要戒饮他人之水，}和\Quote{不要饮他人井中之水}，劝诫我们避开\Quote{情欲}的溪流，使我们得享长寿，且年岁得以加增\BibleRef{箴 9:11}，既因不追逐属他人的快乐，也因转移我们的倾向。\par

贪肴馔与嗜酒，虽为大恶，尚不及爱奢华之甚。\Quote{丰盛的餐桌与重复的杯盏}足以满足贪欲。但爱黄金、紫袍与珠宝者，地上地下的黄金不足，第勒尼安海、印度与埃塞俄比亚的货物、乃至产金的帕克托洛斯河皆不足；即便人成为迈达斯，亦不满足，仍觉贫穷，渴望更多财富。这样的人甘愿与黄金同死。\par

若普路托斯是瞎眼的，那些为他疯狂并与他同感的女人岂非也瞎眼？她们既无限制其贪欲，便走向无耻。剧场、游行、众多观众、在神殿闲逛、街头徘徊，为使众人均可见到，皆为其所需。因为那些夸耀容貌而非内心的人\BibleRef{得前 2:17}为悦人而打扮。正如烙印标明奴隶，华彩标明淫妇。\BibleQuote{纵然你穿上朱红衣，佩戴金饰，用锑石画眼，你的美丽也是枉然}\BibleRef{耶 4:30}，圣言借耶肋米亚说。岂不荒谬：马、鸟及其他动物从草地与草场跳跃欢腾，以鬃毛、天然颜色与斑斓羽毛为自己的装饰而欣喜；而女人，仿佛低于动物，竟以为自己无可爱之处，需要外来的、购买的、涂抹的美？\par

头饰与各种头饰、精巧编发、无限的梳妆方式、昂贵的镜子——她们在其中整饰着装，追逐那些像傻孩子般为体型疯狂之人——乃是无羞耻女人的特征。若有人称她们为妓女，并无错误，因她们将脸化为面具。但我们，圣言命令\BibleQuote{不要看那些看得见的，而要看那些看不见的；因为看得见的是暂时的，看不见的才是永恒的}\BibleRef{格后 4:18}。\par

而超越荒谬界限的是，她们为这虚假的形态发明了镜子，仿佛那是某种卓越之作或杰作。这骗局反倒需要遮上薄纱。按希腊传说，美丽的纳西索斯看见自己的形象并非幸事。若梅瑟曾命令人不可用艺术制造天主的形象，这些女人通过自己的反射模仿自己的容貌以伪造面容，又怎能正确？同样，当先知撒慕尔被派去给叶瑟的一个儿子傅油为王，看见长子英俊高大，便欢喜地取出傅油膏，主对他说：\BibleQuote{不要看他的外貌和身高，我拒绝了他。因为人看眼目，主看内心}\BibleRef{撒上 16:7}。\par

他傅油的不是容貌美之人，而是灵魂美之人。若主视身体的天然美低于灵魂，他对虚假之美又会如何？他彻底拒绝一切谎言。\Quote{因为我们凭信德行走，而非凭眼见}。主借亚巴郎非常清楚地教导：跟随天主的人必须轻视祖国、亲戚、财产及所有财富，使他成为外方人。因此，他也称他为朋友，因他藐视在家中所拥有的产业。因为他出身良好，极为富有，并以自己的三百一十八名仆人征服了掳走罗特的那四位君王。\par

我们发现惟有艾斯德尔正当装饰。新娘为了王的丈夫神秘地装扮；但她的美却成了即将被屠杀之民的赎价。那种装饰使女人成为娼妓，男人成为柔弱的奸夫，悲剧诗人可为证，他如此说：——\par

\Quote{那位审判女神者，\newline{}按阿尔戈斯人的神话，从弗里吉亚\newline{}来到拉刻戴蒙，穿着繁花似锦的衣袍，\newline{}闪耀着黄金与蛮族的奢华，\newline{}爱上并带走了他所爱之人——\newline{}海伦，去往伊达的山谷，发现\newline{}墨涅拉俄斯不在家。}\par

哦，奸淫之美！蛮族的华饰与柔弱的奢华推翻了希腊；拉刻戴蒙的贞洁被衣服、奢华与优雅之美腐化；蛮族的炫耀使宙斯之女成为娼妓。\par

他们没有导师约束其贪欲，也没有人说：\Quote{不可奸淫}；或\Quote{不要贪恋}；或\Quote{不要因贪恋而走向奸淫}；或更甚：\Quote{不要因渴望装饰而影响你的情欲}。\par

他们结局如何？遭受何等祸患？他们不愿约束自己的任性！两大洲因放纵的享乐而震动，一切因一个蛮族少年陷入混乱。整个希腊出海；海洋承载着大洲之重；漫长的战争爆发，激烈的战斗爆发，平原堆满尸体；蛮族放肆攻击舰队；邪恶得胜，那诗意的宙斯之眼注视色雷斯人：——\par

\Quote{蛮族的平原饮尽高贵的鲜血，\newline{}河流的溪流被尸体堵塞。}\par

捶胸哀悼，悲伤荒芜大地；所有山脚、众泉之源的伊达山顶、特洛伊人的城市、阿开亚人的船只，尽皆震动。\par

荷马啊，我们往哪里逃，往哪里站？指给我们一块不受震动的地面吧！——\par

\Quote{不要碰缰绳，无经验的少年，\newline{}也不要坐上车座，既然尚未学会驾车。}\par

天堂喜爱两位御者，唯独他们驾驭火战车。因为心神被快乐带走；而无瑕的理性原则，当未受圣言\OriginalTerm{圣言}{Word}训导时，便滑入放纵，并因其过犯而跌倒，得到应得的报应。天使便是例子，他们为了会消逝的美而弃绝了天主的美，如此便从天堂坠落到地上。\par

舍根人也因污辱了那位圣洁的童贞女而受惩罚遭倾覆。坟墓就是他们的惩罚，他们耻辱的纪念碑引向救恩。\par

% structure: p0043 source=h2 target=subsection reason=relative-heading-rank; base=h2

\subsection{第三章 反对那些妆扮自己的男人}

奢侈已发展到如此地步，不仅女性因此虚荣追求而迷失，连男性也被感染此病。因为他们未脱离爱装饰的心，故不健康；倾向纵欲，从而变得女性化，以不绅士和娼妓般的方式剪发，穿着精美透明的衣服，嚼乳香，散发香水味。看到他们，能说什么呢？就像凭额头判断人的人，会猜测他们是奸夫和女性化的男人，沉溺于两种淫行，恨头发、少头发、厌恶男子气概，像女人一样装饰发型。\Quote{这些反复无常的可怜虫，为不敬的狂妄行为而活，干着鲁莽邪恶的事，}西彼拉说。为了伺候他们，城镇到处充斥着用沥青膏脱毛、刮胡、拔毛的人，服务于这些女性化的家伙。到处开设店铺；精于此等娼妓式淫乱的人公开赚大钱，从那些涂抹自身的人身上，以及让操此业者用各种方式拔除毛发的人身上，他们在旁观者或靠近的人面前，甚至在自己面前，都不感到羞耻，尽管他们是男人。这就是那些沉溺于卑劣情欲的人，全身因沥青膏的猛拉而变得光滑。这样的厚颜无耻简直无可超越。如果他们无所不为，我也将无所不说。狄奥革涅斯在被出卖时，像老师一样责骂其中一个堕落之徒，非常刚毅地说：\Quote{来吧，年轻人，为你自己买一个男人，}用一句含糊的话责罚他的娼妓行为。但身为男人却剃须并柔化自己，多么卑劣！染发、涂抹灰白鬓发、染黄它们，这些是被遗弃的女性化者的行为；女性化的梳妆应被禁止。因为他们以为，像蛇一样，他们通过涂色和更新自己来脱去头上的老年。但即使他们巧妙地医治头发，也逃不过皱纹，也无法通过欺骗时间而躲过死亡。因为并不恐怖，并不恐怖显得年老，当你不能闭眼不看自己确实年老的事实时。\par

人越急迫地走向终点，就越真正可敬，因为天主是他唯一的长者，祂是永恒的古者，万有中最为年长。先知称祂为\BibleQuote{\OriginalTerm{万古常存者}{Ancient of days}；祂的头发如同纯净的羊毛。} 主说：\BibleQuote{除了祂以外，没有人能使头发变白或变黑。} \BibleRef{玛 5:36} 那么，这些不敬虔的人怎能与天主竞争，甚至激烈反对祂，改变祂所变白的头发呢？\BibleQuote{长者的冠冕是丰富的经验，}圣经说。\BibleRef{德 25:6} 他们面容上的白发是丰富经验的绽放。但这些人不尊敬年迈的尊严，不尊敬布满白发的头。不可能，他绝不可能显露出真实的头，倘若他的头是虚假的。\BibleQuote{但你们从基督所学的却不是这样。如果你们真听过祂，并按照在耶稣内的真理在祂内受过教导，就该脱去你们照从前生活的旧人（不是白发老人，而是那依欺骗的情欲而腐败的人），并应更新（不是靠染发和装饰），而是更新你们的心思，穿上新人，就是按照天主的肖像在正义和真正的圣德中受造的新人。} \BibleRef{弗 4:20}\par

但是，对于一个男人来说，为了追求美观而用剃刀梳理头发、刮脸、拔除胡须并使其光滑，这是多么女人气啊！事实上，若非你看到他们赤裸的身体，你准会以为他们是女人。因为他们虽然不被允许佩戴金饰，却出于娇柔的欲望，用金叶缠绕自己的鞋带和衣缘；或者用同样的金属制成某种球形饰物，系在脚踝上，挂在脖子上。这是孱弱男人的发明，他们被拖到女眷的居所，是两栖的、淫乱的禽兽。因为这是一种娼妓式的、不虔敬的陷阱。天主希望女人光滑，只以自然生长的头发为乐，如同马以其鬃毛为乐；但祂却装饰男人，如同狮子，赐予胡须，并赋予他多毛的胸膛作为刚毅的标记——这是力量和统治的记号。同样，天主为保护母鸡而战斗的公鸡装饰了鸡冠，如同头盔；天主对这些毛发如此重视，以致祂命令它们与明辨一同出现在男人身上，并以可敬的外表为乐，用白发尊荣面容的庄重。但是，智慧以及与智慧一同苍白的明辨判断，随着时间成熟，并藉长久经验的活力给老年增添力量，产生白发——可敬智慧的卓越花朵，赢得信任。因此，这男人的标记——胡须，藉此他被视为男人——比厄娃更古老，是更高本性的标记。天主认为他应当在这方面优越，并将毛发散布于男人的全身。祂在形成女人厄娃时，从男人的肋骨中取走了其中所有的光滑和柔软；厄娃在身体上易接受，是他的生育伴侣、家务帮手，而他（因为他已放弃了所有光滑）仍是一个男人，并显示出自己是男人。行动被分配给他，如同受苦被分配给她；因为多毛比光滑更干燥、更温暖。因此，雄性比雌性毛发更多、体温更高；完整的动物比阉割的动物更完全；完美的比不完美的更完全。因此，亵渎男性气概的象征——多毛，是不敬的。但是，使自己光滑的装饰（因为我受到\OriginalTerm{圣言}{Word}的警告），如果是为了吸引男人，那就是娇柔之人的行为；如果是为了吸引女人，那就是奸夫的行为；两者都必须尽可能远离我们的团体。\BibleQuote{就是你们的头发，也都被数过了。}主说；\BibleRef{玛 10:30}下巴上的头发也被数过了，全身的头发也被数过了。因此，绝不可拔除，违反天主的安排，祂已按自己的旨意将它们数算在内。\BibleQuote{你们岂不知道，}宗徒说，\BibleQuote{耶稣基督就在你们内吗？} \BibleRef{格后 13:5}如果我们知道祂居住在我们内，我不知道我们怎敢羞辱祂。但是，使用沥青拔除毛发（我甚至羞于提及与此过程相关的无耻），以及在后仰和弯腰的动作中，因跨步和以可耻姿势后仰而对本性端庄施加的暴力，然而做这些事的人却不以自己为耻，反而在青年中间、在考验男子气概的体育馆中无耻地行事；追随这种违反本性的做法，岂非放纵到极致？因为那些在公开场合从事这些事的人，在家中对任何人几乎不会保持端庄。他们在公开场合的无耻，证明了他们在私下里毫无节制的放纵。因为那在白日否认自己男子气概的人，在夜间将显明地证明自己是一个女人。\BibleQuote{不可有以色列女儿中的娼妓；不可有以色列儿子中的淫乱者。} \BibleRef{申 23:17}\par

但有人说，沥青有益处。不，我说它玷污名声。没有一个持正确见解的人愿意显得像个淫乱者，除非他是那恶习的牺牲品，并图谋玷污自己形体的美。没有人会自愿选择这样做。\BibleQuote{因为天主预知那些蒙召的人，按祂的旨意被预定肖似祂儿子的肖像，}为了这儿子，按蒙福宗徒的话，祂立定\BibleQuote{他为众多弟兄中的长子，} \BibleRef{罗 8:28}那些凌辱与主相似的身体的人，难道不是不敬神吗？\par

要成为美男子的男人，必须装饰人身上最美的东西——他的心智，他每天都应使之更显优雅；他应当拔除的，不是毛发，而是情欲。我怜悯那些被贩卖者控制的男孩，他们被妆饰以受辱。但他们并非自愿受辱，而是被命令，这些可怜的人为了卑鄙的利益而被装饰。但是，那些自愿实行这些事的人——如果被迫，他们若是有男子气概，宁愿死也不去做——是多么可憎呢？\par

但是，生活由于邪恶的放纵已达到了这种放荡的程度，淫荡已弥漫于各城，成了法律。在她们旁边，妓女站在妓院中，为淫乐的欢愉而出卖自己的肉体；男孩们被教导否认自己的性别，扮演女人的角色。\par

奢侈已混乱一切，使人蒙羞。奢靡的考究追求一切，尝试一切，强迫一切，胁迫本性。男人扮演女人的角色，女人扮演男人的角色，违反本性；女人同时是妻子和丈夫：没有一条通路不对淫欲开放；他们混杂的淫荡成为公开的制度，奢侈已在家中扎根。可悲的景象！可怖的行为！这就是你们社交放荡所展示的胜利：这些行径的证据就是娼妓。哀哉如此邪恶！此外，这些可怜人不知道交合的偶然性产生了多少悲剧。因为父亲，不记得自己遗弃的子女，常常在不知情下，与自甘堕落的儿子和做娼妓的女儿交合；而情欲的放纵表明他们是生育他们的人。你们明智的法律允许这些事：人们可以合法地犯罪；他们称这可憎的享乐为无关紧要之事。那些违反本性行奸淫的人自以为没有犯奸淫。报应的正义追赶他们胆大妄为的行径，给他们招致不可避免的灾祸，他们用小钱买来死亡。这些可悲的商贩航行运载淫乱之货，如同酒或油；另一些人，更为可悲，买卖享乐如同买卖面包和酱料，不顾梅瑟的话：\BibleQuote{不可让你的女儿卖淫，以致这地成为娼妓之地，地上充满邪恶。}\BibleRef{肋 19:29}\par

这自古以来就有预言，结果是众所周知的：整个大地如今已充满淫乱和邪恶。我钦佩罗马的古代立法者：他们憎恶行为的柔靡；将身体用作女性用途、违反自然律者，他们依法之正义，判定应受最严厉的惩罚。\par

因为拔除胡须——人自然高贵的装饰——是不合法的。\par

\Quote{刚长初须的青年：因为这样，青春最为优雅。}\par

不久之后，他被傅油，喜悦于那\Quote{降在}预表性的\Quote{油膏}上的胡须，这油膏曾使亚郎受尊荣。\par

那受过正确培育的人，和平在他内安营，他也当与自己的头发保持和平。\par

那么，当女人看到男人做出如此令人发指的事时，那些有强烈纵欲倾向的女人会做出什么呢？我们不应称这些人为男人，而应称他们为淫荡的\OriginalTerm{败类}{\GreekText{βατάλοι}}和柔弱的女人（\OriginalTerm{柔弱}{\GreekText{γύνιδες}}），他们的声音微弱，衣服在质感和颜色上都像女人。从他们的外表、衣服、鞋、身形、步态、发式、眼神，这些家伙明显显露出他们的本质。\BibleQuote{因为从一个人的外表可以认出他，}圣经说，\BibleQuote{从遇见一个人可以认出这人：人的衣服、脚步、牙齿的笑容，都讲述他的事。}\BibleRef{德 19:29}\par

这些人多半拔去其余毛发，只整理头上的头发，几乎像女人一样用发带束住。狮子以它们浓密的毛发为荣，在战斗中毛发是它们的武器；野猪甚至因鬃毛而显得威严；猎人看见它们竖起毛发时就害怕。\par

\Quote{毛茸茸的羊被它们的羊毛压着。}\par

慈悲的圣父使它们的羊毛丰盛，为你所用，人啊，教导你剪取它们的毛料。在民族中，克尔特人和西徐亚人留长发，但不去装饰自己。蛮族的浓密毛发有某种可怖之处；其\OriginalTerm{红褐色}{\GreekText{ξανθόν}}预示战争，色调有点近于血。这两个蛮族都憎恶奢侈。日耳曼人将以莱茵河为证，西徐亚人以马车为证。有时西徐亚人甚至轻看马车：其大小对蛮族而言显得奢华；西徐亚人放弃其奢侈安逸，过着简朴的生活。因为一所比马车更足够且更少累赘的房屋，他取来坐骑，骑上它，被载到想去的地方。当饥饿时，他向马求食；马献出自己的血管，以所有的一切——自己的血——供应主人。对游牧者而言，马同时是运输工具和食物；阿拉伯人的好战青年（这是另一种游牧者）骑在骆驼上。他们坐在繁殖期的骆驼上；骆驼一边吃草一边奔跑，同时驮着主人，并带着房屋。如果蛮族缺少饮水，他们就挤骆驼奶；食物耗尽后，他们甚至不吝惜血，据说像疯狂的狼一样。这些蛮族比野兽温和，受伤时不怀恨，而是勇敢地驰骋沙漠，同时驮着并喂养主人。\par

那些以血为食物的野兽，该灭亡！因为人的身体不过是血所造的肉，触摸血是不合法的。因为人的血已分享了圣言；它因圣神而分受恩宠；如果有人伤害他，必不能逃脱惩罚。人虽身体赤裸，仍可向主祈求。但我赞赏蛮族的简朴：热爱无累赘的生活，蛮族放弃了奢侈。主召叫我们成为这样的人——赤裸于华丽，赤裸于虚荣，脱离我们的罪，只背着生命的木，只追求救恩。\par

% structure: p0061 source=h2 target=subsection reason=relative-heading-rank; base=h2

\subsection{第四章 我们应与谁交往}

但我确实无意中背离了次序，现在必须回归正题，并指责拥有大量家仆的陋习。因为人们避免亲手劳作和自行服务，转而雇用仆人，购买一大群精致的厨师、摆放餐桌的人、以及熟练地将肉切割成块的人。仆从队伍分成许多类别：有些为贪饕服务，有切肉者、调味者、甜点师、蜂蜜蛋糕师、蛋奶冻师；有些则忙于过多的衣物；有些像格里芬一样守卫金银；有些保管银器、擦拭杯盏、准备节日餐桌所需的物品；有些刷马；一大群司酒者竭力服务；还有成群的俊美少年，如同牛群，从他们身上榨取美貌。男女梳妆助手则为女眷服务——有的拿镜子，有的管头饰，有的拿梳子。许多是阉人；这些拉皮条者因被认为无力纵欲而毫无嫌疑地服务于那些想自由享乐的人。但真正的阉人不是无力，而是不愿纵欲。圣言通过先知撒慕尔向悖逆的犹太人作证，当百姓要求立王时，他应许的不是仁爱之主，而是威胁赐予他们一个任性放纵的暴君，祂说：\Quote{祂必取你们的女儿作香料师、厨子、烤面包的，}\BibleRef{撒上 8:13}以战争之法统治，不渴望和平治理。还有许多克耳特人，他们肩扛女轿。但羊毛工、纺线工、织布工、女工及家务劳动却无处可寻。\par

但那些欺骗妇女的人整日与她们厮混，讲些愚蠢的爱情故事，用虚假的言行耗尽身心。经上记着：\Quote{不可随众行恶，也不可附和多数；}\BibleRef{出 23:2}因为智慧显于少数，混乱显于多数。但这些妇女购买轿夫并非出于端庄、不愿被人看见（若出于此意而遮盖自己倒是可称赞的），而是出于奢华，被家仆扛在肩上，并想炫耀一番。\par

于是她们掀开帘子，锐利地环视所有投向她们的目光，显出她们的作风；而且常常从里面探出身来，因危险的躁动而辱没了这表面的端庄。经上记着：\Quote{不可在城里的街上东张西望，也不可在僻静之处闲逛。}\BibleRef{德 9:7}因为那确实是一个僻静之处，即使其中满是放荡之徒，却没有智者在场。\par

这些女人被人抬着穿梭于庙宇，日日献祭和占卜，与算命者、乞讨的司祭和声名狼藉的老妇厮混；她们在酒杯旁低语老妇的耳语，从占卜者那里学习符咒和咒语，以破坏婚姻的纽带。她们或包养某些男人，或被某些男人包养，或由占卜者许诺给她们。她们不知是在欺骗自己，并将自己当作享乐的器皿献给想行淫乱的人；将贞洁换成最污秽的暴行，竟把最可耻的毁灭当作大买卖。许多侍奉这种娼妓式放荡的人，从四面八方渗透进来。因为放荡者像猪一样冲入污秽，冲向船底凹陷处。因此圣经极力劝诫：\Quote{不可让每个人都进入你的家，因为奸诈之人的罗网众多。}\BibleRef{德 11:29}又在另一处说：\Quote{当以义人为你的宾客，并在敬畏上主中存留你的夸耀。}\BibleRef{德 9:16}远离淫乱。\Quote{因为你们应当清楚知道，}宗徒说，\Quote{凡是淫乱的、不洁的、贪心的（就是拜偶像的），在基督和天主的国里都得不到产业。}\BibleRef{弗 5:5}\par

但这些妇女喜爱与柔弱男子交往。成群的\OriginalTerm{淫荡者}{\GreekText{κιναίδες}}涌来，口无遮拦，身体污秽，言语污秽；他们足以充当淫秽的差役、奸淫的帮凶，咯咯笑着、窃窃私语，无耻地用鼻子发出淫荡和奸淫的声音以挑起情欲，尽力以淫词和姿态取悦人，激起笑声——奸淫的前兆。有时，因某种挑衅而激怒时，无论是这些奸淫者，还是那些追随可憎的流氓走向毁灭的人，都会像青蛙一样从鼻子里发出声音，仿佛愤怒居于他们的鼻孔中。但比这些人更高雅一些的，则饲养印度鸟和米底孔雀，与尖头动物一起斜躺；与\Person{萨提尔}玩耍，以怪物为乐。他们听到\Person{忒耳西忒斯}时便发笑；这些妇女购买高价\Person{忒耳西忒斯}们，不以丈夫为荣，反以那些地上的负担为荣，忽视那贞洁的寡妇，她远高于一只\OriginalTerm{梅利塔小狗}{Melitæan pup}，并对正义的老人侧目而视，在我看来，他比用钱买来的怪物更可爱。尽管饲养鹦鹉和鹬鸟，她们却不收留孤儿；她们遗弃自己家中所生的孩子，却收养鸟雏，偏好无理性的受造物胜过有理性的；尽管她们本应照顾那些性格节制的老人，在我看来，老人比猿猴更美丽，能说出比夜莺更好的话语；并应将这句箴言摆在他们面前：\BibleQuote{怜悯穷人的，就是借给主；}\BibleRef{箴 19:17}以及这句：\BibleQuote{你们对我这些最小兄弟中的一个所做的，就是对我做的。}\BibleRef{玛 25:40}但这些妇女却相反，偏爱无知胜过智慧，将财富转化为石头，即珍珠和印度翡翠。她们将财富挥霍在褪色的染料和买来的奴隶身上，如同填塞的禽鸟刮取生命中的粪土。\BibleQuote{贫穷，}经上说，\BibleQuote{使人谦卑。}\BibleRef{箴 10:4}这里的贫穷是指那种吝啬，使富人变穷，无物可施。\par

% structure: p0067 source=h2 target=subsection reason=relative-heading-rank; base=h2

\subsection{第五章 在浴室中的行为}

他们的浴室又是怎样的呢？精工建造的房子，紧凑、便携、透明，覆盖着细麻布。镀金的椅子，银的椅子，以及成千上万的金银器皿，有饮用的、有进食的、有沐浴的，都随身携带。此外，还有炭火盆；因为他们已到达如此放纵的地步，竟在沐浴时用餐并喝醉。还有银制的器具用来炫耀，他们在浴室中显摆它们，从而或许因过度骄傲展示财富，但主要是那任性的无知，借此他们给被妇女征服的柔弱男子打上烙印；至少证明他们自己若无大量器皿就无法会面、无法出汗，而穷苦妇女没有炫耀同样享受沐浴。因此，财富的污秽有丰富的责备遮盖。以此如同诱饵，他们钩住那些对黄金闪光目瞪口呆的可怜虫。因为他们这样炫惑那些喜爱炫耀的人，巧妙地试图赢得情人的钦佩，而后者稍后便会赤身裸体地侮辱她们。她们几乎不会在丈夫面前脱衣，假装合理的谦逊；但其他任何愿意的人，都可以在家中看到她们锁在浴室里赤身裸体。因为在那里，她们在旁观者面前脱衣不以为耻，仿佛在出卖自己的肉体。但\Person{赫西俄德}劝告道：\par

\Quote{不要在女子的浴池里洗皮肤。}\par

浴室对男女开放；他们在那里为放纵的欢乐而脱衣（因为从观看起，人就开始爱上），好像谦逊已在浴室中被洗去。那些尚未完全失去谦逊的人则禁止陌生人入内；但与自己仆人一起沐浴，在奴隶面前赤身裸体，并让他们擦身；给蹲伏的仆人以放纵情欲的自由，允许无惧的触摸。因为那些在沐浴时被引到赤身女主人面前的人，刻意脱衣以胆敢纵欲，因这邪恶的习俗抛却了恐惧。古代的运动员，因羞于展示裸体男子，穿着短裤参加比赛以保持谦逊；但这些妇女，脱去外衣的同时也脱去了谦逊，她们希望显得美丽，却违背自己的意愿仅仅显为邪恶。因为通过身体本身，情欲的放荡清楚地闪耀；如同患水肿的人，皮肤遮盖下的水。两种疾病都从外观而知。因此，男性应为女性提供真理的高尚榜样，应当羞于在她们面前脱衣，并提防这些危险的景象；\BibleQuote{因为好奇观看的，}经上说，\BibleQuote{已经犯了罪。}\BibleRef{玛 5:28}因此，在家中，他们应当以谦逊对待父母和家仆；在路上，对待遇见的人；在浴室，对待妇女；独自一人时，对待自己；并在各处对待那无处不在的\OriginalTerm{圣言}{Word}，\BibleQuote{没有祂，什么也没有。}\BibleRef{若 1:3}因为唯有这样，人才能保持不跌倒，如果他视天主常与他同在。\par

% structure: p0071 source=h2 target=subsection reason=relative-heading-rank; base=h2

\subsection{第六章 唯独基督徒富有}

财富当理性地享有，仁爱地施与，不污秽也不浮夸；对美的爱也不应变为自爱和炫耀；免得有人对我们说：\Quote{他的马、地、仆人或金子值十五塔冷通，而这人本身只值三个铜钱。}\par

那么，直接从妇女身上取下装饰，从主人那里取走仆人，你就会发现主人在步态、容貌或声音上与买来的奴隶毫无差别，他们与自己的奴隶如此相似。但他们不同之处在于比奴隶更虚弱，且教养更病态。\par

这句最好的格言，应当反复重复：\Quote{善人，既然节制而正义，}在天堂积蓄财富。那变卖了自己世俗的财物，分给穷人的人，找到了不朽的宝藏，\Quote{那里既没有蛾子，也没有盗贼。}他真是有福的，\Quote{尽管他卑微、软弱、默默无闻；}他真地富有着最大的财富。\Quote{因此，一个人即使比\Person{Cinyras}和\Person{Midas}更富有，却是邪恶的，}并且傲慢如那个穿着紫袍和细麻衣奢华度日、藐视拉匝禄的人，\Quote{他是可怜的，生活在忧苦中，}而且将不得存活。财富在我看来就像一条蛇，会缠绕手并咬人；除非人懂得如何抓住尾巴尖端而无危险。财富，无论在经验丰富或生疏的手中扭动，都善于附着和咬人；除非人轻视它，巧妙地使用它，以圣言的魅力压碎这受造物，而自己安然无恙。\par

但是，理所当然，唯独那拥有最贵重之物的人，才是真正富有的，尽管不被认可。并非宝石、黄金、衣服、或美貌具有高价，而是德行；德行就是导师吩咐付诸实践的圣言。这圣言摒弃奢侈，却召叫自助为仆人，并赞扬节俭，它是节制的产物。他说：\Quote{领受教诲，而非白银；宁要知识，而非精炼的黄金；因为智慧胜于宝石，任何贵重之物都无法与她的价值相比。}\BibleRef{箴 8:10}又说：\Quote{宁可得着我，胜过黄金、宝石和白银；因为我的出产比精选的银子更好。}\BibleRef{箴 8:19}\par

但若我们必须区分，就应承认，那拥有许多财物、满载黄金如同肮脏的钱袋的人是富有的；但唯有义人是有恩宠的，因为恩宠即秩序，在管理和分配上持守适当且得体的尺度。\Quote{因为有人撒种，收割更多，}\BibleRef{箴 11:24}论到这样的人经上写着：\Quote{他施舍，他赒济穷人；他的仁义存到永远。}因此，不是那拥有并保留的人富有，而是那施舍的人富有；使人幸福的不是拥有，而是施舍；圣神的果实是慷慨。所以，财富在于灵魂。那么，就应承认，美善只属于善人；而基督徒是善人。愚人或放纵者既不能领悟什么是善，也不能拥有它。因此，美善唯独基督徒拥有。没有什么比这些美善更富足的；所以只有他们是富有的。因为正义是真正的财富；圣言比一切财宝更有价值，并非来自牲畜和田地，而是由天主赐予——是不能被夺去的财富。灵魂本身是其财宝。这是对拥有者最好的财产，使成为真正有福的人。因为一个人的愿望是不渴求任何不在我们能力之内的事物，并通过祈求从天主那里获得他虔敬渴望的东西，难道他不是拥有许多，甚至一切，以天主为他的永恒财宝吗？经上说：\BibleQuote{凡求的，就必得到；凡敲门的，就必给他开门。}\BibleRef{玛 7:7}若天主不拒绝任何事物，则一切都属于虔诚的人。\par

% structure: p0077 source=h2 target=subsection reason=relative-heading-rank; base=h2

\subsection{第七章 节俭是基督徒的良好预备}

花在享乐上的美食对人是危险的海难；因为这种多数人耽于逸乐、卑贱的生活与对真正的美和优雅的快乐格格不入。人本性上是直立而庄严的存有，作为唯一者的受造物，追求美善。但是爬行于肚腹的生活缺乏尊严，是可耻、可憎、可笑的。对天主性而言，纵欲是最不相宜的；因为这就是人像麻雀一样贪食，像猪和山羊一样淫荡。因为将快乐视为善，是对美善全然无知的标志。对财富的贪爱使人偏离正确的生活方式，并诱使他不再以羞耻之事为耻；只要像野兽一样，能吃到各种东西、同样地喝，并用各种方式满足他的淫欲，他就无所谓。因此他很少继承天主的国。那么，这些精致的菜肴准备来做什么呢？不过是填满一个肚子吗？贪食的污秽由我们肚子排泄食物残渣的阴沟所证明。他们聚集如此多的侍酒者，而只用一杯就能满足自己是为何故？那些衣箱又是为何？金饰呢？那些东西是为偷衣贼、恶棍和贪婪的眼睛预备的。经上说：\Quote{但愿你施舍和信德不离开你。}\BibleRef{箴 3:5}\par

试看提斯贝人厄里亚，在他身上我们有节俭的美好榜样：当他坐在荆棘下时，天使给他带来食物。\BibleQuote{是一个大麦饼和一罐水。}\BibleRef{列上 19:4}主将这作为最好的给他。那么，我们在通往真理的旅程中，必须轻装上阵。主说：\BibleQuote{不要带钱囊，不要带口袋，不要带鞋；}\BibleRef{路 10:4}也就是说，不要拥有财富，那不过是收在钱囊中的；不要装满自己的仓库，好像把收成装在口袋里，而要分施给有需要的人。不要为马匹和仆役操心，当富人旅行时，它们就像负担一样，被寓言式地称为鞋。\par

因此，我们必须抛弃大量的器皿、银质和金质的饮杯，以及众多的仆人，因为我们已经从导师那里接受了公平而庄严的侍从：自助和简朴。我们必须与圣言合宜地行走；如果有了妻子和儿女，房屋就不是重担，因为学会了随同心地纯正的旅客迁徙。爱丈夫的妻子应当装备得与丈夫相似。那些怀着贞洁的庄重带着节俭的人，拥有前往天堂的美好旅费。正如脚是鞋的尺度，身体也是每个人所有物的尺度。但那些多余的、人们称为装饰和富人家具的东西，是负担，而非身体的装饰。用力攀爬天堂的人，必须携带仁爱的美丽手杖，并通过与困苦者分享而达到真正的安息。因为圣经肯定说：\Quote{灵魂真正的财富是人的赎价，}\BibleRef{箴 13:8}也就是说，如果他是富有的，他将通过施舍而得救。因为正如涌出的井，被抽干后又恢复到原来的水位，同样，施舍，作为爱的仁慈泉源，通过将它的饮料分给口渴的人，又增加和溢满，就像乳汁习惯于流入被吮吸或挤出的乳房中。因为拥有全能的天主、即圣言的人，什么也不缺，从不缺乏所需。因为圣言是无所需要的财富，是一切丰裕的原因。若有人说他常看到义人缺乏食物，这是罕见的，除非那里没有另一个义人。无论如何，让他读下面的经文：\Quote{因为义人不是单靠饼生活，而是靠主的话，}\BibleRef{申 8:3}他乃是真正的饼，天上的饼。因此，善人永远不会陷入困境，只要他保持对天主的宣认完整。因为属于他的是向万有之父祈求并领受他所需要的一切；并享有属于他自己的，如果他持守圣子。这也属于他：感到不缺乏。\par

这位训练我们的圣言赐予我们真正的财富。对于那些通过祂而拥有无所缺乏的特权的人来说，致富并非可羡慕之事。拥有这财富的人将继承天主的国。\par

% structure: p0082 source=h2 target=subsection reason=relative-heading-rank; base=h2

\subsection{第八章 比喻和榜样是正确教导最重要的一部分}

如果你们中有人完全避免奢侈，他将通过节俭的教养，训练自己忍受非自愿的劳作，常常利用自愿的苦楚作为迫害的训练；这样当他遇到强迫的劳作、恐惧和忧伤时，就不会缺乏忍受的实践。\par

因此我们在世上没有家乡，以便我们轻视地上的财物。而节俭在最高程度上是富有的，等于不间断的支出，用于必需之物，且达必要程度。因为\OriginalTerm{费用}{\GreekText{τέλε}}有支出的意思。\par

关于丈夫如何与妻子共同生活，以及自助、家务和仆役的使用；此外，关于婚姻的时间和适合妻子的事，我们在关于婚姻的讲话中已经论述了。现在只保留与纪律相关的内容作描述，因为我们描绘基督徒的生活。大部分已经说过，并以纪律规则的形式确立。剩下的我们将补充；因为榜样对于决定救恩并非小事。\par

悲剧说：\par

\Quote{奥德修斯的配偶没有被杀死\newline{}由忒勒马科斯；因为她没有在丈夫之外再嫁一个丈夫，\newline{}但在家中，婚姻之床保持未受玷污。}\par

他指责可耻的通奸，同时通过对丈夫的爱情展示了贞洁的美好形象。\par

拉刻代蒙人强迫他们的仆人黑劳士（黑劳士是他们仆人的名称）喝醉，在他们节制的面前展示他们的醉态，以作疗治和纠正。\par

因此，观察他们的不雅行为，为了自己不落入类似受责难的行为中，他们训练自己，将醉汉的指责转化为保持自身无过错的益处。\par

因为有些人通过受教而得救；而另一些人自学，或追求或寻求德行。\par

\Quote{他确实是最优秀的人，他亲自洞察万物。}\par

这就是亚伯拉罕，他寻求天主。\par

\Quote{再者，听从善言劝告的人是好的。}\par

这就是那些听从圣言的门徒。因此前者被称为\Quote{朋友}，后者\Quote{宗徒}；前者殷勤寻求，后者宣讲同一天主。两者都是子民，两者都有听众，一个通过寻求得益，另一个通过找到得救。\par

\Quote{但是无论谁，既不能自己洞察，也不听从他人，\newline{}放在心上——他就是无用的人。}\par

另一个子民就是外邦人——无用的；这是不跟随基督的子民。然而，导师是热爱人的，以多种方式帮助，有时劝勉，有时责备。其他人犯了罪，祂向我们展示他们的卑劣，并展示随之而来的惩罚，以爱劝诫我们，通过展示之前遭受苦难者的榜样来劝阻我们远离罪恶。通过这些榜样，祂非常明确地制止了那些已有恶意的人，并阻止了那些敢做类似行为的人；又使另一些人奠定忍耐的基础；阻止某些人作恶；又通过观察类似情况来治愈另一些人，将他们引向更好的事。\par

因为谁在路上跟随一个人，然后看到前者落入坑中，难道不会警惕避免同样危险，小心不跟随他滑倒？哪个运动员，在得知通往荣耀的道路，并看到在他之前争战的人获得了奖赏，难道不为冠冕努力，效法前辈？\par

这样神圣智慧的图像为数众多，但我只提一个例子，并用几句话加以阐述。对索多玛人的审判，对作恶者是惩罚，对听闻者是教训。索多玛人由于过度奢华而陷入不洁，无耻地行奸淫，对男童燃起疯狂的爱恋；全视的圣言——行不虔敬之事的人无法逃脱其注意——将目光投向了他们。那不眠的人性守护者并没有沉默地观察他们的放荡；而是劝阻我们效法他们，训练我们趋向祂自己的节制，并降罚于一些罪人，免得情欲因未受惩罚而挣脱一切恐惧的约束，命令焚烧索多玛，将少许明智的火倾泻于放荡之上；免得情欲因缺乏惩罚而为那些冲入纵欲之门的人敞开大门。因此，索多玛人公正的惩罚对世人成了为人类妥善计算的救恩的图像。因为那些没有犯下与受罚者同样罪行的人，绝不会遭受同样的惩罚。通过避免犯罪，我们就避免了受苦。\Quote{因为我要你们知道，}犹达说，\Quote{天主既然一次拯救了祂的子民离开埃及地，后来就把那些不信的毁灭了；至于那些不守本位、离开自己住处的天使，祂用永远的锁链把他们拘留在幽暗里，等候大日的审判。}接着他又以最具教导性的方式列举了受审判者的象征：\Quote{祸哉他们！因为他们走了该隐的道路，又为巴兰的错谬而贪婪地奔跑，并在科辣黑的悖逆中灭亡。}因为那些不能获得义子特权的人，恐惧会防止他们变得傲慢。惩罚和威胁正是为此目的：因惧怕刑罚而避免犯罪。我可以向你讲述因炫耀而受罚、因虚荣而受罚的事例，不仅是因放荡；并引证对那些因财富而心怀恶念之人的谴责，在这些谴责中，圣言通过恐惧来制止恶行。但为免在论文中过于冗长，我将提出导师的以下训诫，好使你谨防祂的威胁。\par

% structure: p0100 source=h2 target=subsection reason=relative-heading-rank; base=h2

\subsection{第九章　为何要沐浴}

那么，沐浴有四种理由（因为我在演讲中从那里岔开），我们常去洗浴：为清洁、为取暖、为健康，或最后为享乐。为享乐而沐浴应当省略。因为不知羞耻的享乐必须连根拔除；妇女应为清洁和健康而沐浴，男人则只为健康。为取暖而沐浴是多余的，因为可以用其他方式恢复因寒冷而冻僵的身体。频繁沐浴也会削弱体力，松弛身体能量，常引起虚弱和昏厥。因为在某种程度上，身体像树木一样，不仅通过嘴饮水，而且在沐浴时通过所谓的毛孔全身饮水。证明这一点：人们常常口渴，随后进入水中就解了渴。因此，除非沐浴有某种用途，否则我们不应沉迷于它。古人称浴场为\Quote{熨平人的地方}，因为它们比应该的更早地使人的皮肤起皱，并且像煮东西一样，迫使它们过早老化。肌肉像铁一样，被热软化，因此我们需要冷，如同淬火和开刃。也不应总是沐浴；如果一个人稍微疲惫，或者相反，饱胀过饱，应当禁止沐浴，同时考虑身体的年龄和季节。因为并非对所有人都总是有益，正如那些精通此道的人所承认的。但在所有场合我们都称之为生活助手的适当比例，对我们来说就足够了。因为我们不应如此使用沐浴以致需要助手，也不应像常去市场那样在一天中频繁多次沐浴。但让几个人把水浇在我们身上，是对邻人的冒犯，出于对奢侈的喜爱，并且是那些不明白浴场对全体沐浴者同样开放的人所做的。\par

但最必要的是用洁净的圣言洗涤灵魂（有时也洗涤身体，因为污垢聚集并附着其上，有时也为了缓解疲劳）。\Quote{祸哉，你们经师和法利塞人，假善人！}主说，\Quote{你们好像粉刷的坟墓。外面看来美丽，里面却满是死人的骨骸和各种不洁。}\BibleRef{玛 23:27}祂又对同一群人说话：\Quote{祸哉你们！因为你们洗净杯盘的外面，里面却满是污秽。你们先洗净杯子的里面，好使外面也清洁。}\BibleRef{玛 23:25}那么，最好的沐浴是擦去灵魂的污染，是属神的。预言明确论及此事：\Quote{上主将洗净以色列子女的污秽，并从他们中间清除血迹}\BibleRef{依 4:4}——即罪的血迹和先知的被杀。而净化的方式，圣言接着说道：\Quote{用审判的神和焚烧的神。}\BibleRef{依 4:4}肉身的沐浴，即身体的沐浴，仅用水完成，如同在乡下没有浴场时常常那样。\par

% structure: p0103 source=h2 target=subsection reason=relative-heading-rank; base=h2

\subsection{第十章　适于善生的操练}

对男孩来说，即使浴室就在附近，体育场也足够了。甚至对男人来说，偏爱体操运动远胜于浴室，或许也不算坏事，因为这在某些方面有助于年轻人的健康，并能产生努力——一种旨在追求不仅健康体魄而且勇敢灵魂的竞争。当这样做时，如果不使人脱离更有益的事务，那就是愉快且有益的。妇女也不应被剥夺身体运动。但不应鼓励她们参与摔跤或跑步，而是让她们在纺纱、织布以及必要时监督烹饪中锻炼自己。她们还应亲手从储藏室取所需之物。对她们来说，从事磨坊工作并不丢脸。作为管家和助手的妻子，忙于烹饪以使食物合乎丈夫口味，这也不是责备。如果她整理床榻、在丈夫口渴时递上饮品、尽可能整洁地将食物摆上桌子，并以此使自己得到有益健康的锻炼，那么教导者将赞许这样的妇女，她\Quote{伸手辅助有用之事，手扶纺锤，张开手施舍，伸腕给乞丐。}\par

效法\Person{撒辣}的女子不以那最崇高的服务——帮助旅客——为耻。因为\Person{亚巴郎}对她说：\Quote{快，拿三斗细面，揉和，作些饼。}\BibleRef{创 18:6}\Quote{\Person{辣黑耳}，\Person{拉班}的女儿，来了，}经上说，\Quote{带着她父亲的羊。}\BibleRef{创 29:9}这还不够，为教导谦逊，又加上\Quote{因为她牧放她父亲的羊。}圣经中提供了无数此类关于节俭、自助以及运动的例子。就男人而言，让一些人脱衣摔跤；让一些人在阳光下玩小球，尤其是他们称为\OriginalTerm{腓尼达球}{Pheninda}的游戏。对于另一些去乡间散步或下城的人，步行便是足够的运动。如果他们使用锄头，这种农业劳动中的节俭之举并非不雅。\par

我几乎忘了说，著名的\Place{米利都}王\Person{庇塔库斯}曾从事推磨的艰苦劳动。一个男人自己打水、劈自己要用的木柴，这是可敬的。\Person{雅各伯}牧放了托付给他的\Person{拉班}的羊群，他有一根\Quote{安息香木的杖}作为王权标志，其木料旨在改变和改善本性。大声朗读对许多人来说也是一种运动。但我们所允许的这些体育竞赛，不应为了虚荣而进行，而应为了流出男子汉的汗水。我们也不应以狡诈和炫耀的方式摔跤，而要以站立搏斗的方式，通过解脱颈、手和侧身的纠缠。因为这种带有优雅力量的努力更得体、更有男子气概，是为了有益健康的缘故。至于那些以体操中不高尚姿势为业的人，应被辞退。我们必须始终以适度为目标。因为正如劳动应优于食物一样，过度劳动既非常糟糕、非常疲惫，又容易使我们生病。因此，我们既不应完全闲散，也不应完全疲惫。因为我们在饮食方面所规定的，同样适用于一切场合。我们的生活方式不应使我们习惯于享乐和放纵，也不应走向另一个极端，而应介于两者之间，即和谐、节制、远离奢侈和吝啬这两种邪恶。此外，正如我们之前提到的，料理自己的需要是一种没有骄傲的运动——例如，自己穿鞋、自己洗脚、以及当涂油后自己揉搓。为那揉搓过你的人回馈同样的服务，是互惠正义的运动；陪伴生病的朋友、帮助弱者、供给贫乏者，都是适当的运动。\Quote{\Person{亚巴郎}，}经上说，\Quote{在树下为三人备好晚餐，并在他们吃的时候伺候他们。}\BibleRef{创 18:8}同样，捕鱼也是如此，如\Person{伯多禄}的例子，如果我们有暇离开必要的圣言训导。但主赐给门徒更美好的享受，当祂教导他\Quote{捕人}如同捕水中的鱼时。\par

% structure: p0107 source=h2 target=subsection reason=relative-heading-rank; base=h2

\subsection{第十一章 基督徒生活概观}

因此，佩戴金饰和使用较柔软的衣服并非完全禁止。但必须抑制非理性的冲动，免得它们通过过度松弛将我们带走，驱使我们走向享乐。因为纵欲，一旦冲至过饱，便易于尥蹶子、甩鬃毛，并甩掉御者——教导者；祂从远处拉紧缰绳，引领并驱策人这匹马——即灵魂的非理性部分——它狂野地追求快乐、邪恶的欲望、宝石、黄金、服饰的多样以及其他奢侈，走向救恩。\par

最重要的是，我们要牢记那神圣的话语：\Quote{在外邦人中要品行端正，使他们虽然毁谤你们是作恶的，但因看见你们的好行为，便在鉴察的日子归荣耀于天主。}\BibleRef{伯前 2:12}\par

% structure: p0110 source=h3 target=subsubsection reason=relative-heading-rank; base=h2

\subsubsection{衣服}

因此，教导者允许我们使用简单的衣服，并且如我们先前所说，用白色。这样，我们不是使自己适应于花哨的艺术，而是适应于自然本身，并摒除一切欺骗和歪曲真理的东西，从而拥抱真理的一致性和单纯。\par

\Person{索福克勒斯}责备一个青年时说：\par

\Quote{穿着女人的衣服。}\par

因为，如同士兵、水手和统治者的情况一样，节制之人的合适衣着是朴素、合宜且清洁的。因此，在律法中，\Person{梅瑟}关于癞病所立的律法排斥那些具有多种颜色和斑点、如同蛇鳞般杂乱的衣物。因此，他愿意人不再用各种颜色装饰得花里胡哨，而是从头到脚一身洁白，成为洁净；如此，由身体过渡，我们也能摆脱人多变而易变的私欲，热爱真理那不变、明确而单纯的颜色。在这方面效法\Person{梅瑟}的——尤其是\Person{柏拉图}——也赞同那种不超过贞洁妇女手工的织物。白色很能配得上庄重。他在别处也说：\Quote{不可施用染料或纺织，除非是为战争装饰。}\par

因此，白色适合于和平与光明之人。正如与原因紧密相连的标记，其出现指示或更确切地说证明了结果的存在；如烟是火的标记，好的肤色和规律的脉搏是健康的标记；同样，这类衣着也显示出我们习惯的品性。节制是纯洁而简单的；因为纯洁是一种确保行为纯洁、不掺杂卑劣事物的习惯。简单是一种去除多余之物的习惯。\par

厚实的衣物，尤其是未经缩绒的，能保护体内的热量；并非衣物本身有热，而是它阻挡体内散发出的热量，不让其逸出。任何落在其上的热量，它都会吸收并保留，并因温热而反过来温暖身体。因此这主要在冬天穿着。\par

它（节制）也是知足的。知足是一种摒除多余之物的习惯，并且为了不致匮乏，能接纳按照\OriginalTerm{圣言}{Word}为健康与有福生活所需的一切。\par

让妇女穿着朴素合宜的衣裳，比男子所穿的更柔软，但绝非过分暴露或完全沉溺于奢华。衣物应适合年龄、身份、体型、本性和职业。因为神圣的宗徒最美妙地劝告我们：\BibleQuote{穿上耶稣基督，不要为肉欲作准备。}\BibleRef{罗 13:14}\par

% structure: p0119 source=h3 target=subsubsection reason=relative-heading-rank; base=h2

\subsubsection{耳环}

圣言禁止我们因刺穿耳垂而违背自然。为何不也刺穿鼻子呢？——如此，所说的就应验了：\BibleQuote{耳朵上的金环在猪鼻子上，美貌而无见识的妇女也是如此。}\BibleRef{箴 11:22}总之，若有人以为自己因金子而美丽，他就低于金子；低于金子的人就不是金子的主人。而承认自己不如吕底亚矿石装饰，这是多么荒谬！因此，正如金子被母猪的污秽所玷污——母猪用鼻子搅动泥泞——同样，那些因奢侈而过分放荡、被财富抬高的妇女，以纵欲的污点玷污了真正的美丽。\par

% structure: p0121 source=h3 target=subsubsection reason=relative-heading-rank; base=h2

\subsubsection{戒指}

因此，圣言允许她们戴一枚金戒指。但这并非为了装饰，而是为了在她们行使家务管理职责时，封印那些值得在家中安全保存的物品。\par

因为如果所有人都受过良好教育，那么就不需要印章了，如果仆人和主人同样诚实的话。但由于缺乏教育使人易于不诚实，我们便需要印章。\par

但在某些情况下，这种严格可以放宽。因为有时必须为那些不幸未遇上贞洁丈夫的妇女开脱，她们打扮自己是为了取悦丈夫。但只应让她们以赢得丈夫的赞赏为目的。我不愿她们致力于个人展示，而要以贞洁的爱——一种强大而合法的魅力——吸引她们的丈夫。然而，既然丈夫们希望妻子精神不愉快，那么妻子们若愿贞洁，就应致力于逐渐缓解丈夫们非理性的冲动和情欲。她们应被温柔地引向简朴，逐渐使他们习惯于节制。因为端庄不是通过强加沉重负担而产生的，而是通过去除多余。妇女的奢侈品应被禁止，因为这些是快翼之物，产生不稳定的愚行和虚空之乐；妇女们被这些抬高并装上翅膀，常常飞离婚姻的纽带。因此，妇女也应穿着整洁，用贞洁端庄的带子束紧自己，免得因轻飘而偏离真理。所以，男子应当信任妻子，将家务管理托付给她们，因为她们被赐予作此事的助手。\par

若我们在处理公务或从事其他乡间事务时，常常离开妻子，有必要为安全而封印某物，圣言只为此目的允许我们使用印章。其他戒指应丢弃，因为按照圣经：\BibleQuote{教诲是智者的金饰。}\BibleRef{德 21:21}\par

但戴金饰的妇女在我看仿佛害怕，若有人剥去她们的珠宝，她们就会被当作没有装饰的仆人。然而，以灵魂为居所的本真之美所发现的真理的高贵，判断奴隶不是根据买卖，而是根据奴性的心态。我们有责任不是显得自由，而是真正自由，受天主训练，被天主收纳为义子。\par

因此，我们必须采纳一种站姿、动作、步伐、衣着，总之，一种生活方式，各方面尽可能与自由人相称。男子不应将戒指戴在关节上；这是女性的做法，而应戴在小指的根部。这样手在任何需要时都最便于工作，而印章也不会轻易脱落，因被关节的大节所保护。\par

我们的印章应是鸽子、鱼、顺风航行的船、\Person{波利克拉特斯}所用的竖琴、或是\Person{塞琉古}刻为徽号的船锚；若有渔夫，他将记起宗徒，以及从水中拉出的孩子们。因为我们不可描绘偶像的面容——我们被禁止依附它们；也不可刻刀剑或弓，因为我们追随和平；也不可刻酒杯，因我们节制。\par

许多放荡者将他们的情郎或情妇刻在印章上，彷佛要借著不断想起自己的放荡，使纵欲之事永不忘怀。\par

% structure: p0130 source=h3 target=subsubsection reason=relative-heading-rank; base=h2

\subsubsection{头发}

关于头发，以下似乎正确。男人的头应剃净，除非有卷发；但下巴须留须。不要让卷曲的发绺从头上垂落，滑入女子的卷发样式。浓密的胡须对男人已足够。若有人剃去部分胡须，不可使其全光，因为这是可耻的景象。将下巴剃至皮肤应受谴责，近乎拔毛与磨光。例如，圣咏作者喜悦胡须，说：\BibleQuote{如同香膏流在胡须上，\Person{亚郎}的胡须上。}\par

他反复颂扬胡须之美，使面容因主的香膏而发光。\par

剪发并非为雅致，而是因实际需要；头顶的头发不可长得垂下挡住眼睛，胡须亦同样，吃东西时会弄脏，应用剪子修剪，不可用剃刀，因那不够文雅。但下巴的胡须不应扰动，因它不带来麻烦，且赋予面容尊严与严父般的威仪。\par

此外，发型教导许多人不可犯罪，因为它使行迹易于察觉。对于那些不愿公开犯罪的人，一种不引人注意、不显眼的习惯最合其意；他们一旦养成此习惯，便能暗中犯过而不被发现；如此，因与众人无异，他们便肆无忌惮地犯罪到底。剃头不仅表明一个人庄重，且使头颅不易受伤，因它使头颅习惯冷热；并避免由此产生的危害，因为头发像海绵一样吸收这些湿气，从而不断对大脑造成伤害。\par

妇女保护自己的发绺，用简单的发夹将头发在颈后束起，以朴素护理培养贞洁的发绺达致真美，这已足够。因为妖艳的编发与梳成辫子，只会使她们显得丑陋，剪掉头发并拔去那些奸诈的编结；因此她们不敢碰触头部，生怕弄乱头发。甚至入睡时也提心吊胆，唯恐不明如何弄乱辫子形状。\par

然而，使用他人的头发应完全拒绝，用假发遮头、用死发覆盖头颅，是最亵神之举。因为司铎的手按在谁头上？他祝福的是谁？不是那打扮的妇人，而是别人的头发，并通过头发祝福另一个头。如果\BibleQuote{男人是女人的头，而天主是男人的头}，那么他们犯双重罪岂非不敬？因为他们以过量的头发欺骗男人；并尽己所能，通过妖冶打扮来掩饰真理，从而羞辱主。他们也诋毁那真正美丽的头。\par

因此，头发不可染色，白发也不可改变颜色。因为我们也不可变化衣饰。最重要的是，老年（它博得信任）不可隐藏。天主的尊荣标记应在白日下显扬，以赢得年轻人的尊敬。因为有时，在年轻人行为可耻时，皓首的出现如导师般临至，将他们转为清醒，并以光景的光辉麻痹少年的情欲。\par

% structure: p0138 source=h3 target=subsubsection reason=relative-heading-rank; base=h2

\subsubsection{涂脸}

女人也不应用诡诈的狡猾所设的圈套来涂抹自己的脸。让我们向她们展示节制的装饰。因为首先，最美的美是精神的美，正如我们常指出的。当灵魂被圣神装饰，并受祂所发出的灿烂魅力——正义、智慧、刚毅、节制、好善、谦逊——所激励时，世上从未见过比这更鲜艳的色彩。然后，肉体的美也当培育，即肢体和器官的匀称，以及白皙的肤色。健康的装饰在此是合适的，通过它，人工形象依照天主所赐的样式过渡到真理。但饮料的节制和食物的适度，能有效产生合乎自然的美；因为身体不仅由此保持健康，它们也使美显现。因为从火产生光芒和火花；从潮湿产生明亮和优雅；从干燥产生力量和坚固；从气产生呼吸自如和平衡；由此，圣言这匀称而美丽的形象得以装饰。美是健康的自由花朵；后者产生于身体内部，而前者从身体绽放，显现出明显的肤色之美。因此，这些最庄重和健康的习惯，通过锻炼身体，产生真实持久的美，热量将一切湿气和冷气吸引到自己身上。热量在被运动原因搅动时，是一种吸引自身的东西；而当它吸引时，它通过温暖的肉体本身轻柔地呼出丰富的食物，带着一些湿气，但热量过剩。因此，最初的食物被带走。但当身体不动时，所消耗的食物并不附着，而是脱落，就像从冷炉中取出的面包，要么完整，要么只留下下部。因此，那些不通过运动所需的摩擦排出排泄物的人，其\Emph{\OriginalTerm{粪便}{fœces}}过多。他们身上也充满其他多余物，还有汗液，因为食物没有被身体同化，而是流出浪费。由此，情欲也被激发，冗余物通过相应的运动流向\Emph{\OriginalTerm{阴部}{pudenda}}。因此，这种冗余应当被液化和分散以助消化，由此美获得红润的色泽。但对于那些按\Quote{天主的肖像和模样}受造的人来说，采用外来的装饰，以人的有害发明取代天主的创造，从而羞辱原型，这是怪异的。\par

导师命令她们出去时\BibleQuote{穿着端庄，以羞怯和节制装饰自己，}\BibleRef{弟前 2:9} \BibleQuote{服从自己的丈夫；这样，若有人不信从道理，他们可以因妻子的品行，不用言语而被感化；他们看见，}他说，\BibleQuote{你们贞洁的品行。你们的装饰不应是外面的编发、戴金饰、穿美衣；而应是那隐藏于内心的人，以不朽的温柔宁静的精神为装饰，这在天主面前是极宝贵的。}\BibleRef{伯前 3:1}\par

首先，她们亲手劳动为妇女增添真正的美，锻炼身体，并通过自己的努力装饰自己；而不是带来他人制作的、没有装饰的装饰，那是庸俗和轻浮的，而是每个好妇女的装饰，由她自己的双手在必要时供应和编织。因为对于生活合乎天主的妇女来说，穿从市场买来的衣物永远不合适，而应穿自己家里制作的。因为最美丽的是节俭的妻子，她用自己的劳作为自己和丈夫穿上美丽的衣裳；为此所有人都高兴——孩子因母亲，丈夫因妻子，她因他们所有人，而所有人因天主。\par

简而言之，\BibleQuote{贤淑的妇女是卓越的宝库，她不游手好闲地吃饭；仁慈的法令在她舌上；她开口明智而正直；她的儿女起来称她有福，}正如圣言藉\Person{撒罗满}所说：\BibleQuote{她的丈夫也赞美她。因为虔诚的妇女是有福的；让她赞美敬畏主。}\par

又说：\BibleQuote{贤淑的妻子是丈夫的花冠。}\BibleRef{箴 12:4}她们必须尽可能纠正自己的姿态、眼神、步履和言语。因为她们不应像某些人那样，模仿喜剧表演，练习舞蹈演员的娇柔动作，在社会中如舞台上般举止，带着淫荡的动作、滑行的步伐和做作的声音，投出慵懒的眼神，装饰着快乐的诱饵。\Quote{因为妓女的嘴唇滴下蜂蜜；她说话讨人喜欢，润滑你的喉咙。但最终你会发现它比胆汁更苦，比双刃剑更锋利。因为愚昧的脚引那些行此道的人死后进入地狱。}\par

高贵的\Person{三松}被妓女征服，又被另一个女人剪去了男性的力量。但\Person{若瑟}并没有被另一个女人如此引诱。埃及的妓女被征服了。贞洁，以束缚为己任，显得优于放荡的放纵。以下所说的极为精妙：\par

\Quote{总之，我不知如何\newline{}低语，也不娘娘腔地\newline{}歪着脖子走路，\newline{}如同我看见其他人——那些城中众多的\newline{}色鬼，拔掉了头发。}\par

但是女性的动作、放荡和奢侈应当完全禁止。因为走路时的淫荡动作\Quote{和扭捏的步伐}，正如\Person{阿那克里翁}所说，完全是妓女式的。\par

\Quote{依我看，}喜剧说，\Quote{是时候放弃轻浮的步态与奢侈了。}娼妓的步履不趋向真理，因为它们不走近生命的道路。她的踪迹危险，且不易辨认。眼睛尤其要节俭使用，因为失足于脚不如失足于眼。因此，吾主极简要地治愈此疾：\BibleQuote{倘若你的眼使你跌倒，就把它剜出来，}祂说，从根基上拔除情欲。但乏力的顾盼与飞眼——即以眼示意——无非是以眼行淫，情欲通过它们交战。因为全身之中，眼睛最先被毁。\Quote{眼目观赏\OriginalTerm{美好的事物}{\GreekText{καλά}}，使心灵喜悦；}即那学会\OriginalTerm{正确地}{\GreekText{καλῶς}}观看的眼目令人喜悦。\BibleQuote{以眼示意、心怀诡诈的，必招致灾祸于人。}\BibleRef{箴 10:10}他们如此描写阴柔的亚述王萨达纳帕卢斯，坐于卧榻，双腿抬起，摆弄其紫袍，翻起眼白。效法此等行径的女子，以目光自献为娼。\BibleQuote{身体的光是眼睛，}经上说，内在因照耀之光而显露。\BibleQuote{女子的淫乱，表现在眼目的抬起。}\BibleRef{德 26:9}\par

\BibleQuote{所以要治死你们在地上的肢体，即淫乱、污秽、邪情、恶欲和贪婪——贪婪即是偶像崇拜；因这些事，天主的忿怒临到悖逆之子身上，}\BibleRef{哥 3:5}宗徒如此呼喊。\par

但我们却点燃情欲，且不以为耻。\par

有些这类女子咀嚼乳香，走来走去，向近人露齿；另一些则仿佛没有手指，故作姿态，用发簪搔首；这些发簪或以龟甲、或以象牙、或由某种死物制成，费尽心力。还有些人，仿佛具有某种绽放之花，为在观众眼中显得秀丽，以艳丽的膏油涂饰面容。这样的女子被撒罗满称为\BibleQuote{愚昧而狂妄的妇人，}她\BibleQuote{不知羞耻。她坐在自家门口，显眼地在座位上，呼唤一切行路的、直行其道的人；}以她的举止和整个生活明言：\BibleQuote{你们谁最愚笨？让他转到我这里来。}对缺乏智慧的人她劝勉说：\BibleQuote{偷偷触摸甜蜜的饼，偷来的水是甘甜的；}以此暗指隐秘之爱（由此，博奥提亚的品达前来相助说：\Quote{隐秘的追求爱情是甜蜜的}）。但那可怜之人\BibleQuote{不知尘世之子必在她身旁灭亡，她趋向地狱的深渊。}但导师说：\BibleQuote{速速离开，不要停留在那地方；也不要定睛看她；因为你将如此渡过异水，穿越到阿刻戎。}\BibleRef{箴 9:13}因此，主藉依撒意亚这样说：\BibleQuote{因为熙雍的女子昂首而行，眼目顾盼，行走时曳裙，以足戏耍；上主必使熙雍的女子卑微，显露她们的形态}——她们畸形的形态。我以为，跟随贵妇的婢女向她们言语或举止不当，是有罪的。但我以为她们应受主母的纠正。因此，喜剧诗人费勒蒙以极严厉的谴责说：\Quote{你可以跟在美貌婢女身后，她出现在淑女之后；任何人从普拉泰依库姆都可紧随其后，向她睇视。}因为婢女的放荡归咎于主母，让那些试图稍越礼数的人不惧进而逾矩；因为主母容许不当行为，表明她并不反对。而不对放荡者发怒，便是心性倾向同样的确证。\Quote{有其主必有其婢，}正如谚语所说。\par

% structure: p0151 source=h3 target=subsubsection reason=relative-heading-rank; base=h2

\subsubsection{行走}

我们也必须弃绝狂躁的行走方式，选择庄重而从容、但不拖沓的脚步。\par

也不可在大路上昂首阔步，或扭头回望所遇之人（若他们看他），犹如在舞台上趾高气扬、受人指点。也不可在上山时让仆从推着走，如我们所见那些更奢侈者——他们看似强壮，却因灵魂的阴柔而衰弱。\par

一位真正的君子，脸上或身体任何部分，不可有阴柔的痕迹。因此，在他的举动或习惯中，永远不可发现任何男子气概的污点。健康之人也不可利用仆人如驮马驮负自己。因为正如伯多禄所说，命令他们\BibleQuote{以一切敬畏服从主人，不单对良善温柔的，也对乖僻的}\BibleRef{伯前 2:18}；同样，公正、宽容、仁慈正是主人应有的。因为他说：\BibleQuote{总之，你们都要同心，彼此体恤，相爱如弟兄，仁慈，谦卑，等等，好使你们承受祝福，}美好而可羡慕。\par

% structure: p0155 source=h3 target=subsubsection reason=relative-heading-rank; base=h2

\subsubsection{模范少女}

基提翁的芝诺认为应当描绘一位少女的形象，并如此塑造雕像：\Quote{让她的面容洁净，眉毛不下垂，眼睑不张开或翻起。她的颈项不向后仰，身体的肢体不松弛。但身体悬挂的部分应如紧绷一般；有良好规整的心灵为辩论而敏锐，并牢记正确所说；她的态度与动作不给予放荡者任何希望；但应有端庄的光彩与坚定的表情。远离她的，是来自香贩、金匠、羊毛商之店的烦扰，以及来自其他店铺的——那些店铺中，女子装饰得妖冶，终日如坐于娼寮一般的烦恼。}\par

% structure: p0157 source=h3 target=subsubsection reason=relative-heading-rank; base=h2

\subsubsection{娱乐与同伴}

所以，男人不应在理发店和酒馆里虚度光阴，胡言乱语；他们应放弃追逐坐在近旁的女人，以及不断诽谤多人以博一笑。\par

掷骰子游戏应被禁止，追求利益，尤其是通过掷骰子，许多人热衷于此。奢侈的挥霍为闲人发明了这些东西。原因是懒惰，以及对远离真理的轻浮之事的爱好。因为若非如此，就不可能毫无伤害地获得快乐；而每个人对生活方式的偏好是其性情的一面镜子。\par

但看来，只有与善人来往才有益；另一方面，全智的导师藉\Person{梅瑟}的口，认定了与恶人交往如同猪一般，便禁止古代的人民吃猪肉，以指出那些呼求天主的人不应与不洁的人混杂，这些人像猪一样，喜爱肉体的快乐、不洁的食物，并因不洁的淫欲而对邪恶的享乐感到瘙痒。\par

祂接着说：\Quote{不可吃鸢或快翼的猛禽，或鹰。}\BibleRef{肋 11:13}意思是：你不可接近以劫掠为生的人。还有其他事物也是以比喻方式展示的。\par

那么，我们应与谁交往？祂再次以比喻方式说：因为\Quote{凡蹄分两瓣且反刍的动物，是洁净的。}因为分蹄表明正义的平衡，反刍则指向正义的适当食物，即圣言，它从外进入，如食物一般，通过教导，但从心思中召回，如同从胃中召回，进行理性的回忆。属神的人，口中含有圣言，反刍属神的食物；正义正确地分蹄，因为它在今生圣化我们，并送我们踏上通往来世之路。\par

% structure: p0163 source=h3 target=subsubsection reason=relative-heading-rank; base=h2

\subsubsection{公共表演}

导师不会带我们去公共表演；将赛马场和剧场称为\Quote{灾祸之座}也未尝不可；因为那里有反对正义者的邪恶计谋\BibleRef{宗 3:14}，因此反对祂的集会应受诅咒。这些集会确实充满混乱和不义；这些集会的借口是混乱的根源——男女混杂，仅仅为了相看。在这方面，集会已经显出其恶劣：因为眼睛淫荡时，欲望升温；习惯于在闲暇时无耻地注视邻人的眼睛，会点燃情欲。因此，应禁止充满粗鄙和大量闲谈的表演和戏剧。因为有什么卑劣的行为不在剧场中展示？有什么无耻的话语不被小丑们说出？那些享受其中邪恶的人，在家中刻下清晰的影像。另一方面，那些对这些事情有抵抗力、不受影响的人，绝不会在奢侈享乐上失足。\par

因为若有人说他们去观看表演是为了消遣娱乐，我要说，那些认真对待消遣的城市是不明智的；因为那些导致如此致命后果的残酷荣耀竞赛不是娱乐。同样，无意义的金钱花费也不是娱乐，由它们引起的暴乱也不是娱乐。心灵的安宁不能通过热衷追求轻浮之事而获得，因为一个有理智的人绝不会将愉悦置于善之上。\par

% structure: p0166 source=h3 target=subsubsection reason=relative-heading-rank; base=h2

\subsubsection{日常生活中的宗教}

但有人说我们并非所有人都哲学。那么我们岂不都追求生命吗？你说什么？你是如何相信的？请问，你若不爱哲学，怎能爱天主和你的邻人？你若不爱生命，怎能爱自己？有人说，我没有学过文字；但即使你没学过阅读，在听方面你也不能原谅自己，因为听不是学来的。信德不是属于世上智慧人的，而是属于天主的人的；它不需要文字就能教导；它的手册，既粗糙又神圣，被称为爱——一本属灵的书。你有能力聆听天主的智慧，是的，并按照它塑造你的生活。不仅如此，你并不被禁止按照天主体面地处理世事。买卖任何东西的人，不要对所买或所卖的东西报两个价格；而要说明净价，并努力说真话，即使他没得到他要的价，他得到了真理，并在拥有正直方面是富有的。但最重要的是，要远避因所售之物而起誓；因其他事而起誓也应被驱逐。\par

这样，那些常去市场和店铺的人也在哲学。因为\BibleQuote{不可妄呼上主你的天主的名；因为凡妄呼祂名的人，上主决不会视为无罪。}\BibleRef{出 20:7}\par

但那些行事与此相反的人——贪婪者、说谎者、伪善者、以真理做生意的人——主将他们赶出祂父的殿宇，不愿天主的圣殿成为不义交易的场所，无论是言语上的还是物质上的。\par

% structure: p0170 source=h3 target=subsubsection reason=relative-heading-rank; base=h2

\subsubsection{去教堂}

男女去教堂应衣着得体，步履自然，保持沉默，拥有真诚的爱，身体纯洁，心灵纯洁，适于向天主祈祷。此外，女人应遵守以下：她应完全遮盖，除非在家。因为那种着装庄重，保护不被注视。将谦逊和披肩放在眼前的女子永不跌倒；她也不会因露脸而引他人犯罪。因为这是圣言的意愿，因为蒙头祈祷是合宜的。\par

据说\Person{埃涅阿斯}的妻子，由于过分的得体，即使在特洛伊陷落的恐惧中也没有暴露自己；尽管逃离大火，仍蒙着头。\par

% structure: p0173 source=h3 target=subsubsection reason=relative-heading-rank; base=h2

\subsubsection{在教堂外}

这样的人应当是奉献于\Person{基督}的人，他们在整个生活中如此塑造自己，如同在教会中为了庄重而塑造自己那样；他们应当如此，而非看似如此——如此温良、如此虔诚、如此有爱德。但现在我不知道人们如何随着地点改变他们的时尚和举止。正如他们所说，章鱼会附着于岩石而变得与岩石同色；同样，离开集会之后，他们卸下聚会时的灵感，变得与交往的其他人一样。不仅如此，他们卸下做作的庄重面具，显露出他们原本隐藏的样子。在向关于天主的言论表示敬意之后，他们将所听到的留在教会内。而在外面，他们愚蠢地以不虔的游戏和情歌颤音自娱，忙于吹笛、跳舞、酗酒及各种垃圾。这样唱歌和应和的人，正是那些先前歌颂不朽的人，最终变得邪恶，邪恶地唱着这首最为有害的翻案诗：\Quote{我们吃喝吧，因为明天我们就要死了。}但事实上不是明天，而是现在，这些人已经对天主死了；埋葬他们的死者，\BibleRef{玛 8:22}，也就是说，使自己沉入死亡。宗徒坚定地斥责他们：\BibleQuote{不要被迷惑；无论是奸淫者、孪童者、与人苟合者、偷窃者、贪婪者、醉酒者、辱骂者}，以及他加上的其他任何，\BibleQuote{都不能继承天主的国。}\BibleRef{格前 6:9}\par

% structure: p0175 source=h3 target=subsubsection reason=relative-heading-rank; base=h2

\subsubsection{爱德与爱德之吻}

如果我们被召叫进入天主的国，就让我们相称于天国而行，爱天主和我们的近人。但爱德不是藉亲吻证明的，而是藉善意。但有些人，他们只让教会响彻亲吻，内心却没有爱德本身。因为正是这件事，即对亲吻的无耻使用（这亲吻本应是奥秘的），引起了可恶的猜疑和恶意的传言。宗徒称这亲吻是圣的。\BibleRef{罗 16:16}\par

当天国被相称地考验时，我们通过贞洁而紧闭的口传达灵魂的情感，主要由此表达温良的举止。\par

但还有另一种不圣的亲吻，充满毒药，冒充圣洁。难道你们不知道，蜘蛛仅仅接触嘴巴就会使人疼痛吗？而且亲吻常常注入放荡的毒药。因此对我们来说非常明显：亲吻不是爱德。因为这里所指的爱德是天主之爱。\Person{若望}说：\BibleQuote{这就是爱德：我们遵守祂的诫命；}，\BibleQuote{而且祂的诫命并不沉重。}\BibleRef{若一 5:3}并不是我们彼此抚摸嘴巴。但是，在路途中向所爱之人打招呼，充满了愚昧的放肆，正是那些想在旁人面前显眼之人的特征，他们没有丝毫恩宠。因为如果奥秘地\Quote{在内室}中向天主祈祷是合宜的，那么我们也应奥秘地在屋内问候我们的近人（我们被命令同样地爱他仅次于天主），\Quote{赎回时间。}\BibleQuote{因为我们是地上的盐。}\BibleRef{玛 5:13}\BibleQuote{凡清晨大声祝福朋友的人，将被视为与咒骂无异。}\BibleRef{箴 27:14}\par

% structure: p0179 source=h3 target=subsubsection reason=relative-heading-rank; base=h2

\subsubsection{眼的管束}

但首要的是，似乎我们应当避开观看妇女。因为不仅触摸是罪，观看也是；而且受过正确训练的人必须特别避免它们。\BibleQuote{你的眼睛要向前直视，你的睫毛要正直眨眼。}\BibleRef{箴 4:25}因为虽然观看的人有可能保持坚定，但仍须谨慎以防跌倒。因为观看的人有可能滑倒；但不看的人不可能产生情欲。因为对于贞洁者而言，单单纯洁是不够的；他们必须全力以赴，使自己不受指责，杜绝一切怀疑的缘由，以达到贞洁的圆满；这样我们不仅信实，而且显得值得信任。正因为此，也必须防范，如宗徒所说：\BibleQuote{免得有人责怪我们；我们要预备光明的事，不但在主面前，也在人面前。}\BibleRef{格后 8:20}\BibleQuote{但你要转开你的眼不看优雅的妇女，不要注视他人的美貌，}圣经说。\BibleRef{德 9:8}如果你问原因，它会进一步告诉你：\BibleQuote{因为妇女的美貌使许多人迷失，而且为此情欲如火般燃烧；}\BibleRef{德 9:8}这情欲由我们称为爱的火而生，引向因罪而永不熄灭的火。\par

% structure: p0181 source=h2 target=subsection reason=relative-heading-rank; base=h2

\subsection{第十二章 续：带着圣经经文}

我要劝告已婚者绝不要在仆人面前亲吻自己的妻子。因为\Person{亚里士多德}不允许人对奴隶发笑。也绝不能让妻子在他们面前被问候。此外，更好的是，在家中从婚姻开始，我们应当在其中表现出得体。因为它是贞洁的最大纽带，呼出纯真的愉悦。悲剧非常美妙地说道：——\par

\Quote{好啊！好啊！夫人们，那么，在男人中间，\newline{}黄金、帝国或财富的奢华，\newline{}并未带来如此显著的愉悦，\newline{}如同有德之人的正直气质\newline{}与一位尽责的妻子。}\par

那些通晓世俗智慧的人所发出的此类正义训诫，不可拒绝。既然知道了各人的本分，\BibleQuote{你们在寄居的时候，当怀着敬畏度过。因为你们知道，你们不是用可朽坏的东西，如同金银，从你们祖先所传授的虚妄生活中被赎出来的，而是用基督的宝血，如同无瑕疵、无玷污的羔羊之血。}\BibleRef{伯前 1:17}\Quote{因为，}伯多禄说，\BibleQuote{过去我们生活中的时间，足以让我们行了外邦人的意志，当我们行在放荡、情欲、酗酒、狂欢、宴饮和可憎的偶像崇拜中。}\BibleRef{伯前 4:3}我们有十字架作为界限，这十字架是我们的主，藉着它我们从先前的罪恶中被围护和篱笆围住。因此，我们既已重生，就该真诚地固守其上，回归清醒，并圣化我们自己；\BibleQuote{因为主的眼看顾义人，祂的耳听他们的祈祷；但主的面容反对作恶的人。}那么，如果我们追随美善，谁能伤害我们呢？\BibleRef{伯前 3:13}——\Quote{我们}代替\Quote{你们。}但最好的训练是良好的秩序，即完美的礼节，以及稳定有序的力量，在行动中保持所行的一致性。如果我过于严厉地引证这些事，为成就因你们的纠正而来的救恩；那么，这些也是由我所说的，导师说：\Quote{因为勇敢责斥的人，是缔造和平者。}如果你们听我，你们必得救。如果你们不听我所说，那不是我的事。然而这也是我的事这样：\BibleQuote{因为祂愿意罪人悔改，而不愿他死亡。}\BibleRef{则 18:23}\Quote{如果你们听我，你们必吃土地的出产，}导师又说，以\Quote{土地的出产}之称，指美貌、财富、健康、力量、食物。因为那些真正美善的事，是\BibleQuote{耳朵未曾听见，心也未曾想到的}\BibleRef{格前 2:9}，关于那位真正为王的那一位，以及等待我们的真正美善的现实。因为祂是美善之事的赐予者和守护者。关于参与这些美善，祂将同样的事物名称应用于此世，圣言就这样借着天主，将人的软弱从感官事物训练到理解。\par

在家应当遵守什么，以及我们的生活如何规范，导师已经充分说明。祂在指引孩子们走向师傅的路上常对他们说的话，祂提出这些，并以简明的形式引用圣经本身，列出纯粹的诫命，使其适应指导时期，并将对这些诫命的解释指派给师傅。因为祂法律的意图是驱散恐惧，解放自由意志以达致信德。\Quote{听，}祂说，\Quote{孩子，}你已受到正确指导，这是救恩的要点。因为我要启示我的道路，将善的诫命摆在你面前，藉此你将达致救恩。我引领你走救恩之路。离开欺诈之道。\par

\Quote{因为主知道义人的道路，恶人的道路却必灭亡。}\Quote{所以，儿子啊，跟随我所描述的美善道路，向我侧耳倾听。}\BibleQuote{我要将隐藏的、黑暗中的宝藏赐给你，}这是外邦人所未见、但我们所见到的。\BibleRef{依 45:3}智慧的宝藏是无穷的，宗徒赞叹它而说：\BibleQuote{哦，天主智慧和知识的富饶何其深！}\BibleRef{罗 11:33}并由一位天主分施众多宝藏，有些由法律启示，有些由先知启示，有些来自神圣的口，有些来自圣神的七重歌声和谐一致。而主是唯一的，祂藉着这一切是同一的导师。这里是一条全面的诫命，一个涵盖一切的生命劝勉：\BibleQuote{你们愿意人怎样待你们，你们也要怎样待人。}\BibleRef{路 6:31}我们可以将诫命总结为两条，正如主所说：\Quote{你要全心、全灵、全力爱主你的天主；并爱邻人如同自己。}然后祂从中推论：\Quote{法律和先知都系于这两条。}此外，当有人问：\Quote{我该行什么善事才能得永生？}祂回答：\Quote{你知道诫命吗？}那人回答知道，祂说：\Quote{你这样做，就必得救。}尤其显著的是导师在各样有益的诫命中所彰显的爱，使得透过圣经的丰富和安排，发现更为容易。我们有梅瑟所赐的十诫，它透过一个基本的原则，简单而划一，以有益于救恩的方式界定了罪的名目：\Quote{不可奸淫。不可崇拜偶像。不可败坏男童。不可偷盗。不可作假见证。应孝敬你的父亲和母亲。}\EditorialNote{出谷纪20章；申命纪5章}等等。这些事必须遵守，以及其他在读经中所命令的一切。祂藉依撒意亚命令我们：\Quote{你们要洗涤，自洁，从我眼前除掉你们灵魂的邪恶。学习行善，寻求正义，解救受欺压的，为孤儿伸冤，为寡妇辩护。然后来，我们彼此辩论，主说。}我们在别处也会发现许多例子——例如关于祈祷：\Quote{善行是蒙主悦纳的祈祷，}经上说。祈祷的方式也被描述：\Quote{你若看见赤身露体的，就给他衣服穿；不可忽略你同族的人。这样，你的光必如晨光出现，你的伤口必迅速愈合；你的正义在你前面行，天主的荣耀必护卫你。}那么，这种祈祷的果实是什么呢？\Quote{那时你求告，天主必应允你；你还在说话，祂必说:我在这里。}\par

论及禁食，经上记载：\Quote{你们为什么向我禁食呢？}主说：\Quote{难道我所拣选的禁食是这样的吗？是叫人刻苦自己的一天吗？你们不要把头弯得像环，把麻布和灰铺在下面。你们不可如此称之为可悦纳的禁食。}\par

那么，禁食是什么意思呢？\Quote{看，我所拣选的禁食是这样的，主说：解开一切邪恶的捆绑，解除压迫契约的绳结，让受压迫者获得自由，撕毁一切不义的债据。把你的饼分给饥饿的人，领无家可归的穷人到你家中。若看见赤身露体的人，就遮盖他。}\BibleRef{依 57:6}关于祭献也是如此：\Quote{你们献上许多祭物，于我有何益？主说。我已饱享全燔祭和公羊的油脂；羔羊的脂油、公牛和山羊的血，我都不喜爱；也不要你们到我面前来朝拜。谁向你们要求这些？你们不可再践踏我的庭院。若你们奉献细面，虚浮的供品是我所憎恶的。你们的月朔和安息日，我都不能容忍。}\BibleRef{依 1:11}那么，我应怎样向主献祭呢？祂说：\Quote{主的祭献就是破碎的心。}那么，我应怎样加冕，或涂抹香膏，或向主献香呢？经上说：\Quote{馨香的香气就是光荣创造祂的心。}这些就是天主的冠冕、祭献、芬芳的香气和花朵。\par

此外，关于容忍。经上说：\Quote{如果你的兄弟得罪了你，你要规劝他；他若悔改，你就宽恕他。如果他一天七次得罪你，又七次回转来说：\InnerQuote{我后悔了}，你总该宽恕他。}\BibleRef{路 17:3}同样，若翰对士兵们命令\Quote{只该满足于自己的粮饷}；对税吏说\Quote{不可征收超过规定的数目}。对审判者祂说：\Quote{审判时不可偏袒。因为贿赂会蒙蔽明眼人的眼睛，败坏公正的言语。要拯救受屈的人。}\par

对家主们说：\Quote{以不义得来的财物必会减少。}\BibleRef{箴 13:11}\par

论到\Quote{爱}，祂说：\Quote{爱能遮盖许多罪。}\BibleRef{伯前 4:8}\par

论及世俗政权：\Quote{凯撒的归凯撒，天主的归天主。}\par

论起誓和记仇：\Quote{我岂曾吩咐你们的祖先，当他们出离埃及时，献全燔祭和祭品吗？我吩咐他们的是：你们各人心里不可怀恨邻人，也不可喜爱虚假的誓言。}\par

祂也威胁说谎者和骄傲的人，前者如：\Quote{祸哉，那些称恶为善、称善为恶的人！}；后者如：\Quote{祸哉，那些在自己眼中看为智慧、在自己面前看为精明的人。}\BibleRef{依 5:20}\Quote{因为自谦的必被高举，自高的必被贬抑。}\par

祂祝福\Quote{怜悯人的人，因为他们要蒙受怜悯。}\par

智慧宣称愤怒是可怜的事，因为\Quote{它会毁灭智慧人}。现在祂命令我们\Quote{爱你们的仇敌，祝福诅咒你们的人，为虐待你们的人祈祷。}祂又说：\Quote{若有人打你的一边脸颊，把另一边也转给他；若有人拿你的外衣，连内衣也不要阻止他拿。}\BibleRef{玛 5:40}\par

论\OriginalTerm{信德}{faith}，祂说：\Quote{你们在祈祷时，无论求什么，只要相信，就必得到。}\BibleRef{玛 21:22}\Quote{对不信的人，没有什么是可信的}，\Person{品达}如是说。\par

同样，对待仆人应如同对待我们自己，因为他们与我们是同样的人。因为天主对自由人和奴隶是一样的，只要你思考。\par

我们犯罪的手足，不可惩罚，只应规劝。\Quote{因为姑息杖责的，是恨恶自己的儿子。}\BibleRef{箴 13:24}\par

此外，祂彻底驱逐贪爱荣耀之心，说：\Quote{祸哉，你们法利塞人！因为你们喜爱会堂里的首位，和市场中的问安。}\BibleRef{路 11:43}但祂欢迎罪人的悔改——祂喜爱悔改——这悔改在犯罪之后。因为我们所论的这\OriginalTerm{圣言}{Word}是唯一无罪的。因为犯罪是自然且普遍的。但犯罪后回归（天主）并非任何人的特征，而只属于有价值的人。\par

论及慷慨，祂说：\Quote{你们这些蒙祝福的，到我这里来，承受自创世以来为你们预备的国度：因为我饿了，你们给了我吃的；我渴了，你们给了我喝的；我作客，你们收留了我；我赤身露体，你们给了我穿的；我生病，你们看顾了我；我在监里，你们来探望了我。}而我们何时曾向主做过这些事呢？\par

导师自己又要说话，祂喜爱将弟兄们的善行归于自己：\Quote{你们对我这些最小兄弟中的一个所做的，就是对我做的。这些人要进入永生。}\par

圣言的这些法律，这些安慰的话语，不是写在石版上——那是用主的手指所写的——而是铭刻在人的心里，只有在那里它们才能不朽。因此，那些心硬之人的石版被打破了，好让子民的信德得以烙印在软化的心上。\par

然而，两种法律——梅瑟所赐的律法和宗徒所传的律法——都服务于圣言为教导人类。因此，由宗徒所传授的训练之性质，在我看来需要加以论述。在此标题下，我将叙述，或者更确切地说，由导师通过我来叙述；我将再次把那些规条本身放在你们面前，如同在萌芽状态。\par

\Quote{弃绝谎言，各人与邻人说真话，因为我们彼此互为肢体。不可让太阳落在你们的忿怒上，也不可给魔鬼留地步。偷窃的，不要再偷；却要劳力，亲手做正当的事，使自己有所余，可赒济有需要的人。一切毒恨、忿怒、暴躁、喧嚷、毁谤，连同一切恶毒，都该从你们中除掉；彼此要以仁慈相待，心怀温良，互相宽恕，如同天主在基督内宽恕了你们一样。所以你们该效法天主，如同蒙爱的儿女，并在爱德中行走，正如基督也爱了我们。妻子要服从自己的丈夫，如同服从主。丈夫要爱妻子，如同基督也爱了教会。} 那些配合为一的，要彼此相爱，\Quote{如同自己的身体。} \Quote{孩子们，要服从你们的父母。父母不要惹你们的子女发怒，却要以主的训诲和劝诫养育他们。仆人，要战战兢兢、以诚实的心服从你们肉身的主人，如同服从基督；要甘心服事，如同从心里做。你们做主人的，要善待仆人，不要威吓，因为知道他们和你们的主都在天上，而且祂不以貌取人。}\par

\Quote{如果我们生活于圣神，就让我们遵圣神而行。不要贪图虚荣，彼此激怒，彼此嫉妒。彼此分担重担，如此便满了基督的法律。不要自欺；天主是轻慢不得的。我们行善不要厌倦，因为到了适当的时候，我们必要收获，如果我们不松懈的话。}\par

\Quote{你们要彼此和睦。弟兄们，我们劝你们，劝戒那些不守规矩的人，安慰怯懦的人，扶持软弱的人，对众人要忍耐。谁也不要以恶报恶。不要熄灭圣神。不要轻视先知讲论。凡事试验；持守那善的。各样恶事要远避。}\par

\Quote{要恒切祈祷，在此警醒感恩。你们在外人中间，要凭智慧行事，爱惜光阴。你们的言语要常常带着和气，如同用盐调和，使你们知道应当怎样回答各人。}\par

\Quote{要在信仰的话语上滋养自己。操练自己达到虔敬；因为身体的操练只有少许益处，但虔敬在一切事上都有益处，有今生和来生的应许。} \BibleRef{弟前 4:6}\par

\Quote{那些有忠信主人的，不可轻视他们，因为他们是弟兄；反而要服事他们，因为他们是忠信的。} \BibleRef{弟前 6:2}\par

\Quote{施与的，要诚心；治理的，要殷勤；怜悯人的，要喜乐。爱不可虚假。恶要厌恶，善要亲近。以兄弟之爱彼此相亲相爱，恭敬人，彼此争先。殷勤不懒惰，心灵火热，服事主。在指望中要喜乐，在患难中要忍耐，恒切祈祷。一味的慷慨好客，圣人有缺乏，要帮补。} \BibleRef{罗 12:8}\par

这些是从许多教训中取出的少数例子，导师遍览圣经，将它们摆在祂的儿女面前；借此，罪孽可说是被连根拔除，不义被约束。\par

无数这样的命令写在圣经上，关乎特定的人，有些涉及长老，有些涉及主教，有些涉及执事，另一些涉及寡妇，关于他们我们会有机会另作讨论。许多以谜语讲述的，许多以比喻讲述的，可能会让遇到它们的人受益。但教导这些已不再是我的范围，导师说。我们需要一位阐释那些神圣言语的导师，我们必须向祂迈步。\par

而现在，事实上是我停止教导、你们聆听导师的时候了。祂接纳你们这些受过卓越操练的人，将教导你们圣言。教会和新郎——那唯一的导师、圣父的良善计谋、真智慧、知识的圣所——已经高雅地歌唱了。\Quote{祂是赎我们罪的祭献，} \Person{若望}说；耶稣，祂医治我们的灵魂和身体——这构成真正的人。\BibleQuote{不但是为我们的罪，也是为全世界的罪。我们若遵守祂的命令，由此就知道我们认识祂。那说\InnerQuote{我认识祂}，却不遵守祂命令的，是撒谎者，真理不在他内。但遵守祂言语的，天主的爱在他内才得以成全。由此我们知道我们在祂内。那说住在祂内的，就该自己照样行走，如同祂行了。} \BibleRef{若一 2:2} 哦，蒙祂祝福操练的儿女！让我们完成教会美丽的面容；让我们像孩子一样奔向我们的好母亲。如果我们成了圣言的聆听者，让我们光荣那有福的安排，人由此被操练并圣化为天主的子，并在天上拥有他的交往，虽在地上受操练，却在那里接受圣父，就是他在世上学习认识的那一位。圣言行一切事，教导一切事，并操练一切事。\par

马由嚼环引导，牛由轭引导，野兽被套索捕获。但人由圣言改变，借着祂，野兽被驯服，鱼被捕获，鸟被引下。事实上，是祂为马制造嚼环，为牛制造轭，为野兽制造套索，为鱼制造竿，为鸟制造网罗。祂既管理国家，又耕田；命令，帮助，并创造宇宙。\par

\Quote{那里有大地、天空和海洋，\newline{}永远绕行的太阳和圆月，\newline{}以及加冕穹苍的一切星辰。}\par

\Quote{让这水在自己内波动；让这火约束其烈怒；让这空气飘入以太；而这地得坚固，得运动！当我要形成人时，我需要物质，并且在元素中有物质。我与我所形成的同住。如果你认识我，火将成为你的奴隶。}\par

圣言正是如此，导师正是如此，世界和人的创造者；而祂自己，如今是世界的导师，藉着祂的命令，我们和宇宙得以存在，并等候审判。\Quote{因为并非那向人带来隐秘言语的}，如\Person{巴库利德斯}所言，\Quote{将成为智慧之言}；而是\BibleQuote{那无可指责、纯洁、无瑕的天主子女}，如\Person{保禄}所说，\BibleQuote{在弯曲悖谬的世代中，像明灯照耀在世界中}。\par

因此，在这般庆祝圣言之时，我们所要做的，就是向圣言献上我们的祈祷。\par

% structure: p0220 source=h3 target=subsubsection reason=relative-heading-rank; base=h2

\subsubsection{向导师的祈祷}

求祢垂怜，导师，对我们祢的子女，圣父，以色列的御者，圣子与圣父，两者合一，主啊。求祢赐予我们这些遵行祢诫命的人，使我们能完美那肖像的相似，并尽我们所能认识那良善的天主，而非严厉的审判者。并求祢亲自使我们所有在祢和平中度日的人，被迁入祢的国度，安然渡过罪的波涛，被祢的圣神平静地吹送，藉着不可言喻的智慧，日夜趋向那完美的日子；并感恩赞美，赞美感谢那唯一的父与子，圣子与圣父，圣子，导师和教师，与圣神同在，全部合一，在祂内有一切，万有归于祂是一，归于祂是永恒，我们都是祂的肢体，万世是祂的光荣；因为祂是全善、全可爱、全智、全公义的那一位。愿光荣归于祂，自今至永远。阿们。\par

既然导师藉着将我们迁入祂的教会，使我们与祂自己联合，那教导并遍察万有的圣言，我们理当到此地步，向主献上应谢的回报——相称于祂美好训诲的赞美。\par

% structure: p0223 source=h3 target=subsubsection reason=relative-heading-rank; base=h2

\subsubsection{颂歌致基督救主}

I.\par

驯服野马之辔，\newline{}统御我等意志；\newline{}不倦飞鸟之翼，\newline{}稳妥引领飞翔。\newline{}倔强青年之舵，\newline{}坚固对抗逆浪；\newline{}牧者，以智慧牧养\newline{}王家的羊羔：\newline{}使祢纯朴的子女齐聚，\newline{}同声歌唱\newline{}以庄严的诗篇\newline{}赞美之颂\newline{}用无邪的唇舌向基督君王。\par

II.\par

圣人君王，全能圣言\newline{}父最高的主；\newline{}智慧之首与元首；\newline{}一切忧伤的安慰；\newline{}时间与空间的主，\newline{}耶稣，我族的救主；\newline{}牧者，护守我们；\newline{}农夫，耕作我们，\newline{}辔头以约束我们，舵\newline{}以指引我们随祢意；\newline{}至圣羊群的天翼；\newline{}渔人的渔夫，祢带来生命；\newline{}从罪恶的恶海，\newline{}并从波涛的争战，\newline{}收集纯净的鱼，\newline{}以生命的甜饵捕获：\newline{}引导我们，羊群的牧者，\newline{}赋有理智的圣者；\newline{}青年的君王，祢护守他们，\newline{}使他们远避污秽：\newline{}基督的脚步，天上的道路；\newline{}永恒的圣言，无尽的世代；\newline{}永不腐朽的生命；\newline{}仁慈的泉源，德能的施予；\newline{}尊贵的生命属于那些向天主\newline{}高唱赞美诗的人，\newline{}耶稣基督！\par

III.\par

以天上的乳汁滋养，\newline{}赐予我们柔嫩的味觉；\newline{}从恩宠新娘的胸怀\newline{}流出的智慧之乳；\newline{}充满露珠般的圣神\newline{}从理性之胸馏出；\newline{}让我们这些乳儿同声\newline{}用纯洁的唇舌高举赞美诗\newline{}作为我们感恩的奉献，\newline{}清洁纯净，献给基督君王。\newline{}让我们以无玷之心，\newline{}庆祝那大能的孩子。\newline{}我们，基督所生，和平的诗班；\newline{}我们，祂爱的子民，\newline{}让我们歌唱，永不止息，\newline{}向至高和平的天主。\par

我们附上上述颂歌的如下直译：——\par

驯服野马之辔，不倦飞鸟之翼，婴孩的稳妥舵，王家羊羔的牧者，召集祢的纯朴子女以圣洁地赞美，以无邪的口纯洁地歌颂基督，儿童的引导者。圣人君王，至高圣父的全胜圣言，智慧之统御者，忧伤之支持者，在万世中喜乐者，耶稣，人类的救主，牧者，农夫，舵，辔，至圣羊群的天翼，得救之人的渔夫，以甜美的生命从罪恶万恶之海的憎恶波浪中捕获纯洁的鱼——引导我们，理性羊群的牧者；引导无伤的儿童，圣洁的君王，基督的脚步，天上的道路，永存的圣言，无限的世代，永恒的光，仁慈的泉源，德能的施行者；高贵的是那些赞美天主者的生命，基督耶稣，新娘恩宠之甜乳的属天之奶，从祢的智慧中挤出的。以柔嫩之口滋养的婴孩，充满理性之乳的露珠之神，让我们同声歌唱纯朴的赞美，真实的颂歌献给基督君王，生命训诲的神圣报偿；让我们在纯朴中歌唱大能的孩子。和平的诗班，基督所生，纯洁的子民，让我们同声歌唱和平的天主。\par

% structure: p0232 source=h3 target=subsubsection reason=relative-heading-rank; base=h2

\subsubsection{向导师}

导师，我向祢献上一顶花冠，\newline{}用从无瑕草地采摘的言语编织，\newline{}在那里祢牧放祢的羊群；如同蜜蜂，\newline{}那精巧的工人，从众多花朵\newline{}采集珍宝，以便呈上\newline{}甜蜜的奉献至主人手中。\newline{}虽微不足道，我仍是祢的仆人，\newline{}（赞美祢的诫命是合宜的）。\newline{}君王，人类伟大恩赐的施予者，\newline{}善的主宰，圣父，万物的创造者，\newline{}祢以神圣的圣言\newline{}妥善安排了天与天的装饰，独自成就；\newline{}祢带来了阳光与白昼；\newline{}祢为星辰设定了轨道，\newline{}以及大地与海洋如何保持位置；\newline{}当季节在其循环的轨道中，\newline{}冬与夏，春与秋，各按\newline{}井然有序的计划来临；\newline{}祢从混乱的一堆创造了\newline{}这有秩序的球体，并从无形物质\newline{}装饰了宇宙——\newline{}赐予我生命，并使这生命善度，\newline{}享受祢的恩宠；使我行事言谈\newline{}在一切事上遵从祢的圣经；\newline{}祢与祢的同永圣言，全智，\newline{}从祢发出，愿我永远赞美；\newline{}求祢勿赐我贫穷或富裕，而是所需，\newline{}圣父，在人生中，然后安息离世。\par

\SourceRecord{Translated by William Wilson.；Ante-Nicene Fathers；Vol. 2.；Edited by Alexander Roberts, James Donaldson, and A. Cleveland Coxe.；Buffalo, NY: Christian Literature Publishing Co.,；1885.；<http://www.newadvent.org/fathers/02093.htm>.}
